\parte{Párodo}{1–154}
\falante{ΧΟΡΟΣ}
\verso{1}{\notagramatical{οἰχομένων}{οἴχομαι}{partir}{part., pres., méd., gen. pl. masc.}}{Τάδε μὲν Περσῶν τῶν οἰχομένων}
\verso{2}{\notagramatical{αἶαν}{αἶα, -ας, ἡ}{terra}{ac., sing.}}{Ἑλλάδ᾽ ἐς αἶαν πιστὰ καλεῖται,}
\verso{3}{\notagramatical{ἀφνεῶν}{ἀφνεός/\allowbreak -όν}{rico}{adj., gen., pl., neut.}}{καὶ τῶν ἀφνεῶν καὶ πολυχρύσων}
\verso{4}{\notagramatical{ἑδράνων}{ἕδρανον, -ου, τό}{assento, residência}{gen., pl.}\crlf \notagramatical{φύλακες}{φύλαξ, φύλακος, ὁ}{guardião}{nom., pl.}\crlf \preposicaopreposta{κατὰ}{ac.}{segundo}\crlf \notagramatical{πρεσβείαν}{πρεσβεία, -ας, ἡ}{ancianidade}{ac., sing.}}{ἑδράνων φύλακες, κατὰ πρεσβείαν}
\verso{5}{}{οὓς αὐτὸς ἄναξ Ξέρξης βασιλεὺς}
\verso{6}{}{Δαρειογενὴς}
\verso{7}{\notagramatical{εἵλετο}{αἱρέω}{escolher}{aor., ind., méd., 3ª s.}\crlf \notagramatical{ἐφορεύειν}{ἐφορεύω}{vigiar}{inf., pres., at.; + gen.: vigiar, governar algo; χώρας é o complemento}}{εἵλετο χώρας ἐφορεύειν.}
\verso{8}{}{ἀμφὶ δὲ νόστῳ τῷ βασιλείῳ}
\verso{9}{\notagramatical{στρατιᾶς}{στρατιά, -ᾶς, ἡ}{exército}{gen., sing.}\crlf \notagramatical{ἤδη}{ἤδη}{já}{adv.}}{καὶ πολυχρύσου στρατιᾶς ἤδη}
\verso{10}{\notagramatical{ἄγαν}{ἄγαν}{muito, excessivamente}{adv.}}{κακόμαντις ἄγαν ὀρσολοπεῖται}
\verso{11}{\notagramatical{ἔσωθεν}{ἔσωθεν}{por dentro}{adv.}}{θυμὸς ἔσωθεν.}
\verso{12}{\notagramatical{ἰσχὺς}{ἰσχύς, -ύος, ἡ}{força}{nom., sing.}}{πᾶσα γὰρ ἰσχὺς Ἀσιατογενὴς}
\verso{13}{\notagramatical{ᾤχωκε}{οἴχομαι}{}{perf., ind., at., 3ª s.}\crlf \notagramatical{βαΰζει}{βαύζω}{latir, reclamar}{pres., ind., at., 3ª s.}}{ᾤχωκε, νέον δ᾽ ἄνδρα βαΰζει,}
\verso{14}{}{κοὔτε τις ἄγγελος οὔτε τις ἱππεὺς}
\verso{15}{\notagramatical{ἄστυ}{ἄστυ, ἄστεως, τό}{cidade}{ac., sing.}\crlf \notagramatical{ἀφικνεῖται}{ἀφικνέομαι}{chegar}{pres., ind., méd., 3ª s.}}{ἄστυ τὸ Περσῶν ἀφικνεῖται:}
\verso{16}{}{οἵτε τὸ Σούσων ἠδ᾽ Ἀγβατάνων}
\verso{17}{\notagramatical{παλαιὸν}{παλαιός/\allowbreak -ά/\allowbreak -όν}{antigo}{adj., ac., sing., neut.}}{καὶ τὸ παλαιὸν Κίσσιον ἕρκος}
\verso{18}{\notagramatical{προλιπόντες}{προλείπω}{deixar para trás}{part., aor., at., nom., pl., masc.}\crlf \notagramatical{ἔβαν}{βαίνω}{ir, avançar}{aor., ind., at., 3ª p.}\crlf \notagramatical{τοὶ}{οἱ}{}{art., nom. pl. masc.}\crlf \preposicaolemapreposta{ἐφ᾽}{ἐπί}{gen.}{sobre, a bordo de}}{προλιπόντες ἔβαν, τοὶ μὲν ἐφ᾽ ἵππων,}
\verso{19}{}{τοὶ δ᾽ ἐπὶ ναῶν, πεζοί τε βάδην}
\verso{20}{\notagramatical{στῖφος}{στῖφος, -ους, τό}{massa compacta}{ac., sing.}\crlf \notagramatical{παρέχοντες}{παρέχω}{proporcionar}{part., pres., at., nom., pl., masc.}}{πολέμου στῖφος παρέχοντες:}
\verso{21}{}{οἷος Ἀμίστρης ἠδ᾽ Ἀρταφρένης}
\verso{22}{}{καὶ Μεγαβάτης ἠδ᾽ Ἀστάσπης,}
\verso{23}{\notagramatical{ταγοὶ}{ταγός, -οῦ, ὁ}{comandante}{nom., pl.}}{ταγοὶ Περσῶν,}
\verso{24}{\notagramatical{ὕποχοι}{ὕποχος/\allowbreak -ον}{subordinado}{adj., nom., pl., masc.}}{βασιλῆς βασιλέως ὕποχοι μεγάλου,}
\verso{25}{\notagramatical{σοῦνται}{σεύω}{impelir, lançar-se}{pres., ind., méd., 3ª p.}}{σοῦνται, στρατιᾶς πολλῆς ἔφοροι,}
\verso{26}{\notagramatical{τοξοδάμαντές}{τοξοδάμας, τοξοδάμαντος, ὁ}{arqueiro}{nom., pl.}\crlf \notagramatical{ἠδ᾽}{ἠδέ}{e}{conj.}\crlf \notagramatical{ἱπποβάται}{ἱπποβάτης, -ου, ὁ}{cavaleiro}{nom., pl.}}{τοξοδάμαντές τ᾽ ἠδ᾽ ἱπποβάται,}
\verso{27}{}{φοβεροὶ μὲν ἰδεῖν, δεινοὶ δὲ μάχην}
\verso{28}{\notagramatical{εὐτλήμονι}{εὐτλήμων/\allowbreak -ον}{resistente, intrépido}{adj., dat. sing. fem.}}{ψυχῆς εὐτλήμονι δόξῃ:}
\verso{29}{\notagramatical{ἱππιοχάρμης}{ἱππιοχάρμης, -ου, ὁ}{combatente a cavalo}{nom., sing.}}{Ἀρτεμβάρης θ᾽ ἱππιοχάρμης}
\verso{30}{}{καὶ Μασίστης, ὅ τε τοξοδάμας}
\verso{31}{\notagramatical{ἐσθλὸς}{ἐσθλός/\allowbreak -ή/\allowbreak -όν}{valoroso}{adj., nom., sing., masc.}}{ἐσθλὸς Ἰμαῖος, Φαρανδάκης θ᾽,}
\verso{32}{\notagramatical{ἐλατὴρ}{ἐλατήρ, ἐλατῆρος, ὁ}{condutor}{nom., sing.}}{ἵππων τ᾽ ἐλατὴρ Σοσθάνης.}
\verso{33}{\notagramatical{πολυθρέμμων}{πολυθρέμμων/\allowbreak -ον}{que nutre muitos}{adj., nom., sing., masc.}}{ἄλλους δ᾽ ὁ μέγας καὶ πολυθρέμμων}
\verso{34}{}{Νεῖλος ἔπεμψεν: Σουσισκάνης,}
\verso{35}{}{Πηγασταγὼν Αἰγυπτογενής,}
\verso{36}{\notagramatical{ἄρχων}{ἄρχω}{}{part., pres., at., nom., sing., masc.}}{ὅ τε τῆς ἱερᾶς Μέμφιδος ἄρχων}
\verso{37}{\notagramatical{ὠγυγίους}{ὠγύγιος/\allowbreak -α/\allowbreak -ον}{antiquíssimo}{adj., ac., pl., fem.}}{μέγας Ἀρσάμης, τάς τ᾽ ὠγυγίους}
\verso{38}{\notagramatical{ἐφέπων}{ἐφέπω}{governar}{part., pres., at., nom., sing., masc.}}{Θήβας ἐφέπων Ἀριόμαρδος,}
\verso{39}{\notagramatical{ἑλειοβάται}{ἑλειοβάτης, -ου, ὁ}{habitante dos pântanos}{nom., pl.}\crlf \notagramatical{ἐρέται}{ἐρέτης, -ου, ὁ}{remador}{nom., pl.}}{καὶ ἑλειοβάται ναῶν ἐρέται}
\verso{40}{}{δεινοὶ πλῆθός τ᾽ ἀνάριθμοι.}
\verso{41}{\notagramatical{ἁβροδιαίτων}{ἁβροδίαιτος/\allowbreak -ον}{de vida luxuosa}{adj., gen., pl., masc.}}{ἁβροδιαίτων δ᾽ ἕπεται Λυδῶν}
\verso{42}{\notagramatical{ὄχλος}{ὄχλος, -ου, ὁ}{multidão}{nom., sing.}\crlf \notagramatical{ἐπίπαν}{ἐπίπαν}{inteiramente}{adv.}\crlf \notagramatical{ἠπειρογενὲς}{ἠπειρογενής/\allowbreak -ές}{nascido no continente}{adj., ac., sing., neut.}}{ὄχλος, οἵτ᾽ ἐπίπαν ἠπειρογενὲς}
\verso{43}{}{κατέχουσιν ἔθνος, τοὺς Μιτρογαθὴς}
\verso{44}{\notagramatical{δίοποι}{δίοπος, -ου, ὁ}{supervisor}{nom., pl.}}{Ἀρκτεύς τ᾽ ἀγαθός, βασιλῆς δίοποι,}
\verso{45}{\notagramatical{ἐπόχους}{ἔποχος/\allowbreak -ον}{montado, embarcado}{adj., ac., pl., masc.}}{καὶ πολύχρυσοι Σάρδεις ἐπόχους}
\verso{46}{\notagramatical{ἅρμασιν}{ἅρμα, ἅρματος, τό}{carro de guerra}{dat., pl.}}{πολλοῖς ἅρμασιν ἐξορμῶσιν,}
\verso{47}{\notagramatical{δίρρυμά}{δίρρυμος/\allowbreak -ον}{com duas lanças}{adj., ac., pl., neut.}\crlf \notagramatical{τρίρρυμα}{τρίρρυμος/\allowbreak -ον}{com três lanças}{adj., ac., pl., neut.}\crlf \notagramatical{τέλη}{τέλος, -ους, τό}{destacamento}{ac. pl.}}{δίρρυμά τε καὶ τρίρρυμα τέλη,}
\verso{48}{}{φοβερὰν ὄψιν προσιδέσθαι.}
\verso{49}{\notagramatical{στεῦται}{στεῦμαι}{vangloriar-se}{pres., ind., méd., 3ª s.; + inf.: vangloriar-se de que fará; ἀμφιβαλεῖν é inf. fut.}\crlf \notagramatical{πελάτης}{πελάτης, -ου, ὁ}{habitante}{nom. sing.}}{στεῦται δ᾽ ἱεροῦ Τμώλου πελάτης}
\verso{50}{}{ζυγὸν ἀμφιβαλεῖν δούλιον Ἑλλάδι,}
\verso{51}{\notagramatical{ἄκμονες}{ἄκμων, ἄκμονος, ὁ}{bigorna}{nom., pl.}}{Μάρδων, Θάρυβις, λόγχης ἄκμονες,}
\verso{52}{}{καὶ ἀκοντισταὶ Μυσοί: Βαβυλὼν δ᾽}
\verso{53}{\notagramatical{πολύχρυσος}{πολύχρυσος/\allowbreak -ον}{rico em ouro}{adj., nom., sing., fem.}}{ἡ πολύχρυσος πάμμικτον ὄχλον}
\verso{54}{\notagramatical{σύρδην}{σύρδην}{em massa}{adv.}}{πέμπει σύρδην, ναῶν τ᾽ ἐπόχους}
\verso{55}{\notagramatical{τοξουλκῷ}{τοξουλκός/\allowbreak -όν}{que puxa o arco}{adj., dat., sing., neut.}\crlf \notagramatical{λήματι}{λῆμα, λήματος, τό}{ânimo, coragem}{dat., sing.}\crlf \notagramatical{πιστούς}{πιστός/\allowbreak -ή/\allowbreak -όν}{fiel}{adj., ac., pl., masc.}}{καὶ τοξουλκῷ λήματι πιστούς:}
\verso{56}{}{τὸ μαχαιροφόρον τ᾽ ἔθνος ἐκ πάσης}
\verso{57}{}{Ἀσίας ἕπεται}
\verso{58}{\notagramatical{δειναῖς}{δεινός/\allowbreak -ή/\allowbreak -όν}{temível}{adj., dat., pl., fem.}\crlf \preposicaopreposta{ὑπὸ}{dat.}{sob}\crlf \notagramatical{πομπαῖς}{πομπή, -ῆς, ἡ}{escolta}{dat., pl.}}{δειναῖς βασιλέως ὑπὸ πομπαῖς.}
\verso{59}{}{τοιόνδ᾽ ἄνθος Περσίδος αἴας}
\verso{60}{\notagramatical{οἴχεται}{οἴχομαι}{}{pres., ind., méd., 3ª s.}}{οἴχεται ἀνδρῶν,}
\verso{61}{\preposicaoposposta{πέρι}{ac.}{por, a respeito de}\crlf \notagramatical{χθὼν}{χθών, χθονός, ἡ}{terra}{nom., sing.}}{οὓς πέρι πᾶσα χθὼν Ἀσιᾶτις}
\verso{62}{\notagramatical{πόθῳ}{πόθος, -ου, ὁ}{saudade}{dat. sing.; causa: por saudade}}{θρέψασα πόθῳ στένεται μαλερῷ,}
\verso{63}{\notagramatical{τοκέες}{τοκεύς, τοκέως, ὁ}{progenitor}{nom., pl.}\crlf \notagramatical{ἄλοχοί}{ἄλοχος, -ου, ἡ}{esposa}{nom. pl.}\crlf \notagramatical{ἡμερολεγδὸν}{ἡμερολεγδόν}{contando os dias}{adv.}}{τοκέες τ᾽ ἄλοχοί θ᾽ ἡμερολεγδὸν}
\verso{64}{}{τείνοντα χρόνον τρομέονται.}
\pausaestrofe
\verso{65}{\notagramatical{πεπέρακεν}{περάω}{atravessar}{perf., ind., at., 3ª s.}}{πεπέρακεν μὲν ὁ περσέπτολις ἤδη}
\verso{66}{\notagramatical{στρατὸς}{στρατός, -οῦ, ὁ}{exército}{nom., sing.}}{βασίλειος στρατὸς εἰς ἀν-}
\verso{67}{}{τίπορον γείτονα χώραν,}
\verso{68}{\notagramatical{πορθμόν}{πορθμός, -οῦ, ὁ}{estreito marítimo}{ac. sing.}}{λινοδέσμῳ σχεδίᾳ πορ-}
\verso{69}{\notagramatical{ἀμείψας}{ἀμείβω}{passar, mudar}{part., aor., at., nom., sing., masc.}}{θμὸν ἀμείψας}
\verso{70}{}{Ἀθαμαντίδος Ἕλλας,}
\verso{71}{\notagramatical{πολύγομφον}{πολύγομφος/\allowbreak -ον}{preso por muitas cavilhas}{adj., ac. sing. neut.}}{πολύγομφον ὅδισμα}
\verso{72}{\notagramatical{ἀμφιβαλὼν}{ἀμφιβάλλω}{lançar em torno}{part., aor., at., nom., sing., masc.}\crlf \notagramatical{αὐχένι}{αὐχήν, αὐχένος, ὁ}{pescoço}{dat., sing.}}{ζυγὸν ἀμφιβαλὼν αὐχένι πόντου.}
\pausaestrofe
\verso{73}{}{πολυάνδρου δ᾽ Ἀσίας θούριος ἄρχων}
\verso{74}{\notagramatical{ποιμανόριον}{ποιμανόριον, -ου, τό}{rebanho humano}{ac. sing.}}{ἐπὶ πᾶσαν χθόνα ποιμα-}
\verso{75}{}{νόριον θεῖον ἐλαύνει}
\verso{76}{\notagramatical{διχόθεν}{διχόθεν}{em duas divisões}{adv.}\crlf \notagramatical{πεζονόμον}{πεζονόμος/\allowbreak -ον}{que avança por terra}{adj., ac. sing. neut.}}{διχόθεν, πεζονόμον τ᾽ ἔκ}
\verso{77}{}{τε θαλάσσας,}
\verso{78}{\notagramatical{ἐχυροῖσι}{ἐχυρός/\allowbreak -ά/\allowbreak -όν}{forte}{adj., dat. pl. masc.; dat. com πεποιθώς: confiando em}\crlf \notagramatical{πεποιθὼς}{πείθω}{persuadir; perf.: confiar}{part., perf., at., nom., sing., masc.; + dat.: confiar em; perfeito com valor de estado}}{ἐχυροῖσι πεποιθὼς}
\verso{79}{}{στυφελοῖς ἐφέταις, χρυ-}
\verso{80}{\notagramatical{γενεᾶς}{γενεά, -ᾶς, ἡ}{linhagem}{gen. sing.; χρυσογόνου γενεᾶς: de linhagem nascida do ouro}}{σογόνου γενεᾶς ἰσόθεος φώς.}
\pausaestrofe
\verso{81}{\notagramatical{κυανοῦν}{κυάνεος/\allowbreak -α/\allowbreak -ον}{escuro}{adj., ac. sing. neut.; forma contraída}}{κυανοῦν δ᾽ ὄμμασι λεύσσων}
\verso{82}{\notagramatical{φονίου}{φόνιος/\allowbreak -α/\allowbreak -ον}{sanguinário}{adj., gen. sing. masc.}\crlf \notagramatical{δέργμα}{δέργμα, δέργματος, τό}{olhar}{ac., sing.}}{φονίου δέργμα δράκοντος,}
\verso{83}{}{πολύχειρ καὶ πολυναύτας,}
\verso{84}{}{Σύριόν θ᾽ ἅρμα διώκων,}
\verso{85}{\notagramatical{δουρικλύτοις}{δουρικλυτός/\allowbreak -ή/\allowbreak -όν}{célebre pela lança}{adj., dat., pl., masc.}}{ἐπάγει δουρικλύτοις ἀν-}
\verso{86}{}{δράσι τοξόδαμνον Ἄρη.}
\pausaestrofe
\verso{87}{}{δόκιμος δ᾽ οὔτις ὑποστὰς}
\verso{88}{}{μεγάλῳ ῥεύματι φωτῶν}
\verso{89}{\notagramatical{ἐχυροῖς}{ἐχυρός/\allowbreak -ά/\allowbreak -όν}{}{adj., dat. pl. neut.}\crlf \notagramatical{ἕρκεσιν}{ἕρκος, -ους, τό}{recinto, defesa}{dat., pl.}\crlf \notagramatical{εἴργειν}{εἴργω}{impedir, conter}{inf., pres., at.}}{ἐχυροῖς ἕρκεσιν εἴργειν}
\verso{90}{}{ἄμαχον κῦμα θαλάσσας:}
\verso{91}{\notagramatical{ἀπρόσοιστος}{ἀπρόσοιστος/\allowbreak -ον}{inabordável}{adj., nom., sing., masc.}}{ἀπρόσοιστος γὰρ ὁ Περσᾶν}
\verso{92}{\notagramatical{ἀλκίφρων}{ἀλκίφρων/\allowbreak -ον}{de ânimo valente}{adj., nom., sing., masc.}}{στρατὸς ἀλκίφρων τε λαός.}
\pausaestrofe
\verso{93}{\notagramatical{δολόμητιν}{δολόμητις/\allowbreak -ι}{de intenção enganadora}{adj., ac. sing. fem.}\crlf \notagramatical{ἀπάταν}{ἀπάτη, -ης, ἡ}{engano}{ac. sing.; forma dórica}}{δολόμητιν δ᾽ ἀπάταν θεοῦ}
\verso{94}{}{τίς ἀνὴρ θνατὸς ἀλύξει;}
\verso{95}{\notagramatical{κραιπνῷ}{κραιπνός/\allowbreak -ή/\allowbreak -όν}{veloz}{adj., dat., sing., masc.}}{τίς ὁ κραιπνῷ ποδὶ πήδη-}
\verso{96}{\notagramatical{εὐπετῶς}{εὐπετῶς}{facilmente}{adv.}\crlf \notagramatical{ἀνᾴσσων}{ἀναΐσσω}{saltar}{part., pres., at., nom., sing., masc.}}{μα τόδ᾽ εὐπετῶς ἀνᾴσσων;}
\pausaestrofe
\verso{97}{\notagramatical{παρασαίνει}{παρασαίνω}{atrair com lisonjas}{pres., ind., at., 3ª s.}}{φιλόφρων γὰρ παρασαίνει}
\verso{98}{\notagramatical{βροτὸν}{βροτός, -οῦ, ὁ}{mortal}{ac., sing.}}{βροτὸν εἰς ἄρκυας Ἄτα,}
\verso{99}{\notagramatical{τόθεν}{τόθεν}{de lá}{adv.}\crlf \notagramatical{ὕπερθέν}{ὕπερθεν}{por cima}{adv.}}{τόθεν οὐκ ἔστιν ὕπερθέν}
\verso{100}{\notagramatical{νιν}{νιν}{o, a}{pron., ac. sing.; aqui, o homem capturado}\crlf \notagramatical{ἄνατον}{ἄνατος/\allowbreak -ον}{sem dano}{adj., ac. sing. masc.}\crlf \notagramatical{ἐξαλύξαι}{ἐξαλύσκω}{escapar de}{inf., aor., at.}}{νιν ἄνατον ἐξαλύξαι.}
\verso{101}{}{θεόθεν γὰρ κατὰ Μοῖρ᾽}
\verso{102}{\notagramatical{ἐκράτησεν}{κρατέω}{prevalecer}{aor., ind., at., 3ª s.}}{ἐκράτησεν τὸ παλαι-}
\verso{103}{\notagramatical{ἐπέσκηψε}{ἐπισκήπτω}{impor, incumbir}{aor., ind., at., 3ª s.}}{όν, ἐπέσκηψε δὲ Πέρσαις}
\verso{104}{\notagramatical{πυργοδαΐκτους}{πυργοδάικτος/\allowbreak -ον}{que destroem torres}{adj., ac. pl. masc.}}{πολέμους πυργοδαΐκτους}
\verso{105}{\notagramatical{διέπειν}{διέπω}{dirigir}{inf., pres., at.}}{διέπειν ἱππιοχάρμας}
\verso{106}{}{τε κλόνους}
\verso{107}{\notagramatical{ἀναστάσεις}{ἀνάστασις, ἀναστάσεως, ἡ}{devastação}{ac. pl.; aqui, destruição de cidades}}{πόλεών τ᾽ ἀναστάσεις.}
\pausaestrofe
\verso{108}{\notagramatical{ἔμαθον}{μανθάνω}{aprender}{aor., ind., at., 3ª p.; + inf.: aprender a}\crlf \notagramatical{εὐρυπόροιο}{εὐρύπορος/\allowbreak -ον}{de amplas vias}{adj., gen. sing. fem.; gen. poético em -οιο}}{ἔμαθον δ᾽ εὐρυπόροι-}
\verso{109}{}{ο θαλάσσας πολιαι-}
\verso{110}{\notagramatical{λάβρῳ}{λάβρος/\allowbreak -α/\allowbreak -ον}{impetuoso}{adj., dat. sing. neut.}}{νομένας πνεύματι λάβρῳ}
\verso{111}{\notagramatical{ἄλσος}{ἄλσος, -ους, τό}{bosque}{ac. sing.; aqui, o mar figurado como bosque}}{ἐσορᾶν πόντιον ἄλσος,}
\verso{112}{\notagramatical{πίσυνοι}{πίσυνος/\allowbreak -ον}{confiante}{adj., nom. pl. masc.; + dat.: confiando em}\crlf \notagramatical{λεπτοδόμοις}{λεπτοδόμος/\allowbreak -ον}{de delicada feitura}{adj., dat. pl. neut.}\crlf \notagramatical{πείσμασι}{πεῖσμα, πείσματος, τό}{cordame}{dat. pl.}}{πίσυνοι λεπτοδόμοις πεί-}
\verso{113}{}{σμασι λα-}
\verso{114}{}{οπόροις τε μαχαναῖς.}
\pausaestrofe
\verso{115}{\notagramatical{μελαγχίτων}{μελαγχίτων/\allowbreak -ον}{vestido de negro}{adj., nom. sing. fem.}}{ταῦτά μοι μελαγχίτων}
\verso{116}{\notagramatical{ἀμύσσεται}{ἀμύσσω}{dilacerar}{pres., ind., méd./pass., 3ª s.}}{φρὴν ἀμύσσεται φόβῳ,}
\verso{117}{\notagramatical{ὀᾶ}{ὀᾶ}{ai!}{interj.}\crlf \notagramatical{στρατεύματος}{στράτευμα, στρατεύματος, τό}{exército}{gen., sing.}}{ὀᾶ, Περσικοῦ στρατεύματος}
\verso{118}{}{τοῦδε, μὴ πόλις πύθη-}
\verso{119}{\notagramatical{κένανδρον}{κένανδρος/\allowbreak -ον}{sem homens}{adj., ac., sing., neut.}}{ται κένανδρον μέγ᾽ ἄστυ Σουσίδος,}
\pausaestrofe
\verso{120}{\notagramatical{πόλισμα}{πόλισμα, πολίσματος, τό}{cidade}{nom. sing. neut.}}{καὶ τὸ Κισσίων πόλισμ᾽}
\verso{121}{\notagramatical{ᾄσεται}{ᾄδω}{cantar}{fut., ind., méd., 3ª s.}}{ἀντίδουπον ᾄσεται,}
\verso{122}{\notagramatical{γυναικοπληθής}{γυναικοπληθής/\allowbreak -ές}{repleto de mulheres}{adj., nom. sing. masc.}}{ὀᾶ, τοῦτ᾽ ἔπος γυναικοπλη-}
\verso{123}{\notagramatical{ὅμιλος}{ὅμιλος, -ου, ὁ}{multidão}{nom. sing.}}{θὴς ὅμιλος ἀπύων,}
\verso{125}{}{βυσσίνοις δ᾽ ἐν πέπλοις πέσῃ λακίς.}
\pausaestrofe
\verso{126}{\notagramatical{ἱππηλάτας}{ἱππηλάτης, -ου, ὁ}{cavaleiro}{nom. sing.; forma dórica em -ας}}{πᾶς γὰρ ἱππηλάτας}
\verso{127}{\notagramatical{πεδοστιβὴς}{πεδοστιβής/\allowbreak -ές}{que pisa o chão}{adj., nom., sing., masc.}}{καὶ πεδοστιβὴς λεὼς}
\verso{128}{\notagramatical{σμῆνος}{σμῆνος, -ους, τό}{enxame}{nom./ac. sing.; σμῆνος ὥς: como um enxame}\crlf \notagramatical{ἐκλέλοιπεν}{ἐκλείπω}{abandonar}{perf., ind., at., 3ª s.; perf.: abandonou e permanece ausente}}{σμῆνος ὣς ἐκλέλοιπεν μελισ-}
\verso{129}{}{σᾶν σὺν ὀρχάμῳ στρατοῦ,}
\verso{130}{\notagramatical{ἀμφίζευκτον}{ἀμφίζευκτος/\allowbreak -ον}{ligado a ambas as margens}{adj., ac. sing. masc.}\crlf \notagramatical{ἐξαμείψας}{ἐξαμείβω}{atravessar}{part., aor., at., nom., sing., masc.}}{τὸν ἀμφίζευκτον ἐξαμείψας}
\verso{131}{}{ἀμφοτέρας ἅλιον}
\verso{132}{\notagramatical{πρῶνα}{πρών, πρωνός, ὁ}{cabo}{ac., sing.}}{πρῶνα κοινὸν αἴας.}
\pausaestrofe
\verso{133}{\notagramatical{λέκτρα}{λέκτρον, -ου, τό}{leito conjugal}{nom. pl.}}{λέκτρα δ᾽ ἀνδρῶν πόθῳ}
\verso{134}{\notagramatical{πίμπλαται}{πίμπλημι}{encher}{pres., ind., pass., 3ª s.}}{πίμπλαται δακρύμασιν:}
\verso{135}{\notagramatical{ἑκάστα}{ἕκαστος/\allowbreak -η/\allowbreak -ον}{cada um}{adj., nom. sing. fem.; forma dórica}}{Περσίδες δ᾽ ἁβροπενθεῖς ἑκά-}
\verso{136}{\notagramatical{φιλάνορι}{φιλάνωρ/\allowbreak φίλανορ}{que ama o marido}{adj., dat. sing. masc.}}{στα πόθῳ φιλάνορι}
\verso{137}{\notagramatical{αἰχμάεντα}{αἰχμήεις/\allowbreak -εσσα/\allowbreak -εν}{armado de lança}{adj., ac. sing. masc.}\crlf \notagramatical{θοῦρον}{θοῦρος/\allowbreak -α/\allowbreak -ον}{impetuoso}{adj., ac. sing. masc.}}{τὸν αἰχμάεντα θοῦρον εὐνα-}
\verso{138}{}{τῆρ᾽ ἀποπεμψαμένα}
\verso{139}{\notagramatical{λείπεται}{λείπω}{deixar, ficar atrás}{pres., ind., méd./pass., 3ª s.}\crlf \notagramatical{μονόζυξ}{μονόζυξ}{sem par}{adj., nom. sing. fem.}}{λείπεται μονόζυξ.}
\pausaestrofe
\verso{140}{}{ἀλλ᾽ ἄγε, Πέρσαι, τόδ᾽ ἐνεζόμενοι}
\verso{141}{\notagramatical{στέγος}{στέγος, -ους, τό}{edifício, abrigo}{ac. sing.; complemento de ἐνεζόμενοι}}{στέγος ἀρχαῖον,}
\verso{142}{\notagramatical{κεδνὴν}{κεδνός/\allowbreak -ή/\allowbreak -όν}{prudente, fiel}{adj., ac., sing., fem.}}{φροντίδα κεδνὴν καὶ βαθύβουλον}
\verso{143}{}{θώμεθα, χρεία δὲ προσήκει,}
\verso{144}{\notagramatical{ἄρα}{ἄρα}{então, afinal}{partícula; reforça a pergunta}\crlf \notagramatical{πράσσει}{πράσσω}{fazer, realizar}{pres., ind., at., 3ª s.}}{πῶς ἄρα πράσσει Ξέρξης βασιλεὺς}
\verso{145}{}{Δαρειογενής,}
\verso{146}{\notagramatical{πατρωνύμιον}{πατρωνύμιος/\allowbreak -ον}{nomeado pelo pai}{adj., nom., sing., neut.}}{τὸ πατρωνύμιον γένος ἡμέτερον:}
\verso{147}{\notagramatical{ῥῦμα}{ῥῦμα, ῥύματος, τό}{defesa}{nom., sing.}}{πότερον τόξου ῥῦμα τὸ νικῶν,}
\verso{148}{\notagramatical{δορικράνου}{δορίκρανος/\allowbreak -ον}{de ponta de lança}{adj., gen. sing. fem.}}{ἢ δορικράνου}
\verso{149}{\notagramatical{κεκράτηκεν}{κρατέω}{}{perf., ind., at., 3ª s.; perf.: obteve a vitória e a conserva}}{λόγχης ἰσχὺς κεκράτηκεν.}
\verso{150}{\notagramatical{ἥδε}{ὅδε/\allowbreak ἥδε/\allowbreak τόδε}{esta}{pron. demonstr., nom. sing. fem.; aqui, a rainha que se aproxima}}{ἀλλ᾽ ἥδε θεῶν ἴσον ὀφθαλμοῖς}
\verso{151}{}{φάος ὁρμᾶται μήτηρ βασιλέως,}
\verso{152}{\notagramatical{προσπίτνω}{προσπίπτω}{prostrar-se diante}{pres., ind., at., 1ª s.}}{βασίλεια δ᾽ ἐμή, προσπίτνω:}
\verso{153}{\notagramatical{προσφθόγγοις}{πρόσφθογγος/\allowbreak -ον}{de saudação}{adj., dat. pl. masc.}\crlf \notagramatical{χρεὼν}{χρεών}{é preciso}{impessoal; + ac. e inf.: é preciso que alguém faça}}{καὶ προσφθόγγοις δὲ χρεὼν αὐτὴν}
\verso{154}{}{πάντας μύθοισι προσαυδᾶν.}
\parte{Primeiro Episódio}{155–531}
\pausaestrofe
\verso{155}{\notagramatical{ὑπερτάτη}{ὑπέρτατος/\allowbreak -η/\allowbreak -ον}{supremo}{adj., superl., voc. sing. fem.}}{ὦ βαθυζώνων ἄνασσα Περσίδων ὑπερτάτη,}
\verso{156}{}{μῆτερ ἡ Ξέρξου γεραιά, χαῖρε, Δαρείου γύναι:}
\verso{157}{\notagramatical{ἔφυς}{φύω}{ser por natureza}{aor., ind., at., 2ª s.; aqui: foste}}{θεοῦ μὲν εὐνάτειρα Περσῶν, θεοῦ δὲ καὶ μήτηρ ἔφυς,}
\verso{158}{\notagramatical{εἴ τι μὴ}{εἰ μή}{a menos que}{εἴ τι μή: salvo se, por acaso}\crlf \notagramatical{δαίμων}{δαίμων, δαίμονος, ὁ}{divindade, destino}{nom., sing.}\crlf \notagramatical{μεθέστηκε}{μεθίστημι}{mudar de posição}{perf., ind., at., 3ª s.; + dat.: mudar de disposição para com; aqui, στρατῷ}}{εἴ τι μὴ δαίμων παλαιὸς νῦν μεθέστηκε στρατῷ.}
\falante{ΑΤΟΣΣΑ}
\verso{159}{\notagramatical{ἱκάνω}{ἱκάνω}{chegar}{pres., ind., at., 1ª s.}\crlf \notagramatical{δόμους}{δόμος, -ου, ὁ}{casa, palácio}{ac., pl.}}{ταῦτα δὴ λιποῦσ᾽ ἱκάνω χρυσεοστόλμους δόμους}
\verso{160}{\desdobramento{κἀμὸν}{καὶ ἐμόν}\crlf \notagramatical{εὐνατήριον}{εὐνατήριον, -ου, τό}{câmara nupcial}{ac., sing.}}{καὶ τὸ Δαρείου τε κἀμὸν κοινὸν εὐνατήριον.}
\verso{161}{}{καί με καρδίαν ἀμύσσει φροντίς: ἐς δ᾽ ὑμᾶς ἐρῶ}
\verso{162}{\notagramatical{ἐμαυτῆς}{ἐμαυτοῦ/\allowbreak ἐμαυτῆς}{de mim mesma}{pron. reflex., gen. sing. fem.; complemento de ἀδείμαντος: por mim mesma}\crlf \notagramatical{οὖσ᾽}{εἰμί}{ser, estar}{part., pres., nom. sing. fem.}\crlf \notagramatical{ἀδείμαντος}{ἀδείμαντος/\allowbreak -ον}{sem temor}{adj., nom. sing. fem.; + gen.: sem temor por; οὐδαμῶς: de modo algum}}{μῦθον, οὐδαμῶς ἐμαυτῆς οὖσ᾽ ἀδείμαντος, φίλοι,}
\verso{163}{}{μὴ μέγας πλοῦτος κονίσας οὖδας ἀντρέψῃ ποδὶ}
\verso{164}{}{ὄλβον, ὃν Δαρεῖος ἦρεν οὐκ ἄνευ θεῶν τινός.}
\verso{165}{}{ταῦτά μοι διπλῆ μέριμνα φραστός ἐστιν ἐν φρεσίν,}
\verso{166}{\notagramatical{ἀνάνδρων}{ἄνανδρος/\allowbreak -ον}{sem homens}{adj., gen., pl., neut.}\crlf \notagramatical{πλῆθος}{πλῆθος, -ους, τό}{multidão}{ac., sing.}}{μήτε χρημάτων ἀνάνδρων πλῆθος ἐν τιμῇ σέβειν}
\verso{167}{}{μήτ᾽ ἀχρημάτοισι λάμπειν φῶς ὅσον σθένος πάρα.}
\verso{168}{\notagramatical{πλοῦτός}{πλοῦτος, -ου, ὁ}{riqueza}{nom., sing.}\crlf \notagramatical{ἀμεμφής}{ἀμεμφής/\allowbreak -ές}{irrepreensível}{adj., nom., sing., masc.}}{ἔστι γὰρ πλοῦτός γ᾽ ἀμεμφής, ἀμφὶ δ᾽ ὀφθαλμῷ φόβος:}
\verso{169}{}{ὄμμα γὰρ δόμων νομίζω δεσπότου παρουσίαν.}
\verso{170}{\preposicaopreposta{πρὸς}{ac.}{diante de}\crlf \notagramatical{ἐχόντων}{ἔχω}{}{part., pres., at., gen., pl., masc.}}{πρὸς τάδ᾽ ὡς οὕτως ἐχόντων τῶνδε, σύμβουλοι λόγου}
\verso{171}{\notagramatical{γένεσθε}{γίγνομαι}{}{aor., imper., méd., 2ª p.}\crlf \notagramatical{πιστώματα}{πίστωμα, πιστώματος, τό}{garantia}{nom., pl.}}{τοῦδέ μοι γένεσθε, Πέρσαι, γηραλέα πιστώματα:}
\verso{172}{}{πάντα γὰρ τὰ κέδν᾽ ἐν ὑμῖν ἐστί μοι βουλεύματα.}
\falante{ΧΟΡΟΣ}
\verso{173}{\notagramatical{ἴσθι}{οἶδα}{}{perf., imper., at., 2ª s.}\crlf \notagramatical{φράσαι}{φράζω}{explicar, indicar}{inf., aor., at.}}{εὖ τόδ᾽ ἴσθι, γῆς ἄνασσα τῆσδε, μή σε δὶς φράσαι}
\verso{174}{\notagramatical{θέλῃ}{ἐθέλω}{querer}{pres., subj., at., 3ª s.; relativa com ἄν + subj.: tudo o que a capacidade permitir}}{μήτ᾽ ἔπος μήτ᾽ ἔργον ὧν ἂν δύναμις ἡγεῖσθαι θέλῃ:}
\verso{175}{}{εὐμενεῖς γὰρ ὄντας ἡμᾶς τῶνδε συμβούλους καλεῖς.}
\falante{ΑΤΟΣΣΑ}
\verso{176}{}{πολλοῖς μὲν αἰεὶ νυκτέροις ὀνείρασιν}
\verso{177}{}{ξύνειμ᾽, ἀφ᾽ οὗπερ παῖς ἐμὸς στείλας στρατὸν}
\verso{178}{\notagramatical{πέρσαι}{πέρθω}{devastar}{inf., aor., at.}\crlf \notagramatical{θέλων}{ἐθέλω}{}{part., pres., at., nom., sing., masc.}}{Ἰαόνων γῆν οἴχεται πέρσαι θέλων:}
\verso{179}{}{ἀλλ᾽ οὔτι πω τοιόνδ᾽ ἐναργὲς εἰδόμην}
\verso{180}{\notagramatical{εὐφρόνης}{εὐφρόνη, -ης, ἡ}{noite}{gen., sing.}\crlf \notagramatical{λέξω}{λέγω}{}{fut., ind., at., 1ª s.}}{ὡς τῆς πάροιθεν εὐφρόνης: λέξω δέ σοι.}
\verso{181}{}{ἐδοξάτην μοι δύο γυναῖκ᾽ εὐείμονε,}
\verso{182}{\notagramatical{ἠσκημένη}{ἀσκέω}{adornar}{part., perf., pass., nom. sing. fem.}}{ἡ μὲν πέπλοισι Περσικοῖς ἠσκημένη,}
\verso{183}{\notagramatical{μολεῖν}{βλώσκω}{vir}{inf., aor., at.}}{ἡ δ᾽ αὖτε Δωρικοῖσιν, εἰς ὄψιν μολεῖν,}
\verso{184}{}{μεγέθει τε τῶν νῦν ἐκπρεπεστάτα πολύ,}
\verso{185}{\notagramatical{ἀμώμω}{ἄμωμος/\allowbreak -ον}{sem defeito}{adj., nom., dual, fem.}}{κάλλει τ᾽ ἀμώμω, καὶ κασιγνήτα γένους}
\verso{186}{\desdobramento{ταὐτοῦ}{τοῦ αὐτοῦ}}{ταὐτοῦ: πάτραν δ᾽ ἔναιον ἡ μὲν Ἑλλάδα}
\verso{187}{\notagramatical{λαχοῦσα}{λαγχάνω}{obter por sorte}{part., aor., at., nom., sing., fem.}\crlf \notagramatical{γαῖαν}{γαῖα, -ας, ἡ}{terra}{ac., sing.}\crlf \notagramatical{βάρβαρον}{βάρβαρος/\allowbreak -ον}{bárbaro, não grego}{adj., ac., sing., fem.}}{κλήρῳ λαχοῦσα γαῖαν, ἡ δὲ βάρβαρον.}
\verso{188}{}{τούτω στάσιν τιν᾽, ὡς ἐγὼ ᾽δόκουν ὁρᾶν,}
\verso{189}{\notagramatical{τεύχειν}{τεύχω}{produzir}{inf., pres., at.}}{τεύχειν ἐν ἀλλήλαισι: παῖς δ᾽ ἐμὸς μαθὼν}
\verso{190}{\desdobramento{κἀπράυνεν}{καὶ ἐπράυνεν}\crlf \preposicaoposposta{ὕπο}{dat.}{sob}}{κατεῖχε κἀπράυνεν, ἅρμασιν δ᾽ ὕπο}
\verso{191}{\preposicaopreposta{ἐπ᾽}{gen.}{sobre}}{ζεύγνυσιν αὐτὼ καὶ λέπαδν᾽ ἐπ᾽ αὐχένων}
\verso{192}{\desdobramento{χἠ}{καὶ ἡ}}{τίθησι. χἠ μὲν τῇδ᾽ ἐπυργοῦτο στολῇ}
\verso{193}{}{ἐν ἡνίαισί τ᾽ εἶχεν εὔαρκτον στόμα,}
\verso{194}{\notagramatical{ἐσφάδᾳζε}{σφαδάζω}{debater-se}{imperf., ind., at., 3ª s.}}{ἡ δ᾽ ἐσφάδᾳζε, καὶ χεροῖν ἔντη δίφρου}
\verso{195}{\notagramatical{διασπαράσσει}{διασπαράσσω}{despedaçar}{pres., ind., at., 3ª s.}\crlf \notagramatical{ξυναρπάζει}{συναρπάζω}{arrebatar}{pres., ind., at., 3ª s.}}{διασπαράσσει, καὶ ξυναρπάζει βίᾳ}
\verso{196}{}{ἄνευ χαλινῶν καὶ ζυγὸν θραύει μέσον.}
\verso{197}{\notagramatical{πίπτει}{πίπτω}{cair}{pres., ind., at., 3ª s.}}{πίπτει δ᾽ ἐμὸς παῖς, καὶ πατὴρ παρίσταται}
\verso{198}{\notagramatical{οἰκτείρων}{οἰκτείρω}{apiedar-se}{part., pres., at., nom., sing., masc.}}{Δαρεῖος οἰκτείρων σφε: τὸν δ᾽ ὅπως ὁρᾷ}
\verso{199}{\notagramatical{ῥήγνυσιν}{ῥήγνυμι}{rasgar}{pres., ind., at., 3ª s.}}{Ξέρξης, πέπλους ῥήγνυσιν ἀμφὶ σώματι.}
\verso{200}{\notagramatical{νυκτὸς}{νύξ, νυκτός, ἡ}{noite}{gen., sing.}\crlf \notagramatical{εἰσιδεῖν}{εἰσοράω}{olhar, ver}{inf., aor., at.}}{καὶ ταῦτα μὲν δὴ νυκτὸς εἰσιδεῖν λέγω.}
\verso{201}{}{ἐπεὶ δ᾽ ἀνέστην καὶ χεροῖν καλλιρρόου}
\verso{202}{\notagramatical{ἔψαυσα}{ψαύω}{tocar}{aor., ind., at., 1ª s.}}{ἔψαυσα πηγῆς, σὺν θυηπόλῳ χερὶ}
\verso{203}{\notagramatical{προσέστην}{προσίστημι}{aproximar-se de}{aor., ind., at., 1ª s.}}{βωμὸν προσέστην, ἀποτρόποισι δαίμοσιν}
\verso{204}{\notagramatical{θέλουσα}{ἐθέλω}{}{part., pres., at., nom., sing., fem.}\crlf \notagramatical{πέλανον}{πέλανος, -ου, ὁ}{oferenda líquida}{ac., sing.}}{θέλουσα θῦσαι πέλανον, ὧν τέλη τάδε.}
\verso{205}{}{ὁρῶ δὲ φεύγοντ᾽ αἰετὸν πρὸς ἐσχάραν}
\verso{206}{\notagramatical{ἄφθογγος}{ἄφθογγος/\allowbreak -ον}{sem voz}{adj., nom. sing. fem.}}{Φοίβου: φόβῳ δ᾽ ἄφθογγος ἐστάθην, φίλοι:}
\verso{207}{\notagramatical{κίρκον}{κίρκος, -ου, ὁ}{falcão}{ac., sing.}\crlf \notagramatical{εἰσορῶ}{εἰσοράω}{}{pres., ind., at., 1ª s.}}{μεθύστερον δὲ κίρκον εἰσορῶ δρόμῳ}
\verso{208}{}{πτεροῖς ἐφορμαίνοντα καὶ χηλαῖς κάρα}
\verso{209}{\notagramatical{τίλλονθ᾽}{τίλλω}{arrancar}{part., pres., at., ac. sing. masc.}\crlf \notagramatical{πτήξας}{πτήσσω}{encolher-se de medo}{part., aor., at., nom., sing., masc.}\crlf \notagramatical{δέμας}{δέμας, τό}{corpo; indecl.}{ac., sing.}}{τίλλονθ᾽: ὁ δ᾽ οὐδὲν ἄλλο γ᾽ ἢ πτήξας δέμας}
\verso{210}{}{παρεῖχε. ταῦτ᾽ ἔμοιγε δείματ᾽ εἰσιδεῖν,}
\verso{211}{}{ὑμῖν δ᾽ ἀκούειν. εὖ γὰρ ἴστε, παῖς ἐμὸς}
\verso{212}{\notagramatical{πράξας}{πράσσω}{}{part., aor., at., nom., sing., masc.}}{πράξας μὲν εὖ θαυμαστὸς ἂν γένοιτ᾽ ἀνήρ,}
\verso{213}{}{κακῶς δὲ πράξασ—οὐχ ὑπεύθυνος πόλει,}
\verso{214}{\notagramatical{σωθεὶς}{σώζω}{salvar}{part., aor., pass., nom., sing., masc.}\crlf \notagramatical{κοιρανεῖ}{κοιρανέω}{governar}{pres., ind., at., 3ª s.}}{σωθεὶς δ᾽ ὁμοίως τῆσδε κοιρανεῖ χθονός.}
\falante{ΧΟΡΟΣ}
\verso{215}{}{οὔ σε βουλόμεσθα, μῆτερ, οὔτ᾽ ἄγαν φοβεῖν λόγοις}
\verso{216}{\notagramatical{προστροπαῖς}{προστροπή, -ῆς, ἡ}{súplica}{dat., pl.}}{οὔτε θαρσύνειν. θεοὺς δὲ προστροπαῖς ἱκνουμένη,}
\verso{217}{\notagramatical{εἶδες}{ὁράω}{}{aor., ind., at., 2ª s.}}{εἴ τι φλαῦρον εἶδες, αἰτοῦ τῶνδ᾽ ἀποτροπὴν τελεῖν,}
\verso{218}{}{τὰ δ᾽ ἀγάθ᾽ ἐκτελῆ γενέσθαι σοί τε καὶ τέκνοις σέθεν}
\verso{219}{\notagramatical{χοὰς}{χοή, -ῆς, ἡ}{libação}{ac., pl.}}{καὶ πόλει φίλοις τε πᾶσι. δεύτερον δὲ χρὴ χοὰς}
\verso{220}{\notagramatical{φθιτοῖς}{φθιτός/\allowbreak -ή/\allowbreak -όν}{morto}{adj., dat., pl., masc.}\crlf \notagramatical{χέασθαι}{χέω}{verter}{inf., aor., méd.}\crlf \notagramatical{πρευμενῶς}{πρευμενής/\allowbreak πρευμενές}{benévolo}{adv.}}{γῇ τε καὶ φθιτοῖς χέασθαι: πρευμενῶς δ᾽ αἰτοῦ τάδε,}
\verso{221}{}{σὸν πόσιν Δαρεῖον, ὅνπερ φὴς ἰδεῖν κατ᾽ εὐφρόνην,}
\verso{222}{\notagramatical{πέμπειν}{πέμπω}{}{inf., pres., at.}\crlf \notagramatical{φάος}{φάος, -ους, τό}{luz}{ac., sing.}}{ἐσθλά σοι πέμπειν τέκνῳ τε γῆς ἔνερθεν ἐς φάος,}
\verso{223}{\desdobramento{τἄμπαλιν}{τὰ ἔμπαλιν}}{τἄμπαλιν δὲ τῶνδε γαίᾳ κάτοχ᾽ ἀμαυροῦσθαι σκότῳ.}
\verso{224}{\notagramatical{ὢν}{εἰμί}{}{part., pres., at., nom., sing., masc.}\crlf \notagramatical{παρῄνεσα}{παραινέω}{aconselhar}{aor., ind., at., 1ª s.}}{ταῦτα θυμόμαντις ὢν σοὶ πρευμενῶς παρῄνεσα:}
\verso{225}{}{εὖ δὲ πανταχῇ τελεῖν σοι τῶνδε κρίνομεν πέρι.}
\falante{ΑΤΟΣΣΑ}
\verso{226}{}{ἀλλὰ μὴν εὔνους γ᾽ ὁ πρῶτος τῶνδ᾽ ἐνυπνίων κριτὴς}
\verso{227}{\notagramatical{δόμοις}{δόμος, -ου, ὁ}{}{dat., pl.}\crlf \notagramatical{φάτιν}{φάτις, -ιος, ἡ}{declaração}{ac., sing.}}{παιδὶ καὶ δόμοις ἐμοῖσι τήνδ᾽ ἐκύρωσας φάτιν.}
\verso{228}{\notagramatical{ἐκτελοῖτο}{ἐκτελέω}{cumprir}{pres., opt., méd./pass., 3ª s.; opt. de desejo, sem ἄν: que se cumpra}}{ἐκτελοῖτο δὴ τὰ χρηστά: ταῦτα δ᾽, ὡς ἐφίεσαι,}
\verso{229}{}{πάντα θήσομεν θεοῖσι τοῖς τ᾽ ἔνερθε γῆς φίλοις,}
\verso{230}{\notagramatical{μόλωμεν}{βλώσκω}{}{aor., subj., at., 1ª p.; εὖτ᾽ ἄν + subj.: quando chegarmos; evento futuro}\crlf \notagramatical{θέλω}{ἐθέλω}{}{pres., ind., at., 1ª s.}}{εὖτ᾽ ἂν εἰς οἴκους μόλωμεν. κεῖνο δ᾽ ἐκμαθεῖν θέλω,}
\verso{231}{}{ὦ φίλοι, ποῦ τὰς Ἀθήνας φασὶν ἱδρῦσθαι χθονός.}
\falante{ΧΟΡΟΣ}
\verso{232}{\preposicaopreposta{πρὸς}{dat.}{junto a}}{τῆλε πρὸς δυσμαῖς ἄνακτος Ἡλίου φθινασμάτων.}
\falante{ΑΤΟΣΣΑ}
\verso{233}{\notagramatical{ἵμειρ᾽}{ἱμείρω}{desejar}{imperf., ind., at., 3ª s.}\crlf \notagramatical{θηρᾶσαι}{θηράω}{caçar, conquistar}{inf., aor., at.}}{ἀλλὰ μὴν ἵμειρ᾽ ἐμὸς παῖς τήνδε θηρᾶσαι πόλιν;}
\falante{ΧΟΡΟΣ}
\verso{234}{}{πᾶσα γὰρ γένοιτ᾽ ἂν Ἑλλὰς βασιλέως ὑπήκοος.}
\falante{ΑΤΟΣΣΑ}
\verso{235}{\notagramatical{πάρεστιν}{πάρειμι}{estar presente}{pres., ind., at., 3ª s.}}{ὧδέ τις πάρεστιν αὐτοῖς ἀνδροπλήθεια στρατοῦ;}
\falante{ΧΟΡΟΣ}
\verso{236}{\notagramatical{ἔρξας}{ἔρδω}{fazer}{part., aor., at., nom., sing., masc.}}{καὶ στρατὸς τοιοῦτος, ἔρξας πολλὰ δὴ Μήδους κακά.}
\falante{ΑΤΟΣΣΑ}
\verso{237}{}{καὶ τί πρὸς τούτοισιν ἄλλο; πλοῦτος ἐξαρκὴς δόμοις;}
\falante{ΧΟΡΟΣ}
\verso{238}{}{ἀργύρου πηγή τις αὐτοῖς ἐστι, θησαυρὸς χθονός.}
\falante{ΑΤΟΣΣΑ}
\verso{239}{\notagramatical{τοξουλκὸς}{τοξουλκός/\allowbreak -όν}{}{adj., nom., sing., fem.}\crlf \preposicaopreposta{διὰ}{gen.}{por meio de}}{πότερα γὰρ τοξουλκὸς αἰχμὴ διὰ χερῶν αὐτοῖς πρέπει;}
\falante{ΧΟΡΟΣ}
\verso{240}{}{οὐδαμῶς: ἔγχη σταδαῖα καὶ φεράσπιδες σαγαί.}
\falante{ΑΤΟΣΣΑ}
\verso{241}{\notagramatical{ποιμάνωρ}{ποιμάνωρ, ποιμάνορος, ὁ}{pastor, chefe}{nom., sing.}\crlf \desdobramento{κἀπιδεσπόζει}{καὶ ἐπιδεσπόζει}}{τίς δὲ ποιμάνωρ ἔπεστι κἀπιδεσπόζει στρατῷ;}
\falante{ΧΟΡΟΣ}
\verso{242}{}{οὔτινος δοῦλοι κέκληνται φωτὸς οὐδ᾽ ὑπήκοοι.}
\falante{ΑΤΟΣΣΑ}
\verso{243}{\notagramatical{ἐπήλυδας}{ἔπηλυς, ἐπήλυδος, ὁ}{invasor}{ac., pl.}}{πῶς ἂν οὖν μένοιεν ἄνδρας πολεμίους ἐπήλυδας;}
\falante{ΧΟΡΟΣ}
\verso{244}{\notagramatical{φθεῖραι}{φθείρω}{destruir}{inf., aor., at.}}{ὥστε Δαρείου πολύν τε καὶ καλὸν φθεῖραι στρατόν.}
\falante{ΑΤΟΣΣΑ}
\verso{245}{\notagramatical{δεινά}{δεινός/\allowbreak -ή/\allowbreak -όν}{}{adj., ac., pl., neut.}\crlf \notagramatical{τεκοῦσι}{τίκτω}{gerar}{part., aor., at., dat., pl., masc.}}{δεινά τοι λέγεις κιόντων τοῖς τεκοῦσι φροντίσαι.}
\falante{ΧΟΡΟΣ}
\verso{246}{}{ἀλλ᾽ ἐμοὶ δοκεῖν τάχ᾽ εἴσει πάντα ναμερτῆ λόγον.}
\verso{247}{}{τοῦδε γὰρ δράμημα φωτὸς Περσικὸν πρέπει μαθεῖν,}
\verso{248}{\notagramatical{ἐσθλὸν}{ἐσθλός/\allowbreak -ή/\allowbreak -όν}{}{adj., ac., sing., neut.}\crlf \notagramatical{κλύειν}{κλύω}{ouvir}{inf., pres., at.}}{καὶ φέρει σαφές τι πρᾶγος ἐσθλὸν ἢ κακὸν κλύειν.}
\falante{ΑΓΓΕΛΟΣ}
\verso{249}{}{ὦ γῆς ἁπάσης Ἀσιάδος πολίσματα,}
\verso{250}{}{ὦ Περσὶς αἶα καὶ μέγας πλούτου λιμήν,}
\verso{251}{\notagramatical{κατέφθαρται}{καταφθείρω}{destruir por completo}{perf., ind., pass., 3ª s.}}{ὡς ἐν μιᾷ πληγῇ κατέφθαρται πολὺς}
\verso{252}{\notagramatical{ὄλβος}{ὄλβος, -ου, ὁ}{prosperidade}{nom., sing.}\crlf \notagramatical{πεσόν}{πίπτω}{}{part., aor., at., ac., sing., neut.}}{ὄλβος, τὸ Περσῶν δ᾽ ἄνθος οἴχεται πεσόν.}
\verso{253}{}{ὤμοι, κακὸν μὲν πρῶτον ἀγγέλλειν κακά:}
\verso{254}{\notagramatical{πάθος}{πάθος, -ους, τό}{sofrimento}{ac., sing.}}{ὅμως δ᾽ ἀνάγκη πᾶν ἀναπτύξαι πάθος,}
\verso{255}{\notagramatical{ὄλωλε}{ὄλλυμι}{perecer, destruir}{perf., ind., at., 3ª s.}\crlf \notagramatical{βαρβάρων}{βάρβαρος/\allowbreak -ον}{}{adj., gen., pl., masc.}}{Πέρσαι: στρατὸς γὰρ πᾶς ὄλωλε βαρβάρων.}
\falante{ΧΟΡΟΣ}
\verso{256}{}{ἄνι᾽ ἄνια κακὰ}
\verso{257}{\notagramatical{νεόκοτα}{νεόκοτος/\allowbreak -ον}{novo, insólito}{adj., nom., pl., neut.}\crlf \notagramatical{δάι᾽}{δάιος/\allowbreak -α/\allowbreak -ον}{hostil, destruidor}{adj., nom., pl., neut.}}{νεόκοτα καὶ δάι᾽. αἰαῖ,}
\verso{258}{}{διαίνεσθε, Πέρσαι,}
\verso{259}{\notagramatical{ἄχος}{ἄχος, -ους, τό}{dor}{ac., sing.}\crlf \notagramatical{κλύοντες}{κλύω}{}{part., pres., at., nom., pl., masc.}}{τόδ᾽ ἄχος κλύοντες.}
\falante{ΑΓΓΕΛΟΣ}
\verso{260}{}{ὡς πάντα γ᾽ ἔστ᾽ ἐκεῖνα διαπεπραγμένα:}
\verso{261}{\notagramatical{ἀέλπτως}{ἄελπτος/\allowbreak ἄελπτον}{inesperado}{adv.}\crlf \notagramatical{νόστιμον}{νόστιμος/\allowbreak -ον}{do retorno}{adj., ac., sing., neut.}}{αὐτὸς δ᾽ ἀέλπτως νόστιμον βλέπω φάος.}
\falante{ΧΟΡΟΣ}
\verso{262}{}{ἦ μακροβίοτος}
\verso{263}{\notagramatical{ἐφάνθη}{φαίνω}{}{aor., ind., pass., 3ª s.}}{ὅδε γέ τις αἰὼν ἐφάνθη}
\verso{264}{}{γεραιοῖς, ἀκούειν}
\verso{265}{\notagramatical{πῆμ᾽}{πῆμα, πήματος, τό}{desgraça}{ac., sing.}}{τόδε πῆμ᾽ ἄελπτον.}
\falante{ΑΓΓΕΛΟΣ}
\verso{266}{\notagramatical{παρών}{πάρειμι}{}{part., pres., at., nom., sing., masc.}\crlf \desdobramento{κοὐ}{καὶ οὐ}}{καὶ μὴν παρών γε κοὐ λόγους ἄλλων κλύων,}
\verso{267}{\notagramatical{φράσαιμ᾽}{φράζω}{}{aor., opt., at., 1ª s.; opt. + ἄν: poderia relatar}\crlf \notagramatical{ἐπορσύνθη}{πορσύνω}{causar, preparar}{aor., ind., pass., 3ª s.}}{Πέρσαι, φράσαιμ᾽ ἂν οἷ᾽ ἐπορσύνθη κακά.}
\falante{ΧΟΡΟΣ}
\verso{268}{}{ὀτοτοτοῖ, μάταν}
\verso{269}{\notagramatical{παμμιγῆ}{παμμιγής/\allowbreak -ές}{de toda espécie}{adj., nom., pl., neut.}}{τὰ πολέα βέλεα παμμιγῆ}
\verso{270}{}{γᾶσδ᾽ ἀπ᾽ Ἀσίδος ἦλθεν αἴας}
\verso{271}{}{δᾴαν Ἑλλάδα χώραν.}
\falante{ΑΓΓΕΛΟΣ}
\verso{272}{\notagramatical{πλήθουσι}{πλήθω}{estar cheio}{pres., ind., at., 3ª p.}}{πλήθουσι νεκρῶν δυσπότμως ἐφθαρμένων}
\verso{273}{\notagramatical{ἀκταὶ}{ἀκτή, -ῆς, ἡ}{costa}{nom., pl.}}{Σαλαμῖνος ἀκταὶ πᾶς τε πρόσχωρος τόπος.}
\falante{ΧΟΡΟΣ}
\verso{274}{}{ὀτοτοτοῖ, φίλων}
\verso{275}{\notagramatical{πολυβαφῆ}{πολυβαφής/\allowbreak -ές}{muito mergulhado}{adj., ac., pl., neut.}}{ἁλίδονα μέλεα πολυβαφῆ}
\verso{276}{}{κατθανόντα λέγεις φέρεσθαι}
\verso{277}{}{†πλαγκτοῖς ἐν διπλάκεσσιν.}
\falante{ΑΓΓΕΛΟΣ}
\verso{278}{\notagramatical{ἀπώλλυτο}{ἀπόλλυμι}{destruir, perder}{imperf., ind., méd./pass., 3ª s.}}{οὐδὲν γὰρ ἤρκει τόξα, πᾶς δ᾽ ἀπώλλυτο}
\verso{279}{\notagramatical{ἐμβολαῖς}{ἐμβολή, -ῆς, ἡ}{ataque com esporão}{dat., pl.}}{στρατὸς δαμασθεὶς ναΐοισιν ἐμβολαῖς}
\falante{ΧΟΡΟΣ}
\verso{280}{\notagramatical{ἴυζ᾽}{ἰύζω}{soltar gritos}{pres., imper., at., 2ª s.}}{ἴυζ᾽ ἄποτμον Πέρσαις}
\verso{281}{}{δυσαιανῆ βοὰν}
\verso{282}{}{δᾴος, ὡς πάντα παγκάκως}
\verso{283}{}{ἤνυσαν, αἰαῖ, στρατοῦ φθαρέντος.}
\falante{ΑΓΓΕΛΟΣ}
\verso{284}{}{ὦ πλεῖστον ἔχθος ὄνομα Σαλαμῖνος κλύειν:}
\verso{285}{\notagramatical{στένω}{στένω}{gemer}{pres., ind., at., 1ª s.}\crlf \notagramatical{μεμνημένος}{μιμνήσκω}{lembrar}{part., perf., méd./pass., nom., sing., masc.}}{φεῦ, τῶν Ἀθηνῶν ὡς στένω μεμνημένος.}
\falante{ΧΟΡΟΣ}
\verso{286}{}{στυγναί γ᾽ Ἀθᾶναι δᾴοις:}
\verso{287}{\notagramatical{μεμνῆσθαί}{μιμνήσκω}{}{inf., perf., méd./pass.}}{μεμνῆσθαί τοι πάρα}
\verso{288}{}{ὡς πολλὰς Περσίδων μάταν}
\verso{289}{\notagramatical{ἀνάνδρους}{ἄνανδρος/\allowbreak -ον}{}{adj., ac., pl., fem.}}{εὔνιδας ἔκτισσαν ἠδ᾽ ἀνάνδρους.}
\falante{ΑΤΟΣΣΑ}
\verso{290}{\notagramatical{ἐκπεπληγμένη}{ἐκπλήγνυμι}{atordoar}{part., perf., pass., nom. sing. fem.}}{σιγῶ πάλαι δύστηνος ἐκπεπληγμένη}
\verso{291}{\notagramatical{συμφορά}{συμφορά, -ᾶς, ἡ}{desgraça}{nom., sing.}}{κακοῖς: ὑπερβάλλει γὰρ ἥδε συμφορά,}
\verso{292}{\notagramatical{λέξαι}{λέγω}{}{inf., aor., at.}\crlf \notagramatical{πάθη}{πάθος, -ους, τό}{}{ac., pl.}}{τὸ μήτε λέξαι μήτ᾽ ἐρωτῆσαι πάθη.}
\verso{293}{\notagramatical{βροτοῖς}{βροτός, -οῦ, ὁ}{}{dat., pl.}\crlf \notagramatical{φέρειν}{φέρω}{}{inf., pres., at.}}{ὅμως δ᾽ ἀνάγκη πημονὰς βροτοῖς φέρειν}
\verso{294}{\notagramatical{ἀναπτύξας}{ἀναπτύσσω}{expor em detalhe}{part., aor., at., nom., sing., masc.}}{θεῶν διδόντων: πᾶν δ᾽ ἀναπτύξας πάθος}
\verso{295}{\notagramatical{λέξον}{λέγω}{}{aor., imper., at., 2ª s.}\crlf \notagramatical{στένεις}{στένω}{}{pres., ind., at., 2ª s.}}{λέξον καταστάς, κεἰ στένεις κακοῖς ὅμως,}
\verso{296}{\notagramatical{τέθνηκε}{θνήσκω}{morrer}{perf., ind., at., 3ª s.}}{τίς οὐ τέθνηκε, τίνα δὲ καὶ πενθήσομεν}
\verso{297}{\notagramatical{σκηπτουχίᾳ}{σκηπτουχία, -ας, ἡ}{comando do cetro}{dat., sing.}}{τῶν ἀρχελείων, ὅστ᾽ ἐπὶ σκηπτουχίᾳ}
\verso{298}{\notagramatical{ἄνανδρον}{ἄνανδρος/\allowbreak -ον}{}{adj., ac., sing., fem.}\crlf \notagramatical{ἠρήμου}{ἐρημόω}{deixar deserto}{imperf., ind., at., 3ª s.}\crlf \notagramatical{θανών}{θνήσκω}{}{part., aor., at., nom., sing., masc.}}{ταχθεὶς ἄνανδρον τάξιν ἠρήμου θανών.}
\falante{ΑΓΓΕΛΟΣ}
\verso{299}{}{Ξέρξης μὲν αὐτὸς ζῇ τε καὶ βλέπει φάος.}
\falante{ΑΤΟΣΣΑ}
\verso{300}{\notagramatical{εἶπας}{λέγω}{}{aor., ind., at., 2ª s.}}{ἐμοῖς μὲν εἶπας δώμασιν φάος μέγα}
\verso{301}{\notagramatical{μελαγχίμου}{μελάγχιμος/\allowbreak -ον}{de veste negra}{adj., gen., sing., fem.}}{καὶ λευκὸν ἦμαρ νυκτὸς ἐκ μελαγχίμου.}
\falante{ΑΓΓΕΛΟΣ}
\verso{302}{\notagramatical{βραβεὺς}{βραβεύς, βραβέως, ὁ}{comandante, árbitro}{nom., sing.}}{Ἀρτεμβάρης δὲ μυρίας ἵππου βραβεὺς}
\verso{303}{\notagramatical{ἀκτὰς}{ἀκτή, -ῆς, ἡ}{}{ac., pl.}}{στύφλους παρ᾽ ἀκτὰς θείνεται Σιληνιῶν.}
\verso{304}{\desdobramento{χὠ}{καὶ ὁ}}{χὠ χιλίαρχος Δαδάκης πληγῇ δορὸς}
\verso{305}{\notagramatical{νεὼς}{ναῦς, νεώς, ἡ}{navio}{gen., sing.}}{πήδημα κοῦφον ἐκ νεὼς ἀφήλατο:}
\verso{306}{}{Τενάγων τ᾽ ἀριστεὺς Βακτρίων ἰθαγενὴς}
\verso{307}{}{θαλασσόπληκτον νῆσον Αἴαντος πολεῖ.}
\verso{308}{\desdobramento{κἀργήστης}{καὶ Ἀργήστης}}{Λίλαιος, Ἀρσάμης τε κἀργήστης τρίτος,}
\verso{309}{\notagramatical{πελειοθρέμμονα}{πελειοθρέμμων/\allowbreak -ον}{que cria pombas}{adj., ac., sing., fem.}}{οἵδ᾽ ἀμφὶ νῆσον τὴν πελειοθρέμμονα}
\verso{310}{\notagramatical{δινούμενοι}{δινεύω}{girar}{part., pres., méd./pass., nom., pl., masc.}\crlf \notagramatical{κύρισσον}{κυρίσσω}{golpear}{imperf., ind., at., 3ª p.}}{δινούμενοι κύρισσον ἰσχυρὰν χθόνα:}
\verso{311}{}{πηγαῖς τε Νείλου γειτονῶν Αἰγυπτίου}
\verso{312}{}{Ἀρκτεύς, Ἀδεύης, καὶ Φερεσσάκης τρίτος,}
\verso{313}{\notagramatical{πέσον}{πίπτω}{}{aor., ind., at., 3ª p.}}{Φαρνοῦχος, οἵδε ναὸς ἐκ μιᾶς πέσον.}
\verso{314}{}{Χρυσεὺς Μάταλλος μυριόνταρχος θανών,}
\verso{315}{}{ἵππου μελαίνης ἡγεμὼν τρισμυρίας,}
\verso{316}{}{πυρσὴν ζαπληθῆ δάσκιον γενειάδα}
\verso{317}{\notagramatical{ἔτεγγ᾽}{τέγγω}{molhar}{imperf., ind., at., 3ª s.}\crlf \notagramatical{ἀμείβων}{ἀμείβω}{}{part., pres., at., nom., sing., masc.}}{ἔτεγγ᾽, ἀμείβων χρῶτα πορφυρᾷ βαφῇ.}
\verso{318}{}{καὶ Μᾶγος Ἄραβος, Ἀρτάβης τε Βάκτριος,}
\verso{319}{\notagramaticalquebrada{κατέφθιτο}{καταφθίω}{perecer}{mais-que-perf., ind., méd./pass., 3ª s.}}{σκληρᾶς μέτοικος γῆς, ἐκεῖ κατέφθιτο.}
\verso{320}{}{Ἄμιστρις Ἀμφιστρεύς τε πολύπονον δόρυ}
\verso{321}{\notagramatical{νωμῶν}{νωμάω}{manejar}{part., pres., at., nom., sing., masc.}}{νωμῶν, ὅ τ᾽ ἐσθλὸς Ἀριόμαρδος Σάρδεσι}
\verso{322}{\notagramatical{παρασχών}{παρέχω}{}{part., aor., at., nom., sing., masc.}}{πένθος παρασχών, Σεισάμης θ᾽ ὁ Μύσιος,}
\verso{323}{\notagramatical{νεῶν}{ναῦς, νεώς, ἡ}{}{gen., pl.}}{Θάρυβίς τε πεντήκοντα πεντάκις νεῶν}
\verso{324}{}{ταγός, γένος Λυρναῖος, εὐειδὴς ἀνήρ,}
\verso{325}{\notagramatical{θανὼν}{θνήσκω}{}{part., aor., at., nom., sing., masc.}}{κεῖται θανὼν δείλαιος οὐ μάλ᾽ εὐτυχῶς:}
\verso{326}{}{Συέννεσίς τε πρῶτος εἰς εὐψυχίαν,}
\verso{327}{}{Κιλίκων ἔπαρχος, εἷς ἀνὴρ πλεῖστον πόνον}
\verso{328}{}{ἐχθροῖς παρασχών, εὐκλεῶς ἀπώλετο.}
\verso{329}{\notagramatical{ταγῶν}{ταγός, -οῦ, ὁ}{}{gen., pl.}\crlf \notagramatical{ὑπεμνήσθην}{ὑπομιμνήσκω}{lembrar}{aor., ind., pass., 1ª s.}}{τοσόνδε ταγῶν νῦν ὑπεμνήσθην πέρι.}
\verso{330}{}{πολλῶν παρόντων δ᾽ ὀλίγ᾽ ἀπαγγέλλω κακά.}
\falante{ΑΤΟΣΣΑ}
\verso{331}{\notagramatical{κλύω}{κλύω}{}{pres., ind., at., 1ª s.}}{αἰαῖ, κακῶν ὕψιστα δὴ κλύω τάδε,}
\verso{332}{\notagramatical{κωκύματα}{κώκυμα, κωκύματος, τό}{lamento}{ac., pl.}}{αἴσχη τε Πέρσαις καὶ λιγέα κωκύματα.}
\verso{333}{\notagramatical{φράσον}{φράζω}{}{aor., imper., at., 2ª s.}}{ἀτὰρ φράσον μοι τοῦτ᾽ ἀναστρέψας πάλιν:}
\verso{334}{}{πόσον δὲ πλῆθος ἦν νεῶν Ἑλληνίδων,}
\verso{335}{}{ὥστ᾽ ἀξιῶσαι Περσικῷ στρατεύματι}
\verso{336}{\notagramatical{συνάψαι}{συνάπτω}{travar, unir}{inf., aor., at.}}{μάχην συνάψαι ναΐοισιν ἐμβολαῖς;}
\falante{ΑΓΓΕΛΟΣ}
\verso{337}{\notagramatical{ἴσθ᾽}{οἶδα}{}{perf., ind., at., 2ª p.}\crlf \notagramatical{βάρβαρον}{βάρβαρος/\allowbreak -ον}{}{adj., ac., sing., masc.}}{πλήθους μὲν ἂν σάφ᾽ ἴσθ᾽ ἕκατι βάρβαρον}
\verso{338}{}{ναυσὶν κρατῆσαι. καὶ γὰρ Ἕλλησιν μὲν ἦν}
\verso{339}{}{ὁ πᾶς ἀριθμὸς ἐς τριακάδας δέκα}
\verso{340}{}{ναῶν, δεκὰς δ᾽ ἦν τῶνδε χωρὶς ἔκκριτος:}
\verso{341}{\notagramatical{οἶδα}{οἶδα}{}{perf., ind., at., 1ª s.}}{Ξέρξῃ δέ, καὶ γὰρ οἶδα, χιλιὰς μὲν ἦν}
\verso{342}{\notagramatical{ὑπέρκοποι}{ὑπέρκοπος/\allowbreak -ον}{excepcional}{adj., nom., pl., fem.}}{ὧν ἦγε πλῆθος, αἱ δ᾽ ὑπέρκοποι τάχει}
\verso{343}{\notagramatical{ἦσαν}{εἰμί}{}{imperf., ind., at., 3ª p.}}{ἑκατὸν δὶς ἦσαν ἑπτά θ᾽: ὧδ᾽ ἔχει λόγος.}
\verso{344}{\notagramatical{λειφθῆναι}{λείπω}{}{inf., aor., pass.}}{μή σοι δοκοῦμεν τῇδε λειφθῆναι μάχῃ;}
\verso{345}{\notagramatical{κατέφθειρε}{καταφθείρω}{}{aor., ind., at., 3ª s.}}{ἀλλ᾽ ὧδε δαίμων τις κατέφθειρε στρατόν,}
\verso{346}{\notagramatical{βρίσας}{βρίθω}{pesar}{part., aor., at., nom., sing., masc.}}{τάλαντα βρίσας οὐκ ἰσορρόπῳ τύχῃ.}
\verso{347}{\notagramatical{σῴζουσι}{σώζω}{}{pres., ind., at., 3ª p.}}{θεοὶ πόλιν σῴζουσι Παλλάδος θεᾶς.}
\falante{ΑΤΟΣΣΑ}
\verso{348}{\notagramatical{ἀπόρθητος}{ἀπόρθητος/\allowbreak -ον}{não conquistado}{adj., nom., sing., fem.}}{ἔτ᾽ ἆρ᾽ Ἀθηνῶν ἔστ᾽ ἀπόρθητος πόλις;}
\falante{ΑΓΓΕΛΟΣ}
\verso{349}{\notagramatical{ὄντων}{εἰμί}{}{part., pres., at., gen., pl., masc.}}{ἀνδρῶν γὰρ ὄντων ἕρκος ἐστὶν ἀσφαλές.}
\falante{ΑΤΟΣΣΑ}
\verso{350}{\notagramatical{ναυσὶ}{ναῦς, νεώς, ἡ}{}{dat., pl.}\crlf \notagramatical{συμβολῆς}{συμβολή, -ῆς, ἡ}{combate, encontro}{gen., sing.}}{ἀρχὴ δὲ ναυσὶ συμβολῆς τίς ἦν, φράσον:}
\verso{351}{}{τίνες κατῆρξαν, πότερον Ἕλληνες, μάχης,}
\verso{352}{\notagramatical{καταυχήσας}{καταυχέω}{vangloriar-se}{part., aor., at., nom., sing., masc.}}{ἢ παῖς ἐμός, πλήθει καταυχήσας νεῶν;}
\falante{ΑΓΓΕΛΟΣ}
\verso{353}{}{ἦρξεν μέν, ὦ δέσποινα, τοῦ παντὸς κακοῦ}
\verso{354}{\notagramatical{φανεὶς}{φαίνω}{}{part., aor., pass., nom., sing., masc.}\crlf \notagramatical{ἀλάστωρ}{ἀλάστωρ, ἀλάστορος, ὁ}{gênio vingador}{nom., sing.}}{φανεὶς ἀλάστωρ ἢ κακὸς δαίμων ποθέν.}
\verso{355}{}{ἀνὴρ γὰρ Ἕλλην ἐξ Ἀθηναίων στρατοῦ}
\verso{356}{\notagramatical{ἐλθὼν}{ἔρχομαι}{}{part., aor., at., nom., sing., masc.}\crlf \notagramatical{ἔλεξε}{λέγω}{}{aor., ind., at., 3ª s.}}{ἐλθὼν ἔλεξε παιδὶ σῷ Ξέρξῃ τάδε,}
\verso{357}{}{ὡς εἰ μελαίνης νυκτὸς ἵξεται κνέφας,}
\verso{358}{\notagramatical{σέλμασιν}{σέλμα, σέλματος, τό}{banco de remador}{dat., pl.}}{Ἕλληνες οὐ μενοῖεν, ἀλλὰ σέλμασιν}
\verso{359}{}{ναῶν ἐπενθορόντες ἄλλος ἄλλοσε}
\verso{360}{\notagramatical{ἐκσωσοίατο}{ἐκσῴζω}{salvar}{fut., opt., méd., 3ª p.; opt. em discurso indireto após verbo no passado: salvariam a vida}}{δρασμῷ κρυφαίῳ βίοτον ἐκσωσοίατο.}
\verso{361}{}{ὁ δ᾽ εὐθὺς ὡς ἤκουσεν, οὐ ξυνεὶς δόλον}
\verso{362}{}{Ἕλληνος ἀνδρὸς οὐδὲ τὸν θεῶν φθόνον,}
\verso{363}{\notagramatical{προφωνεῖ}{προφωνέω}{proclamar}{pres., ind., at., 3ª s.}}{πᾶσιν προφωνεῖ τόνδε ναυάρχοις λόγον,}
\verso{364}{\notagramatical{φλέγων}{φλέγω}{arder}{part., pres., at., nom., sing., masc.}}{εὖτ᾽ ἂν φλέγων ἀκτῖσιν ἥλιος χθόνα}
\verso{365}{\notagramatical{λήξῃ}{λήγω}{cessar}{aor., subj., at., 3ª s.; εὖτ᾽ ἄν + subj.: quando cessar; evento futuro}\crlf \notagramatical{λάβῃ}{λαμβάνω}{tomar, receber}{aor., subj., at., 3ª s.; εὖτ᾽ ἄν + subj.: quando a escuridão tomar; evento futuro}}{λήξῃ, κνέφας δὲ τέμενος αἰθέρος λάβῃ,}
\verso{366}{}{τάξαι νεῶν μὲν στῖφος ἐν στοίχοις τρισὶν}
\verso{367}{}{ἔκπλους φυλάσσειν καὶ πόρους ἁλιρρόθους,}
\verso{368}{}{ἄλλας δὲ κύκλῳ νῆσον Αἴαντος πέριξ:}
\verso{369}{\notagramatical{μόρον}{μόρος, -ου, ὁ}{destino, morte}{ac., sing.}}{ὡς εἰ μόρον φευξοίαθ᾽ Ἕλληνες κακόν,}
\verso{370}{}{ναυσὶν κρυφαίως δρασμὸν εὑρόντες τινά,}
\verso{371}{\notagramatical{στέρεσθαι}{στέρομαι}{ser privado de}{inf., pres., méd.}\crlf \notagramatical{προκείμενον}{πρόκειμαι}{estar exposto}{part., pres., méd./pass., nom., sing., neut.}}{πᾶσιν στέρεσθαι κρατὸς ἦν προκείμενον.}
\verso{372}{}{τοσαῦτ᾽ ἔλεξε κάρθ᾽ ὑπ᾽ εὐθύμου φρενός:}
\verso{373}{}{οὐ γὰρ τὸ μέλλον ἐκ θεῶν ἠπίστατο.}
\verso{374}{\notagramatical{ἀκόσμως}{ἄκοσμος/\allowbreak ἄκοσμον}{desordenado}{adv.}\crlf \notagramatical{πειθάρχῳ}{πείθαρχος/\allowbreak -ον}{obediente}{adj., dat., sing., fem.}}{οἱ δ᾽ οὐκ ἀκόσμως, ἀλλὰ πειθάρχῳ φρενὶ}
\verso{375}{}{δεῖπνόν <τ᾽> ἐπορσύνοντο, ναυβάτης τ᾽ ἀνὴρ}
\verso{376}{\notagramatical{τροποῦτο}{τροπόω}{prender o remo}{imperf., ind., méd./pass., 3ª s.}\crlf \notagramatical{σκαλμὸν}{σκαλμός, -οῦ, ὁ}{tolete}{ac., sing.}\crlf \notagramatical{εὐήρετμον}{εὐήρετμος/\allowbreak -ον}{de bons remos}{adj., ac., sing., masc.}}{τροποῦτο κώπην σκαλμὸν ἀμφ᾽ εὐήρετμον.}
\verso{377}{}{ἐπεὶ δὲ φέγγος ἡλίου κατέφθιτο}
\verso{378}{}{καὶ νὺξ ἐπῄει, πᾶς ἀνὴρ κώπης ἄναξ}
\verso{379}{}{ἐς ναῦν ἐχώρει πᾶς θ᾽ ὅπλων ἐπιστάτης:}
\verso{380}{\notagramatical{παρεκάλει}{παρακαλέω}{exortar}{imperf., ind., at., 3ª s.}}{τάξις δὲ τάξιν παρεκάλει νεὼς μακρᾶς:}
\verso{381}{}{πλέουσι δ᾽ ὡς ἕκαστος ἦν τεταγμένος,}
\verso{382}{\notagramatical{διάπλοον}{διάπλοος, -ου, ὁ}{travessia marítima}{ac. sing.}}{καὶ πάννυχοι δὴ διάπλοον καθίστασαν}
\verso{383}{\notagramatical{ἄνακτες}{ἄναξ, ἄνακτος, ὁ}{soberano}{nom., pl.}}{ναῶν ἄνακτες πάντα ναυτικὸν λεών.}
\verso{384}{\desdobramento{κοὐ}{καὶ οὐ}}{καὶ νὺξ ἐχώρει, κοὐ μάλ᾽ Ἑλλήνων στρατὸς}
\verso{385}{}{κρυφαῖον ἔκπλουν οὐδαμῇ καθίστατο:}
\verso{386}{\notagramatical{λευκόπωλος}{λευκόπωλος/\allowbreak -ον}{de cavalos brancos}{adj., nom., sing., fem.}}{ἐπεί γε μέντοι λευκόπωλος ἡμέρα}
\verso{387}{\notagramatical{εὐφεγγὴς}{εὐφεγγής/\allowbreak -ές}{luminoso}{adj., nom., sing., fem.}}{πᾶσαν κατέσχε γαῖαν εὐφεγγὴς ἰδεῖν,}
\verso{388}{\notagramatical{κέλαδος}{κέλαδος, -ου, ὁ}{clamor}{nom., sing.}}{πρῶτον μὲν ἠχῇ κέλαδος Ἑλλήνων πάρα}
\verso{389}{\notagramatical{ηὐφήμησεν}{εὐφημέω}{entoar canto auspicioso}{aor., ind., at., 3ª s.}}{μολπηδὸν ηὐφήμησεν, ὄρθιον δ᾽ ἅμα}
\verso{390}{\notagramatical{ἀντηλάλαξε}{ἀνταλαλάζω}{responder com gritos}{aor., ind., at., 3ª s.}}{ἀντηλάλαξε νησιώτιδος πέτρας}
\verso{391}{\notagramatical{βαρβάροις}{βάρβαρος/\allowbreak -ον}{}{adj., dat., pl., masc.}\crlf \notagramatical{παρῆν}{πάρειμι}{}{imperf., ind., at., 3ª s.}}{ἠχώ: φόβος δὲ πᾶσι βαρβάροις παρῆν}
\verso{392}{}{γνώμης ἀποσφαλεῖσιν: οὐ γὰρ ὡς φυγῇ}
\verso{393}{\notagramatical{ἐφύμνουν}{ἐφυμνέω}{entoar}{imperf., ind., at., 3ª p.}}{παιᾶν᾽ ἐφύμνουν σεμνὸν Ἕλληνες τότε,}
\verso{394}{}{ἀλλ᾽ ἐς μάχην ὁρμῶντες εὐψύχῳ θράσει:}
\verso{395}{\notagramatical{ἐπέφλεγεν}{ἐπιφλέγω}{incendiar, inflamar}{imperf., ind., at., 3ª s.}}{σάλπιγξ δ᾽ ἀυτῇ πάντ᾽ ἐκεῖν᾽ ἐπέφλεγεν.}
\verso{396}{\notagramatical{ῥοθιάδος}{ῥοθιάς, ῥοθιάδος, ἡ}{batida de remos}{gen., sing.}}{εὐθὺς δὲ κώπης ῥοθιάδος ξυνεμβολῇ}
\verso{397}{\notagramatical{βρύχιον}{βρύχιος/\allowbreak -ον}{profundo}{adj., ac., sing., fem.}\crlf \notagramatical{κελεύματος}{κέλευμα, κελεύματος, τό}{ordem, clamor}{gen., sing.}}{ἔπαισαν ἅλμην βρύχιον ἐκ κελεύματος,}
\verso{398}{}{θοῶς δὲ πάντες ἦσαν ἐκφανεῖς ἰδεῖν.}
\verso{399}{}{τὸ δεξιὸν μὲν πρῶτον εὐτάκτως κέρας}
\verso{400}{}{ἡγεῖτο κόσμῳ, δεύτερον δ᾽ ὁ πᾶς στόλος}
\verso{401}{}{ἐπεξεχώρει, καὶ παρῆν ὁμοῦ κλύειν}
\verso{402}{\notagramatical{ἴτε}{εἶμι}{ir}{pres., imper., at., 2ª p.}}{πολλὴν βοήν, “ὦ παῖδες Ἑλλήνων ἴτε,}
\verso{403}{}{ἐλευθεροῦτε πατρίδ᾽, ἐλευθεροῦτε δὲ}
\verso{404}{}{παῖδας, γυναῖκας, θεῶν τε πατρῴων ἕδη,}
\verso{405}{}{θήκας τε προγόνων: νῦν ὑπὲρ πάντων ἀγών.”}
\verso{406}{}{καὶ μὴν παρ᾽ ἡμῶν Περσίδος γλώσσης ῥόθος}
\verso{407}{\desdobramento{κοὐκέτ᾽}{καὶ οὐκέτι}}{ὑπηντίαζε, κοὐκέτ᾽ ἦν μέλλειν ἀκμή.}
\verso{408}{}{εὐθὺς δὲ ναῦς ἐν νηὶ χαλκήρη στόλον}
\verso{409}{\notagramatical{ἦρξε}{ἄρχω}{}{aor., ind., at., 3ª s.}}{ἔπαισεν: ἦρξε δ᾽ ἐμβολῆς Ἑλληνικὴ}
\verso{410}{\desdobramento{κἀποθραύει}{καὶ ἀποθραύει}}{ναῦς, κἀποθραύει πάντα Φοινίσσης νεὼς}
\verso{411}{}{κόρυμβ᾽, ἐπ᾽ ἄλλην δ᾽ ἄλλος ηὔθυνεν δόρυ.}
\verso{412}{}{τὰ πρῶτα μέν νυν ῥεῦμα Περσικοῦ στρατοῦ}
\verso{413}{}{ἀντεῖχεν: ὡς δὲ πλῆθος ἐν στενῷ νεῶν}
\verso{414}{}{ἤθροιστ᾽, ἀρωγὴ δ᾽ οὔτις ἀλλήλοις παρῆν,}
\verso{415}{}{αὐτοὶ θ᾽ ὑφ᾽ αὑτῶν ἐμβόλοις χαλκοστόμοις}
\verso{416}{\notagramatical{ἔθραυον}{θραύω}{quebrar}{imperf., ind., at., 3ª p.}}{παίοντ᾽, ἔθραυον πάντα κωπήρη στόλον,}
\verso{417}{\notagramatical{νῆες}{ναῦς, νεώς, ἡ}{}{nom., pl.}}{Ἑλληνικαί τε νῆες οὐκ ἀφρασμόνως}
\verso{418}{\notagramatical{ἔθεινον}{θείνω}{golpear}{imperf., ind., at., 3ª p.}}{κύκλῳ πέριξ ἔθεινον, ὑπτιοῦτο δὲ}
\verso{419}{}{σκάφη νεῶν, θάλασσα δ᾽ οὐκέτ᾽ ἦν ἰδεῖν,}
\verso{420}{\notagramatical{πλήθουσα}{πλήθω}{}{part., pres., at., nom., sing., fem.}\crlf \notagramatical{βροτῶν}{βροτός, -οῦ, ὁ}{}{gen., pl.}}{ναυαγίων πλήθουσα καὶ φόνου βροτῶν.}
\verso{421}{}{ἀκταὶ δὲ νεκρῶν χοιράδες τ᾽ ἐπλήθυον,}
\verso{422}{\notagramatical{ἀκόσμῳ}{ἄκοσμος/\allowbreak -ον}{}{adj., dat., sing., fem.}\crlf \notagramatical{ἠρέσσετο}{ἐρέσσω}{remar; golpear}{imperf., ind., méd./pass., 3ª s.}}{φυγῇ δ᾽ ἀκόσμῳ πᾶσα ναῦς ἠρέσσετο,}
\verso{423}{}{ὅσαιπερ ἦσαν βαρβάρου στρατεύματος.}
\verso{424}{}{τοὶ δ᾽ ὥστε θύννους ἤ τιν᾽ ἰχθύων βόλον}
\verso{425}{}{ἀγαῖσι κωπῶν θραύμασίν τ᾽ ἐρειπίων}
\verso{426}{\notagramatical{ἐρράχιζον}{ῥαχίζω}{quebrar a espinha}{imperf., ind., at., 3ª p.}\crlf \notagramatical{οἰμωγὴ}{οἰμωγή, -ῆς, ἡ}{grito de dor}{nom., sing.}}{ἔπαιον, ἐρράχιζον: οἰμωγὴ δ᾽ ὁμοῦ}
\verso{427}{}{κωκύμασιν κατεῖχε πελαγίαν ἅλα,}
\verso{428}{\notagramatical{κελαινῆς}{κελαινός/\allowbreak -ή/\allowbreak -όν}{negro}{adj., gen., sing., fem.}\crlf \notagramatical{ὄμμ᾽}{ὄμμα, ὄμματος, τό}{olho, olhar}{nom., sing.}}{ἕως κελαινῆς νυκτὸς ὄμμ᾽ ἀφείλετο.}
\verso{429}{}{κακῶν δὲ πλῆθος, οὐδ᾽ ἂν εἰ δέκ᾽ ἤματα}
\verso{430}{\notagramatical{στοιχηγοροίην}{στοιχηγορέω}{enumerar}{pres., opt., at., 1ª s.; εἰ + opt.: mesmo se eu enumerasse; hipótese}\crlf \notagramatical{ἐκπλήσαιμί}{ἐκπίμπλημι}{encher por completo}{aor., opt., at., 1ª s.; οὐκ ἄν + opt.: não conseguiria completar}}{στοιχηγοροίην, οὐκ ἂν ἐκπλήσαιμί σοι.}
\verso{431}{}{εὖ γὰρ τόδ᾽ ἴσθι, μηδάμ᾽ ἡμέρᾳ μιᾷ}
\verso{432}{}{πλῆθος τοσουτάριθμον ἀνθρώπων θανεῖν.}
\falante{ΑΤΟΣΣΑ}
\verso{433}{\notagramatical{ἔρρωγεν}{ῥήγνυμι}{}{perf., ind., at., 3ª s.}}{αἰαῖ, κακῶν δὴ πέλαγος ἔρρωγεν μέγα}
\verso{434}{}{Πέρσαις τε καὶ πρόπαντι βαρβάρων γένει.}
\falante{ΑΓΓΕΛΟΣ}
\verso{435}{\notagramatical{μεσοῦν}{μεσόω}{chegar à metade}{inf., pres., at.}}{εὖ νυν τόδ᾽ ἴσθι, μηδέπω μεσοῦν κακόν:}
\verso{436}{\notagramatical{ἦλθε}{ἔρχομαι}{}{aor., ind., at., 3ª s.}}{τοιάδ᾽ ἐπ᾽ αὐτοὺς ἦλθε συμφορὰ πάθους,}
\verso{437}{}{ὡς τοῖσδε καὶ δὶς ἀντισηκῶσαι ῥοπῇ.}
\falante{ΑΤΟΣΣΑ}
\verso{438}{}{καὶ τίς γένοιτ᾽ ἂν τῆσδ᾽ ἔτ᾽ ἐχθίων τύχη;}
\verso{439}{}{λέξον τίν᾽ αὖ φῂς τήνδε συμφορὰν στρατῷ}
\verso{440}{\notagramatical{ἐλθεῖν}{ἔρχομαι}{}{inf., aor., at.}}{ἐλθεῖν κακῶν ῥέπουσαν ἐς τὰ μάσσονα.}
\falante{ΑΓΓΕΛΟΣ}
\verso{441}{}{Περσῶν ὅσοιπερ ἦσαν ἀκμαῖοι φύσιν,}
\verso{442}{}{ψυχήν τ᾽ ἄριστοι κεὐγένειαν ἐκπρεπεῖς,}
\verso{443}{}{αὐτῷ τ᾽ ἄνακτι πίστιν ἐν πρώτοις ἀεί,}
\verso{444}{\notagramatical{τεθνᾶσιν}{θνήσκω}{}{perf., ind., at., 3ª p.}\crlf \notagramatical{δυσκλεεστάτῳ}{δυσκλεής/\allowbreak -ές}{de má fama}{adj., superl., dat., sing., masc.}}{τεθνᾶσιν αἰσχρῶς δυσκλεεστάτῳ μόρῳ.}
\falante{ΑΤΟΣΣΑ}
\verso{445}{}{οἲ ᾽γὼ τάλαινα συμφορᾶς κακῆς, φίλοι.}
\verso{446}{\notagramatical{ὀλωλέναι}{ὄλλυμι}{}{inf., perf., at.}}{ποίῳ μόρῳ δὲ τούσδε φῂς ὀλωλέναι;}
\falante{ΑΓΓΕΛΟΣ}
\verso{447}{}{νῆσός τις ἐστὶ πρόσθε Σαλαμῖνος τόπων,}
\verso{448}{\notagramatical{ναυσίν}{ναῦς, νεώς, ἡ}{}{dat., pl.}}{βαιά, δύσορμος ναυσίν, ἣν ὁ φιλόχορος}
\verso{449}{\notagramatical{ἐμβατεύει}{ἐμβατεύω}{pisar, frequentar}{pres., ind., at., 3ª s.}\crlf \notagramatical{ποντίας}{πόντιος/\allowbreak -α/\allowbreak -ον}{marinho}{adj., gen., sing., fem.}}{Πὰν ἐμβατεύει, ποντίας ἀκτῆς ἔπι.}
\verso{450}{}{ἐνταῦθα πέμπει τούσδ᾽, ὅπως, ὅταν νεῶν}
\verso{451}{\notagramatical{φθαρέντες}{φθείρω}{}{part., aor., pass., nom., pl., masc.}\crlf \notagramatical{ἐκσῳζοίατο}{ἐκσῴζω}{}{pres., opt., méd./pass., 3ª p.; opt. na subordinada eventual em contexto passado: quando escapassem}}{φθαρέντες ἐχθροὶ νῆσον ἐκσῳζοίατο,}
\verso{452}{}{κτείνοιεν εὐχείρωτον Ἑλλήνων στρατόν,}
\verso{453}{\notagramatical{πόρων}{πόρος, -ου, ὁ}{passagem}{gen., pl.}}{φίλους δ᾽ ὑπεκσῴζοιεν ἐναλίων πόρων,}
\verso{454}{\notagramatical{ἱστορῶν}{ἱστορέω}{informar-se}{part., pres., at., nom., sing., masc.}}{κακῶς τὸ μέλλον ἱστορῶν. ὡς γὰρ θεὸς}
\verso{455}{}{ναῶν ἔδωκε κῦδος Ἕλλησιν μάχης,}
\verso{456}{\notagramatical{φράξαντες}{φράσσω}{cercar, fechar}{part., aor., at., nom., pl., masc.}}{αὐθημερὸν φράξαντες εὐχάλκοις δέμας}
\verso{457}{\notagramatical{ἐξέθρῳσκον}{ἐκθρῴσκω}{saltar para fora}{imperf., ind., at., 3ª p.}}{ὅπλοισι ναῶν ἐξέθρῳσκον: ἀμφὶ δὲ}
\verso{458}{\notagramatical{κυκλοῦντο}{κυκλόω}{cercar}{imperf., ind., méd./pass., 3ª p.}\crlf \notagramatical{ἀμηχανεῖν}{ἀμηχανέω}{estar sem saída}{inf., pres., at.}}{κυκλοῦντο πᾶσαν νῆσον, ὥστ᾽ ἀμηχανεῖν}
\verso{459}{}{ὅποι τράποιντο. πολλὰ μὲν γὰρ ἐκ χερῶν}
\verso{460}{\notagramatical{ἠράσσοντο}{ἀράσσω}{golpear}{imperf., ind., méd./pass., 3ª p.}}{πέτροισιν ἠράσσοντο, τοξικῆς τ᾽ ἀπὸ}
\verso{461}{\notagramatical{προσπίτνοντες}{προσπίπτω}{}{part., pres., at., nom., pl., masc.}\crlf \notagramatical{ὤλλυσαν}{ὄλλυμι}{}{imperf., ind., at., 3ª p.}}{θώμιγγος ἰοὶ προσπίτνοντες ὤλλυσαν:}
\verso{462}{}{τέλος δ᾽ ἐφορμηθέντες ἐξ ἑνὸς ῥόθου}
\verso{463}{\notagramatical{κρεοκοποῦσι}{κρεοκοπέω}{cortar em pedaços}{pres., ind., at., 3ª p.}}{παίουσι, κρεοκοποῦσι δυστήνων μέλη,}
\verso{464}{\notagramatical{ἐξαπέφθειραν}{ἐξαποφθείρω}{aniquilar}{aor., ind., at., 3ª p.}}{ἕως ἁπάντων ἐξαπέφθειραν βίον.}
\verso{465}{\notagramatical{ὁρῶν}{ὁράω}{}{part., pres., at., nom., sing., masc.}}{Ξέρξης δ᾽ ἀνῴμωξεν κακῶν ὁρῶν βάθος:}
\verso{466}{\notagramatical{εἶχε}{ἔχω}{}{imperf., ind., at., 3ª s.}\crlf \notagramatical{εὐαγῆ}{εὐαγής/\allowbreak -ές}{aberto, visível}{adj., ac., sing., fem.}}{ἕδραν γὰρ εἶχε παντὸς εὐαγῆ στρατοῦ,}
\verso{467}{}{ὑψηλὸν ὄχθον ἄγχι πελαγίας ἁλός:}
\verso{468}{\desdobramento{κἀνακωκύσας}{καὶ ἀνακωκύσας}}{ῥήξας δὲ πέπλους κἀνακωκύσας λιγύ,}
\verso{469}{\notagramatical{πεζῷ}{πεζός/\allowbreak -ή/\allowbreak -όν}{a pé}{adj., dat., sing., neut.}\crlf \notagramatical{παραγγείλας}{παραγγέλλω}{dar ordens}{part., aor., at., nom., sing., masc.}}{πεζῷ παραγγείλας ἄφαρ στρατεύματι,}
\verso{470}{}{ἵησ᾽ ἀκόσμῳ ξὺν φυγῇ. τοιάνδε σοι}
\verso{471}{\notagramatical{στένειν}{στένω}{}{inf., pres., at.}}{πρὸς τῇ πάροιθε συμφορὰν στένειν πάρα.}
\falante{ΑΤΟΣΣΑ}
\verso{472}{\notagramatical{δαῖμον}{δαίμων, δαίμονος, ὁ}{}{voc., sing.}}{ὦ στυγνὲ δαῖμον, ὡς ἄρ᾽ ἔψευσας φρενῶν}
\verso{473}{}{Πέρσας: πικρὰν δὲ παῖς ἐμὸς τιμωρίαν}
\verso{474}{\desdobramento{κοὐκ}{καὶ οὐκ}\crlf \notagramatical{ἀπήρκεσαν}{ἀπαρκέω}{bastar}{aor., ind., at., 3ª p.}}{κλεινῶν Ἀθηνῶν ηὗρε, κοὐκ ἀπήρκεσαν}
\verso{475}{}{οὓς πρόσθε Μαραθὼν βαρβάρων ἀπώλεσεν:}
\verso{476}{\notagramatical{πράξειν}{πράσσω}{}{inf., fut., at.}}{ὧν ἀντίποινα παῖς ἐμὸς πράξειν δοκῶν}
\verso{477}{\notagramatical{πημάτων}{πῆμα, πήματος, τό}{}{gen., pl.}\crlf \notagramatical{ἐπέσπασεν}{ἐπισπάω}{atrair para si}{aor., ind., at., 3ª s.}}{τοσόνδε πλῆθος πημάτων ἐπέσπασεν.}
\verso{478}{}{σὺ δ᾽ εἰπέ, ναῶν αἳ πεφεύγασιν μόρον,}
\verso{479}{\notagramatical{ἔλειπες}{λείπω}{}{imperf., ind., at., 2ª s.}\crlf \notagramatical{οἶσθα}{οἶδα}{}{perf., ind., at., 2ª s.}\crlf \notagramatical{σημῆναι}{σημαίνω}{dar sinal}{inf., aor., at.}}{ποῦ τάσδ᾽ ἔλειπες: οἶσθα σημῆναι τορῶς;}
\falante{ΑΓΓΕΛΟΣ}
\verso{480}{}{ναῶν γε ταγοὶ τῶν λελειμμένων σύδην}
\verso{481}{\notagramatical{οὖρον}{οὖρος, -ου, ὁ}{vento favorável}{ac., sing.}}{κατ᾽ οὖρον οὐκ εὔκοσμον αἴρονται φυγήν:}
\verso{482}{}{στρατὸς δ᾽ ὁ λοιπὸς ἔν τε Βοιωτῶν χθονὶ}
\verso{483}{\notagramatical{διώλλυθ᾽}{διόλλυμι}{destruir por completo}{imperf., ind., méd./pass., 3ª s.}\crlf \notagramatical{κρηναῖον}{κρηναῖος/\allowbreak -α/\allowbreak -ον}{de fonte}{adj., ac., sing., neut.}\crlf \notagramatical{γάνος}{γάνος, -ους, τό}{frescor, brilho}{ac., sing.}}{διώλλυθ᾽, οἱ μὲν ἀμφὶ κρηναῖον γάνος}
\verso{484}{}{δίψῃ πονοῦντες, οἱ δ᾽ ὑπ᾽ ἄσθματος κενοὶ}
\verso{485}{}{διεκπερῶμεν ἔς τε Φωκέων χθόνα}
\verso{486}{}{καὶ Δωρίδ᾽ αἶαν, Μηλιᾶ τε κόλπον, οὗ}
\verso{487}{\notagramatical{ἄρδει}{ἄρδω}{irrigar}{pres., ind., at., 3ª s.}}{Σπερχειὸς ἄρδει πεδίον εὐμενεῖ ποτῷ:}
\verso{488}{\desdobramento{κἀντεῦθεν}{καὶ ἐντεῦθεν}}{κἀντεῦθεν ἡμᾶς γῆς Ἀχαιίδος πέδον}
\verso{489}{}{καὶ Θεσσαλῶν πόλεις ὑπεσπανισμένους}
\verso{490}{\notagramatical{θάνον}{θνήσκω}{}{aor., ind., at., 3ª p.}}{βορᾶς ἐδέξαντ᾽: ἔνθα δὴ πλεῖστοι θάνον}
\verso{491}{}{δίψῃ τε λιμῷ τ᾽: ἀμφότερα γὰρ ἦν τάδε.}
\verso{492}{}{Μαγνητικὴν δὲ γαῖαν ἔς τε Μακεδόνων}
\verso{493}{}{χώραν ἀφικόμεσθ᾽, ἐπ᾽ Ἀξίου πόρον,}
\verso{494}{\notagramatical{ἕλειον}{ἕλειος/\allowbreak -α/\allowbreak -ον}{pantanoso}{adj., ac., sing., masc.}\crlf \notagramatical{δόνακα}{δόναξ, δόνακος, ὁ}{junco}{ac., sing.}}{Βόλβης θ᾽ ἕλειον δόνακα, Πάγγαιόν τ᾽ ὄρος,}
\verso{495}{}{Ἠδωνίδ᾽ αἶαν: νυκτὶ δ᾽ ἐν ταύτῃ θεὸς}
\verso{496}{\notagramatical{ὦρσε}{ὄρνυμι}{suscitar}{aor., ind., at., 3ª s.}\crlf \notagramatical{πήγνυσιν}{πήγνυμι}{congelar, fixar}{pres., ind., at., 3ª s.}}{χειμῶν᾽ ἄωρον ὦρσε, πήγνυσιν δὲ πᾶν}
\verso{497}{}{ῥέεθρον ἁγνοῦ Στρυμόνος. θεοὺς δέ τις}
\verso{498}{}{τὸ πρὶν νομίζων οὐδαμοῦ τότ᾽ ηὔχετο}
\verso{499}{\notagramatical{λιταῖσι}{λιτή, -ῆς, ἡ}{prece}{dat., pl.}}{λιταῖσι, γαῖαν οὐρανόν τε προσκυνῶν.}
\verso{500}{\notagramatical{θεοκλυτῶν}{θεοκλυτέω}{invocar os deuses}{part., pres., at., nom., sing., masc.}}{ἐπεὶ δὲ πολλὰ θεοκλυτῶν ἐπαύσατο}
\verso{501}{\notagramatical{περᾷ}{περάω}{}{pres., ind., at., 3ª s.}}{στρατός, περᾷ κρυσταλλοπῆγα διὰ πόρον:}
\verso{502}{\notagramatical{σκεδασθῆναι}{σκεδάννυμι}{dispersar}{inf., aor., pass.}}{χὤστις μὲν ἡμῶν πρὶν σκεδασθῆναι θεοῦ}
\verso{503}{\notagramatical{σεσωσμένος}{σώζω}{}{part., perf., méd./pass., nom., sing., masc.}\crlf \notagramatical{κυρεῖ}{κυρέω}{encontrar, alcançar}{pres., ind., at., 3ª s.}}{ἀκτῖνας ὡρμήθη, σεσωσμένος κυρεῖ.}
\verso{504}{}{φλέγων γὰρ αὐγαῖς λαμπρὸς ἡλίου κύκλος}
\verso{505}{\notagramatical{διῆκε}{διήκω}{atravessar}{imperf., ind., at., 3ª s.}}{μέσον πόρον διῆκε, θερμαίνων φλογί:}
\verso{506}{\notagramatical{πῖπτον}{πίπτω}{}{imperf., ind., at., 3ª p.}}{πῖπτον δ᾽ ἐπ᾽ ἀλλήλοισιν: ηὐτύχει δέ τοι}
\verso{507}{}{ὅστις τάχιστα πνεῦμ᾽ ἀπέρρηξεν βίου.}
\verso{508}{\desdobramento{κἄτυχον}{καὶ ἔτυχον}}{ὅσοι δὲ λοιποὶ κἄτυχον σωτηρίας,}
\verso{509}{\notagramatical{περάσαντες}{περάω}{}{part., aor., at., nom., pl., masc.}}{Θρῄκην περάσαντες μόγις πολλῷ πόνῳ,}
\verso{510}{}{ἥκουσιν ἐκφυγόντες, οὐ πολλοί τινες,}
\verso{511}{\notagramatical{ἑστιοῦχον}{ἑστιοῦχος/\allowbreak -ον}{que possui um lar}{adj., ac., sing., fem.}}{ἐφ᾽ ἑστιοῦχον γαῖαν: ὡς στένειν πόλιν}
\verso{512}{}{Περσῶν, ποθοῦσαν φιλτάτην ἥβην χθονός.}
\verso{513}{\notagramatical{ἐκλείπω}{ἐκλείπω}{}{pres., ind., at., 1ª s.}\crlf \notagramatical{λέγων}{λέγω}{}{part., pres., at., nom., sing., masc.}}{ταῦτ᾽ ἔστ᾽ ἀληθῆ: πολλὰ δ᾽ ἐκλείπω λέγων}
\verso{514}{}{κακῶν ἃ Πέρσαις ἐγκατέσκηψεν θεός.}
\falante{ΧΟΡΟΣ}
\verso{515}{}{ὦ δυσπόνητε δαῖμον, ὡς ἄγαν βαρὺς}
\verso{516}{}{ποδοῖν ἐνήλλου παντὶ Περσικῷ γένει.}
\falante{ΑΤΟΣΣΑ}
\verso{517}{\notagramatical{διαπεπραγμένου}{διαπράσσω}{consumar}{part., perf., pass., gen. sing. masc.}}{οἲ ᾽γὼ τάλαινα διαπεπραγμένου στρατοῦ:}
\verso{518}{}{ὦ νυκτὸς ὄψις ἐμφανὴς ἐνυπνίων,}
\verso{519}{}{ὡς κάρτα μοι σαφῶς ἐδήλωσας κακά.}
\verso{520}{}{ὑμεῖς δὲ φαύλως αὔτ᾽ ἄγαν ἐκρίνατε.}
\verso{521}{\notagramatical{ἐκύρωσεν}{κυρόω}{confirmar}{aor., ind., at., 3ª s.}}{ὅμως δ᾽, ἐπειδὴ τῇδ᾽ ἐκύρωσεν φάτις}
\verso{522}{}{ὑμῶν, θεοῖς μὲν πρῶτον εὔξασθαι θέλω:}
\verso{523}{}{ἔπειτα γῇ τε καὶ φθιτοῖς δωρήματα}
\verso{524}{\notagramatical{λαβοῦσα}{λαμβάνω}{}{part., aor., at., nom., sing., fem.}}{ἥξω λαβοῦσα πέλανον ἐξ οἴκων ἐμῶν,—}
\verso{525}{}{ἐπίσταμαι μὲν ὡς ἐπ᾽ ἐξειργασμένοις,}
\verso{526}{\notagramatical{λῷον}{λωίων/\allowbreak -ον}{melhor}{adj., compar., nom., sing., neut.}\crlf \notagramatical{πέλοι}{πέλω}{ser, tornar-se}{pres., opt., at., 3ª s.; εἴ + opt.: na esperança de que algo seja melhor}}{ἀλλ᾽ ἐς τὸ λοιπὸν εἴ τι δὴ λῷον πέλοι.}
\verso{527}{}{ὑμᾶς δὲ χρὴ ᾽πὶ τοῖσδε τοῖς πεπραγμένοις}
\verso{528}{\notagramatical{πιστὰ}{πιστός/\allowbreak -ή/\allowbreak -όν}{}{adj., ac., pl., neut.}}{πιστοῖσι πιστὰ ξυμφέρειν βουλεύματα:}
\verso{529}{\notagramatical{μόλῃ}{βλώσκω}{}{aor., subj., at., 3ª s.; ἐάν + subj.: se vier; condição futura}}{καὶ παῖδ᾽, ἐάν περ δεῦρ᾽ ἐμοῦ πρόσθεν μόλῃ,}
\verso{530}{\notagramatical{παρηγορεῖτε}{παρηγορέω}{consolar}{pres., imper., at., 2ª p.}\crlf \notagramatical{προπέμπετ᾽}{προπέμπω}{acompanhar, enviar}{pres., imper., at., 2ª p.}}{παρηγορεῖτε, καὶ προπέμπετ᾽ ἐς δόμους,}
\verso{531}{}{μὴ καί τι πρὸς κακοῖσι προσθῆται κακόν.}
\parte{Primeiro Estásimo}{532–597}
\falante{ΧΟΡΟΣ}
\verso{532}{}{ὦ Ζεῦ βασιλεῦ, νῦν <μὲν> Περσῶν}
\verso{533}{}{τῶν μεγαλαύχων καὶ πολυάνδρων}
\verso{534}{\notagramatical{ὀλέσας}{ὄλλυμι}{}{part., aor., at., nom., sing., masc.}}{στρατιὰν ὀλέσας}
\verso{535}{}{ἄστυ τὸ Σούσων ἠδ᾽ Ἀγβατάνων}
\verso{536}{\notagramatical{δνοφερῷ}{δνοφερός/\allowbreak -ά/\allowbreak -όν}{sombrio}{adj., dat., sing., neut.}\crlf \notagramatical{κατέκρυψας}{κατακρύπτω}{ocultar}{aor., ind., at., 2ª s.}}{πένθει δνοφερῷ κατέκρυψας:}
\verso{537}{}{πολλαὶ δ᾽ ἀταλαῖς χερσὶ καλύπτρας}
\verso{538}{\notagramatical{κατερεικόμεναι}{κατερείκω}{rasgar em pedaços}{part., pres., méd./pass., nom., pl., fem.}}{κατερεικόμεναι}
\verso{539}{\notagramatical{διαμυδαλέοις}{διαμυδαλέος/\allowbreak -α/\allowbreak -ον}{encharcado}{adj., dat., pl., neut.}}{διαμυδαλέοις δάκρυσι κόλπους}
\verso{540}{\notagramatical{τέγγουσ᾽}{τέγγω}{}{pres., ind., at., 3ª p.}}{τέγγουσ᾽, ἄλγους μετέχουσαι.}
\verso{541}{}{αἱ δ᾽ ἁβρόγοοι Περσίδες ἀνδρῶν}
\verso{542}{}{ποθέουσαι ἰδεῖν ἀρτιζυγίαν,}
\verso{543}{}{λέκτρων εὐνὰς ἁβροχίτωνας,}
\verso{544}{\notagramatical{χλιδανῆς}{χλιδανός/\allowbreak -ή/\allowbreak -όν}{luxuoso, delicado}{adj., gen., sing., fem.}}{χλιδανῆς ἥβης τέρψιν, ἀφεῖσαι,}
\verso{545}{\notagramatical{γόοις}{γόος, -ου, ὁ}{lamento}{dat., pl.}}{πενθοῦσι γόοις ἀκορεστοτάτοις.}
\verso{546}{\desdobramento{κἀγὼ}{καὶ ἐγώ}\crlf \notagramatical{οἰχομένων}{οἴχομαι}{}{part., pres., méd., gen., pl., masc.}}{κἀγὼ δὲ μόρον τῶν οἰχομένων}
\verso{547}{}{αἴρω δοκίμως πολυπενθῆ.}
\pausaestrofe
\verso{548}{\notagramatical{στένει}{στένω}{}{pres., ind., at., 3ª s.}}{νῦν δὴ πρόπασα μὲν στένει}
\verso{549}{}{γαῖ᾽ Ἀσὶς ἐκκεκενωμένα.}
\verso{550}{}{Ξέρξης μὲν ἄγαγεν, ποποῖ,}
\verso{551}{}{Ξέρξης δ᾽ ἀπώλεσεν, τοτοῖ,}
\verso{552}{\notagramatical{ἐπέσπε}{ἐφέπω}{conduzir}{aor., ind., at., 3ª s.}}{Ξέρξης δὲ πάντ᾽ ἐπέσπε δυσφρόνως}
\verso{553}{\notagramatical{βαρίδεσσι}{βᾶρις, βάριδος, ἡ}{embarcação}{dat., pl.}\crlf \notagramatical{ποντίαις}{πόντιος/\allowbreak -α/\allowbreak -ον}{}{adj., dat., pl., fem.}}{βαρίδεσσι ποντίαις.}
\verso{554}{}{τίπτε Δαρεῖος μὲν οὕ-}
\verso{555}{}{τω τότ᾽ ἀβλαβὴς ἐπῆν}
\verso{556}{}{τόξαρχος πολιήταις,}
\verso{557}{}{Σουσίδαις φίλος ἄκτωρ;}
\pausaestrofe
\verso{558}{\notagramatical{πεζοὺς}{πεζός/\allowbreak -ή/\allowbreak -όν}{}{adj., ac., pl., masc.}}{πεζοὺς δὲ καὶ θαλασσίους}
\verso{559}{\notagramatical{λινόπτεροι}{λινόπτερος/\allowbreak -ον}{com asas de linho}{adj., nom., pl., fem.}}{λινόπτεροι κυανώπιδες}
\verso{560}{}{νᾶες μὲν ἄγαγον, ποποῖ,}
\verso{561}{\notagramatical{ἀπώλεσαν}{ἀπόλλυμι}{}{aor., ind., at., 3ª p.}}{νᾶες δ᾽ ἀπώλεσαν, τοτοῖ,}
\verso{562}{\notagramatical{πανωλέθροισιν}{πανώλεθρος/\allowbreak -ον}{inteiramente funesto}{adj., dat., pl., fem.}}{νᾶες πανωλέθροισιν ἐμβολαῖς,}
\verso{563}{}{ἠδ᾽ Ἰαόνων χέρες.}
\verso{564}{}{τυτθὰ δ᾽ ἐκφυγεῖν ἄνακτ᾽}
\verso{565}{\notagramatical{εἰσακούομεν}{εἰσακούω}{escutar}{pres., ind., at., 1ª p.}}{αὐτὸν εἰσακούομεν}
\verso{566}{\notagramatical{πεδιήρεις}{πεδιήρης/\allowbreak -ες}{que ocupa a planície}{adj., ac., pl., fem.}}{Θρῄκης ἂμ πεδιήρεις}
\verso{567}{}{δυσχίμους τε κελεύθους.}
\pausaestrofe
\verso{568}{\notagramatical{πρωτόμοιροι}{πρωτόμοιρος/\allowbreak -ον}{primeiro a morrer}{adj., nom., pl., masc.}}{τοὶ δ᾽ ἄρα πρωτόμοιροι, φεῦ,}
\verso{569}{}{ληφθέντες πρὸς ἀνάγκας, ἠέ,}
\verso{570}{}{ἀκτὰς ἀμφὶ Κυχρείας, ὀᾶ,}
\verso{571}{\notagramatical{στένε}{στένω}{}{pres., imper., at., 2ª s.}}{<ἔρρουσι:> στένε καὶ δακνά-}
\verso{572}{\notagramatical{ζου}{δακνάζομαι}{ser dilacerado}{pres., imper., méd., 2ª s.}\crlf \notagramatical{ἀμβόασον}{ἀναβοάω}{clamar}{aor., imper., at., 2ª s.}}{ζου, βαρὺ δ᾽ ἀμβόασον}
\verso{573}{}{οὐράνι᾽ ἄχη, ὀᾶ:}
\verso{574}{}{τεῖνε δὲ δυσβάυκτον}
\verso{575}{}{βοᾶτιν τάλαιναν αὐδάν.}
\pausaestrofe
\verso{576}{\notagramatical{γναπτόμενοι}{κνάπτω}{dilacerar}{part., pres., méd./pass., nom., pl., masc.}}{γναπτόμενοι δὲ δίνᾳ, φεῦ,}
\verso{577}{\notagramatical{σκύλλονται}{σκύλλω}{maltratar}{pres., ind., méd./pass., 3ª p.}}{σκύλλονται πρὸς ἀναύδων, ἠέ,}
\verso{578}{}{παίδων τᾶς ἀμιάντου, ὀᾶ.}
\verso{579}{}{πενθεῖ δ᾽ ἄνδρα δόμος στερη-}
\verso{580}{\notagramatical{θείς}{στερέω}{privar}{part., aor., pass., nom., sing., masc.}}{θείς, τοκέες τ᾽ ἄπαιδες}
\verso{581}{}{δαιμόνι᾽ ἄχη, ὀᾶ,}
\verso{582}{\notagramatical{δυρόμενοι}{ὀδύρομαι}{lamentar}{part., pres., méd., nom., pl., masc.}}{δυρόμενοι γέροντες}
\verso{583}{\notagramatical{κλύουσιν}{κλύω}{}{pres., ind., at., 3ª p.}}{τὸ πᾶν δὴ κλύουσιν ἄλγος.}
\pausaestrofe
\verso{584}{}{τοὶ δ᾽ ἀνὰ γᾶν Ἀσίαν δὴν}
\verso{585}{\notagramatical{περσονομοῦνται}{περσονομέομαι}{ser governado pelos persas}{pres., ind., méd., 3ª p.}}{οὐκέτι περσονομοῦνται,}
\verso{586}{\notagramatical{δασμοφοροῦσιν}{δασμοφορέω}{pagar tributo}{pres., ind., at., 3ª p.}}{οὐδ᾽ ἔτι δασμοφοροῦσιν}
\verso{587}{\notagramatical{δεσποσύνοισιν}{δεσπόσυνος/\allowbreak -η/\allowbreak -ον}{do senhor}{adj., dat., pl., fem.}}{δεσποσύνοισιν ἀνάγκαις,}
\verso{588}{\notagramatical{προπίτνοντες}{προπίπτω}{cair diante de}{part., pres., at., nom., pl., masc.}}{οὐδ᾽ ἐς γᾶν προπίτνοντες}
\verso{589}{\notagramatical{ἅζονται}{ἅζομαι}{reverenciar}{pres., ind., méd., 3ª p.}}{ἅζονται: βασιλεία}
\verso{590}{\notagramatical{διόλωλεν}{διόλλυμι}{}{perf., ind., at., 3ª s.}}{γὰρ διόλωλεν ἰσχύς.}
\pausaestrofe
\verso{591}{\notagramatical{βροτοῖσιν}{βροτός, -οῦ, ὁ}{}{dat., pl.}}{οὐδ᾽ ἔτι γλῶσσα βροτοῖσιν}
\verso{592}{}{ἐν φυλακαῖς: λέλυται γὰρ}
\verso{593}{\notagramatical{βάζειν}{βάζω}{falar}{inf., pres., at.}}{λαὸς ἐλεύθερα βάζειν,}
\verso{594}{}{ὡς ἐλύθη ζυγὸν ἀλκᾶς.}
\verso{595}{\notagramatical{αἱμαχθεῖσα}{αἱμάσσω}{ensanguentar}{part., aor., pass., nom., sing., fem.}}{αἱμαχθεῖσα δ᾽ ἄρουραν}
\verso{596}{}{Αἴαντος περικλύστα}
\verso{597}{}{νᾶσος ἔχει τὰ Περσᾶν.}
\parte{Segundo Episódio}{598–622}
\falante{ΑΤΟΣΣΑ}
\verso{598}{}{φίλοι, κακῶν μὲν ὅστις ἔμπειρος κυρεῖ,}
\verso{599}{\notagramatical{κλύδων}{κλύδων, -ος, ὁ}{vagalhão}{nom., sing.}}{ἐπίσταται βροτοῖσιν ὡς ὅταν κλύδων}
\verso{600}{\notagramatical{δειμαίνειν}{δειμαίνω}{temer}{inf., pres., at.}}{κακῶν ἐπέλθῃ πάντα δειμαίνειν φιλεῖ:}
\verso{601}{\notagramatical{εὐροῇ}{εὐροέω}{fluir bem}{pres., subj., at., 3ª s.; ὅταν + subj.: sempre que a sorte corre bem}}{ὅταν δ᾽ ὁ δαίμων εὐροῇ, πεποιθέναι}
\verso{602}{\notagramatical{οὐριεῖν}{οὐρίζω}{soprar favoravelmente}{inf., fut., at.}}{τὸν αὐτὸν αἰεὶ δαίμον᾽ οὐριεῖν τύχην.}
\verso{603}{}{ἐμοὶ γὰρ ἤδη πάντα μὲν φόβου πλέα}
\verso{604}{\desdobramento{τἀνταῖα}{τὰ ἀνταῖα}\crlf \notagramatical{φαίνεται}{φαίνω}{}{pres., ind., méd./pass., 3ª s.}}{ἐν ὄμμασιν τἀνταῖα φαίνεται θεῶν,}
\verso{605}{\notagramatical{βοᾷ}{βοάω}{gritar}{pres., ind., at., 3ª s.}\crlf \notagramatical{παιώνιος}{παιώνιος/\allowbreak -α/\allowbreak -ον}{fúnebre; aqui, de lamento}{adj., nom., sing., masc.}}{βοᾷ δ᾽ ἐν ὠσὶ κέλαδος οὐ παιώνιος:}
\verso{606}{}{τοία κακῶν ἔκπληξις ἐκφοβεῖ φρένας.}
\verso{607}{}{τοιγὰρ κέλευθον τήνδ᾽ ἄνευ τ᾽ ὀχημάτων}
\verso{608}{}{χλιδῆς τε τῆς πάροιθεν ἐκ δόμων πάλιν}
\verso{609}{\notagramatical{ἔστειλα}{στέλλω}{enviar, equipar}{aor., ind., at., 1ª s.}}{ἔστειλα, παιδὸς πατρὶ πρευμενεῖς χοὰς}
\verso{610}{\notagramatical{φέρουσ᾽}{φέρω}{}{part., pres., at., nom., sing., fem.}\crlf \notagramatical{μειλικτήρια}{μειλικτήριος/\allowbreak -ον}{propiciatório}{adj., ac., pl., neut.}}{φέρουσ᾽, ἅπερ νεκροῖσι μειλικτήρια,}
\verso{611}{}{βοός τ᾽ ἀφ᾽ ἁγνῆς λευκὸν εὔποτον γάλα,}
\verso{612}{\notagramatical{ἀνθεμουργοῦ}{ἀνθεμουργός/\allowbreak -όν}{que trabalha nas flores}{adj., gen., sing., fem.}}{τῆς τ᾽ ἀνθεμουργοῦ στάγμα, παμφαὲς μέλι,}
\verso{613}{}{λιβάσιν ὑδρηλαῖς παρθένου πηγῆς μέτα,}
\verso{614}{\notagramatical{ἀκήρατόν}{ἀκήρατος/\allowbreak -ον}{puro, intacto}{adj., ac., sing., neut.}\crlf \notagramatical{μητρὸς}{μήτηρ, μητρός, ἡ}{mãe}{gen., sing.}}{ἀκήρατόν τε μητρὸς ἀγρίας ἄπο}
\verso{615}{}{ποτὸν παλαιᾶς ἀμπέλου γάνος τόδε:}
\verso{616}{}{τῆς τ᾽ αἰὲν ἐν φύλλοισι θαλλούσης βίον}
\verso{617}{\notagramatical{εὐώδης}{εὐώδης/\allowbreak -ες}{perfumado}{adj., nom. sing. masc.}}{ξανθῆς ἐλαίας καρπὸς εὐώδης πάρα,}
\verso{618}{\notagramatical{παμφόρου}{πάμφορος/\allowbreak -ον}{que tudo produz}{adj., gen., sing., fem.}\crlf \notagramatical{γαίας}{γαῖα, -ας, ἡ}{}{ac., pl.}}{ἄνθη τε πλεκτά, παμφόρου γαίας τέκνα.}
\verso{619}{}{ἀλλ᾽, ὦ φίλοι, χοαῖσι ταῖσδε νερτέρων}
\verso{620}{\notagramatical{ἐπευφημεῖτε}{ἐπευφημέω}{entoar invocação}{pres., imper., at., 2ª p.}}{ὕμνους ἐπευφημεῖτε, τόν τε δαίμονα}
\verso{621}{}{Δαρεῖον ἀγκαλεῖσθε, γαπότους δ᾽ ἐγὼ}
\verso{622}{\notagramatical{προπέμψω}{προπέμπω}{}{fut., ind., at., 1ª s.}}{τιμὰς προπέμψω τάσδε νερτέροις θεοῖς.}
\parte{Segundo Estásimo / Canto de Evocação}{623–680}
\falante{ΧΟΡΟΣ}
\verso{623}{}{βασίλεια γύναι, πρέσβος Πέρσαις,}
\verso{624}{\notagramatical{πέμπε}{πέμπω}{}{pres., imper., at., 2ª s.}}{σύ τε πέμπε χοὰς θαλάμους ὑπὸ γῆς,}
\verso{625}{}{ἡμεῖς θ᾽ ὕμνοις αἰτησόμεθα}
\verso{626}{\notagramatical{πομποὺς}{πομπός, -οῦ, ὁ}{guia}{ac., pl.}}{φθιμένων πομποὺς}
\verso{627}{\notagramatical{εἶναι}{εἰμί}{}{inf., pres., at.}}{εὔφρονας εἶναι κατὰ γαίας.}
\verso{628}{\notagramatical{δαίμονες}{δαίμων, δαίμονος, ὁ}{}{voc., pl.}}{ἀλλά, χθόνιοι δαίμονες ἁγνοί,}
\verso{629}{}{Γῆ τε καὶ Ἑρμῆ, βασιλεῦ τ᾽ ἐνέρων,}
\verso{630}{}{πέμψατ᾽ ἔνερθεν ψυχὴν ἐς φῶς:}
\verso{631}{\notagramatical{ἄκος}{ἄκος, -εος, τό}{remédio}{ac., sing.}\crlf \notagramatical{οἶδε}{οἶδα}{}{perf., ind., at., 3ª s.}}{εἰ γάρ τι κακῶν ἄκος οἶδε πλέον,}
\verso{632}{\notagramatical{εἴποι}{λέγω}{dizer}{aor., opt., at., 3ª s.; ἄν + opt.: poderia dizer}}{μόνος ἂν θνητῶν πέρας εἴποι.}
\pausaestrofe
\verso{633}{\notagramatical{ἀίει}{ἀίω}{ouvir}{pres., ind., at., 3ª s.}\crlf \notagramaticalquebrada{μακαρίτας}{μακαρίτης, -ου, ὁ}{bem-aventurado, falecido}{nom., sing.}}{ἦ ῥ᾽ ἀίει μου μακαρίτας}
\verso{634}{}{ἰσοδαίμων βασιλεὺς}
\verso{635}{\notagramatical{βάρβαρ᾽}{βάρβαρος/\allowbreak -ον}{}{adj., ac., pl., neut.}}{βάρβαρ᾽ σαφηνῆ}
\verso{636}{\notagramatical{παναίολ᾽}{παναίολος/\allowbreak -ον}{multiforme}{adj., ac., pl., neut.}}{ἱέντος τὰ παναίολ᾽ αἰ-}
\verso{637}{\notagramatical{ανῆ}{αἰανής/\allowbreak -ές}{doloroso}{adj., ac., pl., neut.}\crlf \notagramatical{δύσθροα}{δύσθρους/\allowbreak -ουν}{de som lamentoso}{adj., ac., pl., neut.}\crlf \notagramatical{βάγματ᾽}{βάγμα, βάγματος, τό}{clamor}{ac., pl.}}{ανῆ δύσθροα βάγματ᾽, ἢ}
\verso{638}{}{παντάλαν᾽ ἄχη διαμβοάσω;}
\verso{639}{\notagramatical{κλύει}{κλύω}{}{pres., ind., at., 3ª s.}}{νέρθεν ἆρα κλύει μου;}
\pausaestrofe
\verso{640}{}{ἀλλὰ σύ μοι Γᾶ τε καὶ ἄλλοι}
\verso{641}{}{χθονίων ἁγεμόνες}
\verso{642}{\notagramatical{μεγαυχῆ}{μεγαυχής/\allowbreak -ές}{glorioso, altivo}{adj., ac., sing., masc.}}{δαίμονα μεγαυχῆ}
\verso{643}{\notagramatical{ἰόντ᾽}{εἶμι}{}{part., pres., at., ac., sing., masc.}}{ἰόντ᾽ αἰνέσατ᾽ ἐκ δόμων,}
\verso{644}{}{Περσᾶν Σουσιγενῆ θεόν:}
\verso{645}{\notagramatical{πέμπετε}{πέμπω}{}{pres., imper., at., 2ª p.}}{πέμπετε δ᾽ ἄνω <τὸν> οἷον οὔπω}
\verso{646}{}{Περσὶς αἶ᾽ ἐκάλυψεν.}
\pausaestrofe
\verso{647}{}{ἦ φίλος ἁνήρ, [ἦ] φίλος ὄχθος:}
\verso{648}{}{φίλα γὰρ κέκευθεν ἤθη.}
\verso{649}{}{Ἀιδωνεὺς δ᾽ ἀναπομπ-}
\verso{650}{\notagramatical{ἀνείης}{ἀνίημι}{enviar para cima}{aor., opt., at., 2ª s.; opt. de desejo, sem ἄν: que deixes subir}}{ὸς ἀνείης, Ἀιδωνεύς,}
\verso{651}{}{θεῖον ἀνάκτορα Δαριᾶνα. ἠέ.}
\pausaestrofe
\verso{652}{\notagramatical{ἀπώλλυ}{ἀπόλλυμι}{}{imperf., ind., at., 3ª s.}}{οὐδὲ γὰρ ἄνδρας ποτ᾽ ἀπώλλυ}
\verso{653}{}{πολεμοφθόροισιν ἄταις,}
\verso{654}{}{θεομήστωρ δ᾽ ἐκικλή-}
\verso{655}{\notagramatical{σκετο}{κικλήσκω}{chamar}{imperf., ind., méd./pass., 3ª s.}}{σκετο Πέρσαις, θεομήστωρ δ᾽}
\verso{656}{}{ἔσκεν, ἐπεὶ στρατὸν εὖ ποδούχει. ἠέ.}
\pausaestrofe
\verso{657}{}{βαλήν, ἀρχαῖος βαλήν,}
\verso{658}{\notagramatical{ἴθι}{εἶμι}{}{pres., imper., at., 2ª s.}}{ἴθι, ἱκοῦ:}
\verso{659}{\notagramatical{ἔλθ᾽}{ἔρχομαι}{}{aor., imper., at., 2ª s.}}{ἔλθ᾽ ἐπ᾽ ἄκρον κόρυμβον ὄχθου,}
\verso{660}{\notagramatical{κροκόβαπτον}{κροκόβαπτος/\allowbreak -ον}{tingido de açafrão}{adj., ac., sing., fem.}\crlf \notagramatical{εὔμαριν}{εὔμαρις, εὐμάριδος, ἡ}{calçado macio}{ac., sing.}\crlf \notagramatical{ἀείρων}{ἀείρω}{erguer}{part., pres., at., nom., sing., masc.}}{κροκόβαπτον ποδὸς εὔμαριν ἀείρων,}
\verso{661}{}{βασιλείου τιάρας}
\verso{662}{}{φάλαρον πιφαύσκων.}
\verso{663}{\notagramatical{βάσκε}{βάσκω}{vir, ir}{pres., imper., at., 2ª s.}}{βάσκε πάτερ ἄκακε Δαριάν, οἶ.}
\pausaestrofe
\verso{664}{\notagramatical{αἰανῆ}{αἰανής/\allowbreak -ές}{}{adj., ac., pl., neut.}\crlf \notagramatical{κλύῃς}{κλύω}{}{pres., subj., at., 2ª s.; ὅπως + subj.: para que ouças; finalidade}}{ὅπως αἰανῆ κλύῃς}
\verso{665}{}{νέα τ᾽ ἄχη,}
\verso{666}{}{δέσποτα δεσποτᾶν φάνηθι.}
\verso{667}{\notagramatical{ἀχλὺς}{ἀχλύς, -ύος, ἡ}{névoa, trevas}{nom., sing.}\crlf \notagramatical{πεπόταται}{ποτάομαι}{voar, pairar}{perf., ind., méd., 3ª s.}}{Στυγία γάρ τις ἐπ᾽ ἀχλὺς πεπόταται:}
\verso{668}{}{νεολαία γὰρ ἤδη}
\verso{670}{\notagramatical{ὄλωλεν}{ὄλλυμι}{}{perf., ind., at., 3ª s.}}{κατὰ πᾶσ᾽ ὄλωλεν.}
\verso{671}{}{βάσκε πάτερ ἄκακε Δαριάν, οἶ.}
\pausaestrofe
\verso{672}{}{αἰαῖ αἰαῖ:}
\verso{673}{\notagramatical{πολύκλαυτε}{πολύκλαυστος/\allowbreak -ον}{muito chorado}{adj., voc., sing., masc.}}{ὦ πολύκλαυτε φίλοισι θανών,}
\verso{675}{}{†τί τάδε δυνάτα δυνάτα}
\verso{676}{\notagramatical{ἁμάρτια}{ἁμάρτιον, -ου, τό}{erro}{nom., pl.}}{περὶ τᾷ σᾷ δίδυμα διαγόεν ἁμάρτια;†}
\verso{677}{}{πᾶσαι <γὰρ> γᾷ τᾷδ᾽ ἐξ-}
\verso{678}{\notagramatical{τρίσκαλμοι}{τρίσκαλμος/\allowbreak -ον}{de três ordens de remos}{adj., nom., pl., fem.}}{έφθινται τρίσκαλμοι}
\verso{680}{\notagramatical{ἄναες}{ἄναυς/\allowbreak -υν}{sem navios}{adj., nom., pl., fem.}}{νᾶες ἄναες ἄναες.}
\parte{Terceiro Episódio}{681–851}
\falante{ΕΙΔΩΛΟΝ ΔΑΡΕΙΟΥ}
\verso{681}{\notagramatical{πιστῶν}{πιστός/\allowbreak -ή/\allowbreak -όν}{}{adj., gen., pl., masc.}}{ὦ πιστὰ πιστῶν ἥλικές θ᾽ ἥβης ἐμῆς}
\verso{682}{}{Πέρσαι γεραιοί, τίνα πόλις πονεῖ πόνον;}
\verso{683}{}{στένει, κέκοπται, καὶ χαράσσεται πέδον.}
\verso{684}{\notagramatical{ἄκοιτιν}{ἄκοιτις, -ιος, ἡ}{esposa}{ac., sing.}}{λεύσσων δ᾽ ἄκοιτιν τὴν ἐμὴν τάφου πέλας}
\verso{685}{\notagramatical{ταρβῶ}{ταρβέω}{temer}{pres., ind., at., 1ª s.}\crlf \notagramatical{πρευμενὴς}{πρευμενής/\allowbreak -ές}{}{adj., nom., sing., masc.}}{ταρβῶ, χοὰς δὲ πρευμενὴς ἐδεξάμην.}
\verso{686}{}{ὑμεῖς δὲ θρηνεῖτ᾽ ἐγγὺς ἑστῶτες τάφου}
\verso{687}{\notagramatical{ψυχαγωγοῖς}{ψυχαγωγός/\allowbreak -όν}{que evoca almas}{adj., dat., pl., masc.}\crlf \notagramatical{ὀρθιάζοντες}{ὀρθιάζω}{erguer a voz}{part., pres., at., nom., pl., masc.}}{καὶ ψυχαγωγοῖς ὀρθιάζοντες γόοις}
\verso{688}{}{οἰκτρῶς καλεῖσθέ μ᾽: ἐστὶ δ᾽ οὐκ εὐέξοδον,}
\verso{689}{}{ἄλλως τε πάντως χοἰ κατὰ χθονὸς θεοὶ}
\verso{690}{\notagramatical{λαβεῖν}{λαμβάνω}{}{inf., aor., at.}\crlf \notagramatical{μεθιέναι}{μεθίημι}{deixar ir}{inf., pres., at.}}{λαβεῖν ἀμείνους εἰσὶν ἢ μεθιέναι.}
\verso{691}{\notagramatical{ἐνδυναστεύσας}{ἐνδυναστεύω}{exercer poder}{part., aor., at., nom., sing., masc.}}{ὅμως δ᾽ ἐκείνοις ἐνδυναστεύσας ἐγὼ}
\verso{692}{}{ἥκω: τάχυνε δ᾽, ὡς ἄμεμπτος ὦ χρόνου.}
\verso{693}{\notagramatical{νεοχμὸν}{νεοχμός/\allowbreak -ή/\allowbreak -όν}{novo, inesperado}{adj., nom., sing., neut.}\crlf \notagramatical{ἐμβριθὲς}{ἐμβριθής/\allowbreak -ές}{grave, opressivo}{adj., nom., sing., neut.}}{τί δ᾽ ἐστὶ Πέρσαις νεοχμὸν ἐμβριθὲς κακόν;}
\falante{ΧΟΡΟΣ}
\verso{694}{}{σέβομαι μὲν προσιδέσθαι,}
\verso{695}{}{σέβομαι δ᾽ ἀντία λέξαι}
\verso{696}{\notagramatical{τάρβει}{τάρβος, -ους, τό}{temor}{dat., sing.}}{σέθεν ἀρχαίῳ περὶ τάρβει.}
\falante{ΔΑΡΕΙΟΣ}
\verso{697}{\notagramatical{ἦλθον}{ἔρχομαι}{}{aor., ind., at., 1ª s.}\crlf \notagramatical{πεπεισμένος}{πείθω}{}{part., perf., méd./pass., nom., sing., masc.}}{ἀλλ᾽ ἐπεὶ κάτωθεν ἦλθον σοῖς γόοις πεπεισμένος,}
\verso{698}{}{μή τι μακιστῆρα μῦθον, ἀλλὰ σύντομον λέγων}
\verso{699}{\notagramatical{εἰπὲ}{λέγω}{}{aor., imper., at., 2ª s.}\crlf \notagramatical{πέραινε}{περαίνω}{concluir, relatar}{pres., imper., at., 2ª s.}\crlf \notagramatical{μεθείς}{μεθίημι}{}{part., aor., at., nom., sing., masc.}}{εἰπὲ καὶ πέραινε πάντα, τὴν ἐμὴν αἰδῶ μεθείς.}
\falante{ΧΟΡΟΣ}
\verso{700}{}{δίεμαι μὲν χαρίσασθαι,}
\verso{701}{\notagramatical{δίεμαι}{δίεμαι}{temer}{pres., ind., méd./pass., 1ª s.}}{δίεμαι δ᾽ ἀντία φάσθαι,}
\verso{702}{\notagramatical{λέξας}{λέγω}{}{part., aor., at., nom., sing., masc.}}{λέξας δύσλεκτα φίλοισιν.}
\falante{ΔΑΡΕΙΟΣ}
\verso{703}{\notagramatical{παλαιὸν}{παλαιός/\allowbreak -ά/\allowbreak -όν}{}{adj., nom., sing., neut.}\crlf \notagramatical{ἀνθίσταται}{ἀνθίστημι}{resistir}{pres., ind., méd., 3ª s.}}{ἀλλ᾽ ἐπεὶ δέος παλαιὸν σοὶ φρενῶν ἀνθίσταται,}
\verso{704}{}{τῶν ἐμῶν λέκτρων γεραιὰ ξύννομ᾽ εὐγενὲς γύναι,}
\verso{705}{\notagramatical{κλαυμάτων}{κλαῦμα, κλαύματος, τό}{pranto}{gen., pl.}\crlf \notagramatical{λήξασα}{λήγω}{}{part., aor., at., nom., sing., fem.}\crlf \notagramatical{γόων}{γόος, -ου, ὁ}{}{gen., pl.}}{κλαυμάτων λήξασα τῶνδε καὶ γόων σαφές τί μοι}
\verso{706}{}{λέξον. ἀνθρώπεια δ᾽ ἄν τοι πήματ᾽ ἂν τύχοι βροτοῖς.}
\verso{707}{\notagramatical{χέρσου}{χέρσος, -ου, ἡ}{terra firme}{gen., sing.}}{πολλὰ μὲν γὰρ ἐκ θαλάσσης, πολλὰ δ᾽ ἐκ χέρσου κακὰ}
\verso{708}{\notagramatical{γίγνεται}{γίγνομαι}{}{pres., ind., méd., 3ª s.}\crlf \notagramatical{θνητοῖς}{θνητός/\allowbreak -ή/\allowbreak -όν}{mortal}{adj., dat., pl., masc.}}{γίγνεται θνητοῖς, ὁ μάσσων βίοτος ἢν ταθῇ πρόσω.}
\falante{ΑΤΟΣΣΑ}
\verso{709}{}{ὦ βροτῶν πάντων ὑπερσχὼν ὄλβον εὐτυχεῖ πότμῳ,}
\verso{710}{\notagramatical{ἔλευσσες}{λεύσσω}{olhar}{imperf., ind., at., 2ª s.}}{ὡς ἕως τ᾽ ἔλευσσες αὐγὰς ἡλίου ζηλωτὸς ὢν}
\verso{711}{\notagramatical{διήγαγες}{διάγω}{passar a vida}{aor., ind., at., 2ª s.}}{βίοτον εὐαίωνα Πέρσαις ὡς θεὸς διήγαγες,}
\verso{712}{\notagramatical{θανόντα}{θνήσκω}{}{part., aor., at., ac., sing., masc.}}{νῦν τέ σε ζηλῶ θανόντα, πρὶν κακῶν ἰδεῖν βάθος.}
\verso{713}{}{πάντα γάρ, Δαρεῖ᾽, ἀκούσει μῦθον ἐν βραχεῖ χρόνῳ:}
\verso{714}{\notagramatical{διαπεπόρθηται}{διαπορθέω}{devastar}{perf., ind., pass., 3ª s.}\crlf \notagramatical{εἰπεῖν}{λέγω}{}{inf., aor., at.}}{διαπεπόρθηται τὰ Περσῶν πράγμαθ᾽, ὡς εἰπεῖν ἔπος.}
\falante{ΔΑΡΕΙΟΣ}
\verso{715}{}{τίνι τρόπῳ; λοιμοῦ τις ἦλθε σκηπτός, ἢ στάσις πόλει;}
\falante{ΑΤΟΣΣΑ}
\verso{716}{\notagramatical{κατέφθαρται}{καταφθείρω}{}{perf., ind., méd./pass., 3ª s.}}{οὐδαμῶς: ἀλλ᾽ ἀμφ᾽ Ἀθήνας πᾶς κατέφθαρται στρατός.}
\falante{ΔΑΡΕΙΟΣ}
\verso{717}{\notagramatical{ἐστρατηλάτει}{στρατηλατέω}{comandar exército}{imperf., ind., at., 3ª s.}}{τίς δ᾽ ἐμῶν ἐκεῖσε παίδων ἐστρατηλάτει; φράσον.}
\falante{ΑΤΟΣΣΑ}
\verso{718}{\notagramatical{κενώσας}{κενόω}{esvaziar}{part., aor., at., nom., sing., masc.}}{θούριος Ξέρξης, κενώσας πᾶσαν ἠπείρου πλάκα.}
\falante{ΔΑΡΕΙΟΣ}
\verso{719}{\notagramatical{πεζὸς}{πεζός/\allowbreak -ή/\allowbreak -όν}{}{adj., nom., sing., masc.}\crlf \notagramatical{ἐμώρανεν}{μωραίνω}{agir loucamente}{aor., ind., at., 3ª s.}}{πεζὸς ἢ ναύτης δὲ πεῖραν τήνδ᾽ ἐμώρανεν τάλας;}
\falante{ΑΤΟΣΣΑ}
\verso{720}{}{ἀμφότερα: διπλοῦν μέτωπον ἦν δυοῖν στρατευμάτοιν.}
\falante{ΔΑΡΕΙΟΣ}
\verso{721}{\notagramatical{ἤνυσεν}{ἀνύω}{realizar}{aor., ind., at., 3ª s.}\crlf \notagramatical{περᾶν}{περάω}{}{inf., pres., at.}}{πῶς δὲ καὶ στρατὸς τοσόσδε πεζὸς ἤνυσεν περᾶν;}
\falante{ΑΤΟΣΣΑ}
\verso{722}{\notagramatical{ἔζευξεν}{ζεύγνυμι}{jungir}{aor., ind., at., 3ª s.}\crlf \notagramatical{ἔχειν}{ἔχω}{}{inf., pres., at.}}{μηχαναῖς ἔζευξεν Ἕλλης πορθμόν, ὥστ᾽ ἔχειν πόρον.}
\falante{ΔΑΡΕΙΟΣ}
\verso{723}{\notagramatical{κλῇσαι}{κλείω}{fechar}{inf., aor., at.}}{καὶ τόδ᾽ ἐξέπραξεν, ὥστε Βόσπορον κλῇσαι μέγαν;}
\falante{ΑΤΟΣΣΑ}
\verso{724}{\notagramatical{δαιμόνων}{δαίμων, δαίμονος, ἡ}{}{gen., pl.}\crlf \notagramatical{ξυνήψατο}{συνάπτω}{}{aor., ind., méd., 3ª s.}}{ὧδ᾽ ἔχει: γνώμης δέ πού τις δαιμόνων ξυνήψατο.}
\falante{ΔΑΡΕΙΟΣ}
\verso{725}{}{φεῦ, μέγας τις ἦλθε δαίμων, ὥστε μὴ φρονεῖν καλῶς.}
\falante{ΑΤΟΣΣΑ}
\verso{726}{}{ὡς ἰδεῖν τέλος πάρεστιν οἷον ἤνυσεν κακόν.}
\falante{ΔΑΡΕΙΟΣ}
\verso{727}{\notagramatical{πράξασιν}{πράσσω}{}{part., aor., at., dat., pl., masc.}}{καὶ τί δὴ πράξασιν αὐτοῖς ὧδ᾽ ἐπιστενάζετε;}
\falante{ΑΤΟΣΣΑ}
\verso{728}{\notagramatical{κακωθεὶς}{κακόω}{maltratar}{part., aor., pass., nom., sing., masc.}\crlf \notagramatical{πεζὸν}{πεζός/\allowbreak -ή/\allowbreak -όν}{}{adj., ac., sing., masc.}\crlf \notagramatical{ὤλεσε}{ὄλλυμι}{}{aor., ind., at., 3ª s.}}{ναυτικὸς στρατὸς κακωθεὶς πεζὸν ὤλεσε στρατόν.}
\falante{ΔΑΡΕΙΟΣ}
\verso{729}{}{ὧδε παμπήδην δὲ λαὸς πᾶς κατέφθαρται δορί;}
\falante{ΑΤΟΣΣΑ}
\verso{730}{}{πρὸς τάδ᾽ ὡς Σούσων μὲν ἄστυ πᾶν κενανδρίαν στένειν.}
\falante{ΔΑΡΕΙΟΣ}
\verso{731}{\notagramatical{κεδνῆς}{κεδνός/\allowbreak -ή/\allowbreak -όν}{}{adj., gen., sing., fem.}\crlf \desdobramento{κἀπικουρίας}{καὶ ἐπικουρίας}}{ὦ πόποι κεδνῆς ἀρωγῆς κἀπικουρίας στρατοῦ.}
\falante{ΑΤΟΣΣΑ}
\verso{732}{}{Βακτρίων δ᾽ ἔρρει πανώλης δῆμος, οὐδέ τις γέρων.}
\falante{ΔΑΡΕΙΟΣ}
\verso{733}{}{ὦ μέλεος, οἵαν ἄρ᾽ ἥβην ξυμμάχων ἀπώλεσεν.}
\falante{ΑΤΟΣΣΑ}
\verso{734}{}{μονάδα δὲ Ξέρξην ἔρημόν φασιν οὐ πολλῶν μέτα—}
\falante{ΔΑΡΕΙΟΣ}
\verso{735}{}{πῶς τε δὴ καὶ ποῖ τελευτᾶν; ἔστι τις σωτηρία;}
\falante{ΑΤΟΣΣΑ}
\verso{736}{\notagramatical{ζευκτηρίαν}{ζευκτήριος/\allowbreak -α/\allowbreak -ον}{que liga, de junção}{adj., ac., sing., fem.}}{ἄσμενον μολεῖν γέφυραν γαῖν δυοῖν ζευκτηρίαν.}
\falante{ΔΑΡΕΙΟΣ}
\verso{737}{\notagramatical{σεσῶσθαι}{σώζω}{}{inf., perf., méd./pass.}\crlf \notagramatical{ἐτήτυμον}{ἐτήτυμος/\allowbreak -ον}{verdadeiro}{adj., nom., sing., neut.}}{καὶ πρὸς ἤπειρον σεσῶσθαι τήνδε, τοῦτ᾽ ἐτήτυμον;}
\falante{ΑΤΟΣΣΑ}
\verso{738}{}{ναί: λόγος κρατεῖ σαφηνὴς τοῦτό γ᾽, οὐδ᾽ ἔνι στάσις.}
\falante{ΔΑΡΕΙΟΣ}
\verso{739}{}{φεῦ, ταχεῖά γ᾽ ἦλθε χρησμῶν πρᾶξις, ἐς δὲ παῖδ᾽ ἐμὸν}
\verso{740}{\notagramatical{θεσφάτων}{θέσφατος/\allowbreak -ον}{anunciado pelos deuses}{adj., gen., pl., neut.}}{Ζεὺς ἐπέσκηψεν τελευτὴν θεσφάτων: ἐγὼ δέ που}
\verso{741}{\notagramatical{ηὔχουν}{αὐχέω}{vangloriar-se}{imperf., ind., at., 1ª s.}}{διὰ μακροῦ χρόνου τάδ᾽ ηὔχουν ἐκτελευτήσειν θεούς:}
\verso{742}{\desdobramento{χὠ}{καὶ ὁ}\crlf \notagramatical{συνάπτεται}{συνάπτω}{}{pres., ind., méd./pass., 3ª s.}}{ἀλλ᾽ ὅταν σπεύδῃ τις αὐτός, χὠ θεὸς συνάπτεται.}
\verso{743}{}{νῦν κακῶν ἔοικε πηγὴ πᾶσιν ηὑρῆσθαι φίλοις.}
\verso{744}{\notagramatical{κατειδὼς}{κάτοιδα}{saber bem}{part., perf., at., nom., sing., masc.}}{παῖς δ᾽ ἐμὸς τάδ᾽ οὐ κατειδὼς ἤνυσεν νέῳ θράσει:}
\verso{745}{}{ὅστις Ἑλλήσποντον ἱρὸν δοῦλον ὣς δεσμώμασιν}
\verso{746}{\notagramatical{σχήσειν}{ἔχω}{}{inf., fut., at.}}{ἤλπισε σχήσειν ῥέοντα, Βόσπορον ῥόον θεοῦ:}
\verso{747}{\notagramatical{μετερρύθμιζε}{μεταρρυθμίζω}{transformar}{imperf., ind., at., 3ª s.}\crlf \notagramatical{σφυρηλάτοις}{σφυρήλατος/\allowbreak -ον}{forjado a martelo}{adj., dat., pl., fem.}}{καὶ πόρον μετερρύθμιζε, καὶ πέδαις σφυρηλάτοις}
\verso{748}{}{περιβαλὼν πολλὴν κέλευθον ἤνυσεν πολλῷ στρατῷ,}
\verso{749}{\notagramatical{θνητὸς}{θνητός/\allowbreak -ή/\allowbreak -όν}{}{adj., nom., sing., masc.}}{θνητὸς ὢν δὲ θεῶν ἁπάντων ᾤετ᾽, οὐκ εὐβουλίᾳ,}
\verso{750}{\notagramatical{κρατήσειν}{κρατέω}{}{inf., fut., at.}}{καὶ Ποσειδῶνος κρατήσειν. πῶς τάδ᾽ οὐ νόσος φρενῶν}
\verso{751}{}{εἶχε παῖδ᾽ ἐμόν; δέδοικα μὴ πολὺς πλούτου πόνος}
\verso{752}{}{οὑμὸς ἀνθρώποις γένηται τοῦ φθάσαντος ἁρπαγή.}
\falante{ΑΤΟΣΣΑ}
\verso{753}{}{ταῦτά τοι κακοῖς ὁμιλῶν ἀνδράσιν διδάσκεται}
\verso{754}{}{θούριος Ξέρξης: λέγουσι δ᾽ ὡς σὺ μὲν μέγαν τέκνοις}
\verso{755}{\notagramatical{ἀνανδρίας}{ἀνανδρία, -ας, ἡ}{covardia}{gen. sing.}\crlf \preposicaoposposta{ὕπο}{gen.}{por causa de}}{πλοῦτον ἐκτήσω ξὺν αἰχμῇ, τὸν δ᾽ ἀνανδρίας ὕπο}
\verso{756}{}{ἔνδον αἰχμάζειν, πατρῷον δ᾽ ὄλβον οὐδὲν αὐξάνειν.}
\verso{757}{\notagramatical{ὀνείδη}{ὄνειδος, -ους, τό}{desonra}{ac., pl.}}{τοιάδ᾽ ἐξ ἀνδρῶν ὀνείδη πολλάκις κλύων κακῶν}
\verso{758}{}{τήνδ᾽ ἐβούλευσεν κέλευθον καὶ στράτευμ᾽ ἐφ᾽ Ἑλλάδα.}
\falante{ΔΑΡΕΙΟΣ}
\verso{759}{\notagramatical{ἐξειργασμένον}{ἐξεργάζομαι}{realizar}{part., perf., méd., nom., sing., neut.}}{τοιγάρ σφιν ἔργον ἐστὶν ἐξειργασμένον}
\verso{760}{\notagramatical{ἀείμνηστον}{ἀείμνηστος/\allowbreak -ον}{sempre lembrado}{adj., ac., sing., neut.}}{μέγιστον, ἀείμνηστον, οἷον οὐδέπω}
\verso{761}{\notagramatical{ἐξεκείνωσεν}{ἐκκενόω}{esvaziar}{aor., ind., at., 3ª s.}}{τόδ᾽ ἄστυ Σούσων ἐξεκείνωσεν πεσόν,}
\verso{762}{\notagramatical{ὤπασεν}{ὀπάζω}{conceder}{aor., ind., at., 3ª s.}}{ἐξ οὗτε τιμὴν Ζεὺς ἄναξ τήνδ᾽ ὤπασεν,}
\verso{763}{\notagramatical{μηλοτρόφου}{μηλοτρόφος/\allowbreak -ον}{que cria rebanhos}{adj., gen., sing., fem.}}{ἕν᾽ ἄνδρ᾽ ἁπάσης Ἀσίδος μηλοτρόφου}
\verso{764}{\notagramatical{ταγεῖν}{ταγέω}{governar}{inf., pres., at.}\crlf \notagramatical{ἔχοντα}{ἔχω}{}{part., pres., at., ac., sing., masc.}\crlf \notagramatical{εὐθυντήριον}{εὐθυντήριος/\allowbreak -α/\allowbreak -ον}{que dirige}{adj., ac., sing., neut.}}{ταγεῖν, ἔχοντα σκῆπτρον εὐθυντήριον.}
\verso{765}{}{Μῆδος γὰρ ἦν ὁ πρῶτος ἡγεμὼν στρατοῦ:}
\verso{766}{}{ἄλλος δ᾽ ἐκείνου παῖς τόδ᾽ ἔργον ἤνυσεν:}
\verso{767}{\notagramatical{θυμὸν}{θυμός, -οῦ, ὁ}{ânimo, coração}{ac., sing.}\crlf \notagramatical{ᾠακοστρόφουν}{οἰακοστροφέω}{manejar o leme}{imperf., ind., at., 3ª p.}}{φρένες γὰρ αὐτοῦ θυμὸν ᾠακοστρόφουν.}
\verso{768}{}{τρίτος δ᾽ ἀπ᾽ αὐτοῦ Κῦρος, εὐδαίμων ἀνήρ,}
\verso{769}{\notagramatical{ἄρξας}{ἄρχω}{}{part., aor., at., nom., sing., masc.}}{ἄρξας ἔθηκε πᾶσιν εἰρήνην φίλοις:}
\verso{770}{}{Λυδῶν δὲ λαὸν καὶ Φρυγῶν ἐκτήσατο,}
\verso{771}{\notagramatical{ἤλασεν}{ἐλαύνω}{conduzir, avançar}{aor., ind., at., 3ª s.}}{Ἰωνίαν τε πᾶσαν ἤλασεν βίᾳ.}
\verso{772}{\notagramatical{ἤχθηρεν}{ἐχθαίρω}{odiar}{aor., ind., at., 3ª s.}}{θεὸς γὰρ οὐκ ἤχθηρεν, ὡς εὔφρων ἔφυ.}
\verso{773}{}{Κύρου δὲ παῖς τέταρτος ηὔθυνε στρατόν.}
\verso{774}{\notagramatical{ἦρξεν}{ἔρδω}{}{aor., ind., at., 3ª s.}}{πέμπτος δὲ Μάρδος ἦρξεν, αἰσχύνη πάτρᾳ}
\verso{775}{}{θρόνοισί τ᾽ ἀρχαίοισι: τὸν δὲ σὺν δόλῳ}
\verso{776}{}{Ἀρταφρένης ἔκτεινεν ἐσθλὸς ἐν δόμοις,}
\verso{777}{}{ξὺν ἀνδράσιν φίλοισιν, οἷς τόδ᾽ ἦν χρέος.}
\verso{778}{}{[ἕκτος δὲ Μάραφις, ἕβδομος δ᾽ Ἀρταφρένης.]}
\verso{779}{\desdobramento{κἀγὼ}{καὶ ἐγώ}\crlf \notagramatical{πάλου}{πάλος, -ου, ὁ}{sorteio}{gen., sing.}}{κἀγὼ πάλου τ᾽ ἔκυρσα τοῦπερ ἤθελον,}
\verso{780}{\desdobramento{κἀπεστράτευσα}{καὶ ἐπεστράτευσα}}{κἀπεστράτευσα πολλὰ σὺν πολλῷ στρατῷ.}
\verso{781}{}{ἀλλ᾽ οὐ κακὸν τοσόνδε προσέβαλον πόλει.}
\verso{782}{}{Ξέρξης δ᾽ ἐμὸς παῖς ὢν νέος νέα φρονεῖ,}
\verso{783}{\desdobramento{κοὐ}{καὶ οὐ}}{κοὐ μνημονεύει τὰς ἐμὰς ἐπιστολάς:}
\verso{784}{\notagramatical{ἴστ᾽}{οἶδα}{}{perf., ind., at., 2ª p.}}{εὖ γὰρ σαφῶς τόδ᾽ ἴστ᾽, ἐμοὶ ξυνήλικες,}
\verso{785}{\notagramatical{ἔσχομεν}{ἔχω}{}{aor., ind., at., 1ª p.}}{ἅπαντες ἡμεῖς, οἳ κράτη τάδ᾽ ἔσχομεν,}
\verso{786}{\notagramatical{φανεῖμεν}{φαίνω}{mostrar, aparecer}{aor., opt., pass., 1ª p.; οὐκ ἄν + opt.: não apareceríamos}\crlf \notagramatical{πήματ᾽}{πῆμα, πήματος, τό}{}{ac., pl.}\crlf \notagramatical{ἔρξαντες}{ἔρδω}{}{part., aor., at., nom., pl., masc.}}{οὐκ ἂν φανεῖμεν πήματ᾽ ἔρξαντες τόσα.}
\falante{ΧΟΡΟΣ}
\verso{787}{\notagramatical{ἄναξ}{ἄναξ, ἄνακτος, ὁ}{}{voc., sing.}\crlf \notagramatical{καταστρέφεις}{καταστρέφω}{submeter}{pres., ind., at., 2ª s.}}{τί οὖν, ἄναξ Δαρεῖε, ποῖ καταστρέφεις}
\verso{788}{}{λόγων τελευτήν; πῶς ἂν ἐκ τούτων ἔτι}
\verso{789}{\notagramatical{πράσσοιμεν}{πράσσω}{}{pres., opt., at., 1ª p.; πῶς ἄν + opt.: como poderíamos agir?}}{πράσσοιμεν ὡς ἄριστα Περσικὸς λεώς;}
\falante{ΔΑΡΕΙΟΣ}
\verso{790}{}{εἰ μὴ στρατεύοισθ᾽ ἐς τὸν Ἑλλήνων τόπον,}
\verso{791}{}{μηδ᾽ εἰ στράτευμα πλεῖον ᾖ τὸ Μηδικόν.}
\verso{792}{\notagramatical{πέλει}{πέλω}{}{pres., ind., at., 3ª s.}}{αὐτὴ γὰρ ἡ γῆ ξύμμαχος κείνοις πέλει.}
\falante{ΧΟΡΟΣ}
\verso{793}{\notagramatical{ἔλεξας}{λέγω}{}{aor., ind., at., 2ª s.}}{πῶς τοῦτ᾽ ἔλεξας, τίνι τρόπῳ δὲ συμμαχεῖ;}
\falante{ΔΑΡΕΙΟΣ}
\verso{794}{}{κτείνουσα λιμῷ τοὺς ὑπερπόλλους ἄγαν.}
\falante{ΧΟΡΟΣ}
\verso{795}{\notagramatical{εὐσταλῆ}{εὐσταλής/\allowbreak -ές}{bem equipado}{adj., ac., sing., masc.}}{ἀλλ᾽ εὐσταλῆ τοι λεκτὸν ἀροῦμεν στόλον.}
\falante{ΔΑΡΕΙΟΣ}
\verso{796}{}{ἀλλ᾽ οὐδ᾽ ὁ μείνας νῦν ἐν Ἑλλάδος τόποις}
\verso{797}{\notagramatical{κυρήσει}{κυρέω}{}{fut., ind., at., 3ª s.}\crlf \notagramatical{νοστίμου}{νόστιμος/\allowbreak -ον}{}{adj., gen., sing., fem.}}{στρατὸς κυρήσει νοστίμου σωτηρίας.}
\falante{ΧΟΡΟΣ}
\verso{798}{}{πῶς εἶπας; οὐ γὰρ πᾶν στράτευμα βαρβάρων}
\verso{799}{}{περᾷ τὸν Ἕλλης πορθμὸν Εὐρώπης ἄπο;}
\falante{ΔΑΡΕΙΟΣ}
\verso{800}{\notagramatical{παῦροί}{παῦρος/\allowbreak -η/\allowbreak -ον}{pouco}{adj., nom., pl., masc.}}{παῦροί γε πολλῶν, εἴ τι πιστεῦσαι θεῶν}
\verso{801}{\notagramatical{θεσφάτοισιν}{θέσφατος/\allowbreak -ον}{}{adj., dat., pl., neut.}\crlf \notagramatical{πεπραγμένα}{πράσσω}{}{part., perf., méd./pass., ac., pl., neut.}}{χρὴ θεσφάτοισιν, ἐς τὰ νῦν πεπραγμένα}
\verso{802}{}{βλέψαντα: συμβαίνει γὰρ οὐ τὰ μέν, τὰ δ᾽ οὔ.}
\verso{803}{}{κεἴπερ τάδ᾽ ἐστί, πλῆθος ἔκκριτον στρατοῦ}
\verso{804}{\notagramatical{λείπει}{λείπω}{}{pres., ind., at., 3ª s.}}{λείπει κεναῖσιν ἐλπίσιν πεπεισμένος.}
\verso{805}{}{μίμνουσι δ᾽ ἔνθα πεδίον Ἀσωπὸς ῥοαῖς}
\verso{806}{\notagramatical{πίασμα}{πίασμα, πιάσματος, τό}{substância que fertiliza}{nom., sing.}}{ἄρδει, φίλον πίασμα Βοιωτῶν χθονί:}
\verso{807}{}{οὗ σφιν κακῶν ὕψιστ᾽ ἐπαμμένει παθεῖν,}
\verso{808}{\notagramatical{ὕβρεως}{ὕβρις, ὕβρεως, ἡ}{desmedida, insolência}{gen., sing.}\crlf \desdobramento{κἀθέων}{καὶ ἀθέων}}{ὕβρεως ἄποινα κἀθέων φρονημάτων:}
\verso{809}{\notagramatical{μολόντες}{βλώσκω}{}{part., aor., at., nom., pl., masc.}}{οἳ γῆν μολόντες Ἑλλάδ᾽ οὐ θεῶν βρέτη}
\verso{810}{\notagramatical{συλᾶν}{συλάω}{saquear}{inf., pres., at.}\crlf \notagramatical{πιμπράναι}{πίμπρημι}{incendiar}{inf., pres., at.}}{ᾐδοῦντο συλᾶν οὐδὲ πιμπράναι νεώς:}
\verso{811}{}{βωμοὶ δ᾽ ἄιστοι, δαιμόνων θ᾽ ἱδρύματα}
\verso{812}{\notagramatical{πρόρριζα}{πρόρριζος/\allowbreak -ον}{desde a raiz}{adj., nom., pl., neut.}\crlf \notagramatical{ἐξανέστραπται}{ἐξαναστρέφω}{derrubar inteiramente}{perf., ind., méd./pass., 3ª s.}}{πρόρριζα φύρδην ἐξανέστραπται βάθρων.}
\verso{813}{}{τοιγὰρ κακῶς δράσαντες οὐκ ἐλάσσονα}
\verso{814}{\desdobramento{κοὐδέπω}{καὶ οὐδέπω}}{πάσχουσι, τὰ δὲ μέλλουσι, κοὐδέπω κακῶν}
\verso{815}{\notagramatical{ἀπέσβηκ᾽}{ἀποσβέννυμι}{extinguir}{perf., ind., at., 3ª s.}\crlf \notagramatical{ἐκπιδύεται}{ἐκπιδύομαι}{jorrar}{pres., ind., méd., 3ª s.}}{κρηνὶς ἀπέσβηκ᾽, ἀλλ᾽ ἔτ᾽ ἐκπιδύεται.}
\verso{816}{}{τόσος γὰρ ἔσται πέλανος αἱματοσταγὴς}
\verso{817}{}{πρὸς γῇ Πλαταιῶν Δωρίδος λόγχης ὕπο:}
\verso{818}{\notagramatical{θῖνες}{θίς, θινός, ἡ}{montão}{nom., pl.}\crlf \notagramatical{τριτοσπόρῳ}{τριτόσπορος/\allowbreak -ον}{de terceira geração}{adj., dat., sing., fem.}}{θῖνες νεκρῶν δὲ καὶ τριτοσπόρῳ γονῇ}
\verso{819}{}{ἄφωνα σημανοῦσιν ὄμμασιν βροτῶν}
\verso{820}{\notagramatical{ὄντα}{εἰμί}{}{part., pres., at., ac., sing., masc.}}{ὡς οὐχ ὑπέρφευ θνητὸν ὄντα χρὴ φρονεῖν.}
\verso{821}{\notagramatical{ἐξανθοῦσ᾽}{ἐξανθέω}{florescer}{part., pres., at., nom., pl., fem.}\crlf \notagramatical{ἐκάρπωσεν}{καρπόω}{produzir fruto}{aor., ind., at., 3ª s.}}{ὕβρις γὰρ ἐξανθοῦσ᾽ ἐκάρπωσεν στάχυν}
\verso{822}{}{ἄτης, ὅθεν πάγκλαυτον ἐξαμᾷ θέρος.}
\verso{823}{}{τοιαῦθ᾽ ὁρῶντες τῶνδε τἀπιτίμια}
\verso{824}{\notagramatical{μέμνησθ᾽}{μιμνήσκω}{}{perf., imper., méd., 2ª p.; + gen.: lembrar-se de; Ἀθηνῶν, Ἑλλάδος}}{μέμνησθ᾽ Ἀθηνῶν Ἑλλάδος τε, μηδέ τις}
\verso{825}{}{ὑπερφρονήσας τὸν παρόντα δαίμονα}
\verso{826}{\notagramatical{ἐκχέῃ}{ἐκχέω}{derramar}{aor., subj., at., 3ª s.; μηδέ + subj.: que ninguém dissipe; proibição}}{ἄλλων ἐρασθεὶς ὄλβον ἐκχέῃ μέγαν.}
\verso{827}{\notagramatical{κολαστὴς}{κολαστής, -οῦ, ὁ}{castigador}{nom., sing.}\crlf \notagramatical{ὑπερκόμπων}{ὑπέρκομπος/\allowbreak -ον}{demasiado arrogante}{adj., gen., pl., neut.}}{Ζεύς τοι κολαστὴς τῶν ὑπερκόμπων ἄγαν}
\verso{828}{}{φρονημάτων ἔπεστιν, εὔθυνος βαρύς.}
\verso{829}{\notagramatical{σωφρονεῖν}{σωφρονέω}{ser prudente}{inf., pres., at.}}{πρὸς ταῦτ᾽ ἐκεῖνον, σωφρονεῖν κεχρημένον,}
\verso{830}{\notagramatical{νουθετήμασιν}{νουθέτημα, νουθετήματος, τό}{advertência}{dat., pl.}}{πινύσκετ᾽ εὐλόγοισι νουθετήμασιν,}
\verso{831}{\notagramatical{λῆξαι}{λήγω}{}{inf., aor., at.}\crlf \notagramatical{θεοβλαβοῦνθ᾽}{θεοβλαβέω}{ser ferido por um deus}{part., pres., at., ac., sing., masc.}\crlf \notagramatical{ὑπερκόμπῳ}{ὑπέρκομπος/\allowbreak -ον}{}{adj., dat., sing., neut.}}{λῆξαι θεοβλαβοῦνθ᾽ ὑπερκόμπῳ θράσει.}
\verso{832}{}{σὺ δ᾽, ὦ γεραιὰ μῆτερ ἡ Ξέρξου φίλη,}
\verso{833}{\notagramatical{ἐλθοῦσ᾽}{ἔρχομαι}{}{part., aor., at., nom., sing., fem.}}{ἐλθοῦσ᾽ ἐς οἴκους κόσμον ὅστις εὐπρεπὴς}
\verso{834}{\notagramatical{λαβοῦσ᾽}{λαμβάνω}{}{part., aor., at., nom., sing., fem.}}{λαβοῦσ᾽ ὑπαντίαζε παιδί. παντὶ γὰρ}
\verso{835}{\notagramatical{λακίδες}{λακίς, λακίδος, ἡ}{rasgão}{nom., pl.}}{κακῶν ὑπ᾽ ἄλγους λακίδες ἀμφὶ σώματι}
\verso{836}{\notagramatical{στημορραγοῦσι}{στημορραγέω}{rasgar as vestes do peito}{pres., ind., at., 3ª p.}\crlf \notagramatical{ἐσθημάτων}{ἔσθημα, ἐσθήματος, τό}{veste}{gen., pl.}}{στημορραγοῦσι ποικίλων ἐσθημάτων.}
\verso{837}{}{ἀλλ᾽ αὐτὸν εὐφρόνως σὺ πράυνον λόγοις:}
\verso{838}{}{μόνης γάρ, οἶδα, σοῦ κλύων ἀνέξεται.}
\verso{839}{\notagramatical{ζόφον}{ζόφος, -ου, ὁ}{trevas}{ac., sing.}}{ἐγὼ δ᾽ ἄπειμι γῆς ὑπὸ ζόφον κάτω.}
\verso{840}{}{ὑμεῖς δέ, πρέσβεις, χαίρετ᾽, ἐν κακοῖς ὅμως}
\verso{841}{}{ψυχὴν διδόντες ἡδονῇ καθ᾽ ἡμέραν,}
\verso{842}{\notagramatical{θανοῦσι}{θνήσκω}{}{part., aor., at., dat., pl., masc.}}{ὡς τοῖς θανοῦσι πλοῦτος οὐδὲν ὠφελεῖ.}
\falante{ΧΟΡΟΣ}
\verso{843}{\notagramatical{παρόντα}{πάρειμι}{}{part., pres., at., ac., pl., neut.}}{ἦ πολλὰ καὶ παρόντα καὶ μέλλοντ᾽ ἔτι}
\verso{844}{\notagramatical{ἀκούσας}{ἀκούω}{}{part., aor., at., nom., sing., masc.}\crlf \notagramatical{βαρβάροισι}{βάρβαρος/\allowbreak -ον}{}{adj., dat., pl., masc.}\crlf \notagramatical{πήματα}{πῆμα, πήματος, τό}{}{ac., pl.}}{ἤλγησ᾽ ἀκούσας βαρβάροισι πήματα.}
\falante{ΑΤΟΣΣΑ}
\verso{845}{}{ὦ δαῖμον, ὥς με πόλλ᾽ ἐσέρχεται κακῶν}
\verso{846}{}{ἄλγη, μάλιστα δ᾽ ἥδε συμφορὰ δάκνει,}
\verso{847}{}{ἀτιμίαν γε παιδὸς ἀμφὶ σώματι}
\verso{848}{\notagramatical{κλύουσαν}{κλύω}{}{part., pres., at., ac., sing., fem.}\crlf \notagramatical{ἀμπέχει}{ἀμπέχω}{cobrir, vestir}{pres., ind., at., 3ª s.}}{ἐσθημάτων κλύουσαν, ἥ νιν ἀμπέχει.}
\verso{849}{}{ἀλλ᾽ εἶμι, καὶ λαβοῦσα κόσμον ἐκ δόμων}
\verso{850}{}{ὑπαντιάζειν παιδί μου πειράσομαι.}
\verso{851}{\notagramatical{προδώσομεν}{προδίδωμι}{abandonar, trair}{fut., ind., at., 1ª p.}}{οὐ γὰρ τὰ φίλτατ᾽ ἐν κακοῖς προδώσομεν.}
\parte{Terceiro Estásimo}{852–907}
\falante{ΧΟΡΟΣ}
\verso{852}{}{ὦ πόποι ἦ μεγάλας ἀγαθᾶς τε πο-}
\verso{853}{\notagramatical{λισσονόμου}{πολισσονόμος/\allowbreak -ον}{que governa cidades}{adj., gen., sing., fem.}\crlf \notagramatical{ἐπεκύρσαμεν}{ἐπικύρω}{alcançar, obter}{aor., ind., at., 1ª p.}}{λισσονόμου βιοτᾶς ἐπεκύρσαμεν,}
\verso{854}{}{εὖθ᾽ ὁ γηραιὸς}
\verso{855}{\notagramatical{πανταρκὴς}{πανταρκής/\allowbreak -ές}{todo-poderoso}{adj., nom., sing., masc.}}{πανταρκὴς ἀκάκας}
\verso{856}{}{ἄμαχος βασιλεὺς}
\verso{857}{\notagramatical{ἆρχε}{ἄρχω}{}{imperf., ind., at., 3ª s.}}{ἰσόθεος Δαρεῖος ἆρχε χώρας.}
\pausaestrofe
\verso{858}{\notagramatical{στρατιὰς}{στρατιά, -ᾶς, ἡ}{}{ac., pl.}}{πρῶτα μὲν εὐδοκίμους στρατιὰς ἀπε-}
\verso{859}{}{φαινόμεθ᾽, ἠδὲ πολίσματα πύργινα}
\verso{860}{}{πάντ᾽ ἐπέρθομεν.}
\verso{861}{\notagramatical{νόστοι}{νόστος, -ου, ὁ}{retorno}{nom., pl.}}{νόστοι δ᾽ ἐκ πολέμων}
\verso{862}{\notagramatical{ἀπαθεῖς}{ἀπαθής/\allowbreak -ές}{sem sofrimento}{adj., ac., pl., masc.}}{ἀπόνους ἀπαθεῖς}
\verso{863}{\notagramatical{πράσσοντας}{πράσσω}{}{part., pres., at., ac., pl., masc.}}{<ἀνέρασ> εὖ πράσσοντας ἆγον οἴκους.}
\pausaestrofe
\verso{864}{\notagramatical{εἷλε}{αἱρέω}{}{aor., ind., at., 3ª s.}\crlf \notagramatical{πόρον}{πόρω}{dar, proporcionar}{aor., ind., at., 1ª s.}}{ὅσσας δ᾽ εἷλε πόλεις πόρον}
\verso{865}{}{οὐ διαβὰς Ἅλυος ποταμοῖο,}
\verso{866}{\notagramatical{συθείς}{σεύω}{}{part., aor., pass., nom., sing., masc.}}{οὐδ᾽ ἀφ᾽ ἑστίας συθείς,}
\verso{867}{}{οἷαι Στρυμονίου πελά-}
\verso{868}{}{γους Ἀχελωίδες εἰσὶ πάροικοι}
\verso{870}{\notagramatical{ἐπαύλων}{ἔπαυλος, -ου, ὁ}{morada}{gen., pl.}}{Θρῃκίων ἐπαύλων,}
\pausaestrofe
\verso{871}{}{λίμνας τ᾽ ἔκτοθεν αἳ κατὰ}
\verso{872}{\notagramatical{ἐληλαμέναι}{ἐλαύνω}{}{part., perf., méd./pass., nom., pl., fem.}}{χέρσον ἐληλαμέναι πέρι πύργον}
\verso{873}{}{τοῦδ᾽ ἄνακτος ἄιον,}
\verso{875}{}{Ἕλλας τ᾽ ἀμφὶ πόρον πλατὺν}
\verso{876}{\notagramatical{μυχία}{μύχιος/\allowbreak -α/\allowbreak -ον}{interior, recôndito}{adj., nom., sing., fem.}}{εὐχόμεναι, μυχία τε Προποντίς,}
\verso{877}{\notagramatical{στόμωμα}{στόμωμα, στομώματος, τό}{boca, entrada}{nom., sing.}}{καὶ στόμωμα Πόντου:}
\pausaestrofe
\verso{878}{}{νᾶσοί θ᾽ αἳ κατὰ πρῶν᾽}
\verso{880}{}{ἅλιον περίκλυστοι}
\verso{881}{\notagramatical{προσήμεναι}{πρόσημαι}{situar-se junto}{part., pres., méd., nom. pl. fem.}}{τᾷδε γᾷ προσήμεναι}
\verso{882}{}{οἵα Λέσβος ἐλαι-}
\verso{883}{}{όφυτός τε Σάμος, Χίος,}
\verso{884}{}{ἠδὲ Πάρος, Νάξος, Μύκο-}
\verso{885}{\notagramatical{συνάπτουσ᾽}{συνάπτω}{}{part., pres., at., nom., sing., fem.}}{νος, Τήνῳ τε συνάπτουσ᾽}
\verso{886}{\notagramatical{ἀγχιγείτων}{ἀγχιγείτων/\allowbreak -ον}{vizinho próximo}{adj., nom., sing., fem.}}{Ἄνδρος ἀγχιγείτων,}
\pausaestrofe
\verso{887}{\notagramatical{ἀγχιάλους}{ἀγχίαλος/\allowbreak -ον}{junto ao mar}{adj., ac., pl., fem.}}{καὶ τὰς ἀγχιάλους}
\verso{888}{\notagramatical{ἐκράτυνε}{κρατύνω}{dominar}{imperf., ind., at., 3ª s.}}{ἐκράτυνε μεσάκτους,}
\verso{890}{}{Λῆμνον, Ἰκάρου θ᾽ ἕδος,}
\verso{891}{}{καὶ Ῥόδον ἠδὲ Κνίδον}
\verso{892}{}{Κυπρίας τε πόλεις, Πάφον,}
\verso{893}{}{ἠδὲ Σόλους, Σαλαμῖνά τε,}
\verso{894}{}{τᾶς νῦν ματρόπολις τῶνδ᾽}
\verso{896}{}{αἰτία στεναγμῶν.}
\pausaestrofe
\verso{897}{\notagramatical{εὐκτεάνους}{εὐκτέανος/\allowbreak -ον}{rico em bens}{adj., ac., pl., fem.}}{καὶ τὰς εὐκτεάνους κατὰ}
\verso{898}{}{κλῆρον Ἰαόνιον πολυάνδρους}
\verso{900}{\notagramatical{ἐκράτει}{κρατέω}{}{imperf., ind., at., 3ª s.}}{Ἑλλάνων ἐκράτει σφετέραις φρεσίν.}
\verso{901}{}{ἀκάματον δὲ παρῆν σθένος}
\verso{902}{\notagramatical{τευχηστήρων}{τευχηστήρ, τευχηστῆρος, ὁ}{guerreiro armado}{gen., pl.}}{ἀνδρῶν τευχηστήρων}
\verso{903}{\notagramatical{παμμίκτων}{παμμιγής/\allowbreak -ές}{}{adj., gen., pl., masc.}}{παμμίκτων τ᾽ ἐπικούρων.}
\verso{904}{}{νῦν δ᾽ οὐκ ἀμφιλόγως}
\verso{905}{\notagramatical{θεότρεπτα}{θεότρεπτος/\allowbreak -ον}{causado pelos deuses}{adj., ac., pl., neut.}}{θεότρεπτα τάδ᾽ αὖ}
\verso{906}{}{φέρομεν πολέμοισι}
\verso{907}{\notagramatical{δμαθέντες}{δαμάζω}{subjugar}{part., aor., pass., nom., pl., masc.}}{δμαθέντες μεγάλως}
\parte{Éxodo}{908–1077}
\verso{908}{\notagramatical{ποντίαισιν}{πόντιος/\allowbreak -α/\allowbreak -ον}{}{adj., dat., pl., fem.}}{πλαγαῖσι ποντίαισιν.}
\falante{ΞΕΡΞΗΣ}
\verso{909}{}{ἰώ, δύστηνος ἐγὼ στυγερᾶς μοίρας}
\verso{910}{\notagramatical{κυρήσας}{κυρέω}{}{part., aor., at., nom., sing., masc.}\crlf \notagramatical{ἀτεκμαρτοτάτης}{ἀτέκμαρτος/\allowbreak -ον}{imprevisível}{adj., superl., gen., sing., fem.}}{τῆσδε κυρήσας ἀτεκμαρτοτάτης,}
\verso{911}{}{ὡς ὠμοφρόνως δαίμων ἐνέβη}
\verso{912}{\notagramatical{τλήμων}{τλήμων/\allowbreak -ον}{sofredor}{adj., nom., sing., masc.}}{Περσῶν γενεᾷ: τί πάθω τλήμων;}
\verso{913}{\notagramatical{γυίων}{γυῖον, -ου, τό}{membro do corpo}{gen., pl.}}{λέλυται γὰρ ἐμοὶ γυίων ῥώμη}
\verso{914}{\notagramatical{ἐσιδόντ᾽}{εἰσοράω}{}{part., aor., at., ac., sing., masc.}}{τήνδ᾽ ἡλικίαν ἐσιδόντ᾽ ἀστῶν.}
\verso{915}{\desdobramento{κἀμὲ}{καὶ ἐμέ}}{εἴθ᾽ ὄφελεν, Ζεῦ, κἀμὲ μετ᾽ ἀνδρῶν}
\verso{916}{}{τῶν οἰχομένων}
\verso{917}{}{θανάτου κατὰ μοῖρα καλύψαι.}
\falante{ΧΟΡΟΣ}
\verso{918}{}{ὀτοτοῖ, βασιλεῦ, στρατιᾶς ἀγαθῆς}
\verso{919}{}{καὶ περσονόμου τιμῆς μεγάλης,}
\verso{920}{}{κόσμου τ᾽ ἀνδρῶν,}
\verso{921}{\notagramatical{ἐπέκειρεν}{ἐπικείρω}{ceifar}{aor., ind., at., 3ª s.}}{οὓς νῦν δαίμων ἐπέκειρεν.}
\falante{ΞΕΡΞΗΣ}
\verso{922}{\notagramatical{αἰάζει}{αἰάζω}{lamentar}{pres., ind., at., 3ª s.}\crlf \notagramatical{ἐγγαίαν}{ἔγγαιος/\allowbreak -ον}{subterrâneo}{adj., ac., sing., fem.}}{γᾶ δ᾽ αἰάζει τὰν ἐγγαίαν}
\verso{923}{}{ἥβαν Ξέρξᾳ κταμέναν Ἅιδου}
\verso{924}{}{σάκτορι Περσᾶν: ᾁδοβάται γὰρ}
\verso{925}{\notagramatical{φῶτες}{φώς, φωτός, ὁ}{homem}{nom., pl.}}{πολλοὶ φῶτες, χώρας ἄνθος,}
\verso{926}{\notagramatical{ταρφύς}{ταρφύς/\allowbreak -εῖα/\allowbreak -ύ}{numeroso, denso}{adj., nom. pl. masc.}}{τοξοδάμαντες, πάνυ ταρφύς τις}
\verso{927}{\notagramatical{ἐξέφθινται}{ἐκφθίνω}{perecer}{perf., ind., méd., 3ª p.}}{μυριὰς ἀνδρῶν, ἐξέφθινται.}
\verso{928}{\notagramatical{κεδνᾶς}{κεδνός/\allowbreak -ή/\allowbreak -όν}{}{adj., gen., sing., fem.}}{αἰαῖ αἰαῖ κεδνᾶς ἀλκᾶς.}
\verso{929}{}{Ἀσία δὲ χθών, βασιλεῦ γαίας,}
\verso{930}{}{αἰνῶς αἰνῶς}
\verso{931}{}{ἐπὶ γόνυ κέκλιται.}
\pausaestrofe
\verso{932}{\notagramatical{αἰακτὸς}{αἰακτός/\allowbreak -ή/\allowbreak -όν}{lamentável}{adj., nom., sing., masc.}}{ὅδ᾽ ἐγών, οἰοῖ, αἰακτὸς}
\verso{933}{}{μέλεος γέννᾳ γᾷ τε πατρῴᾳ}
\verso{934}{\notagramatical{ἐγενόμαν}{γίγνομαι}{}{aor., ind., méd., 1ª s.}}{κακὸν ἄρ᾽ ἐγενόμαν.}
\falante{ΧΟΡΟΣ}
\verso{935}{}{πρόσφθογγόν σοι νόστου τὰν}
\verso{936}{}{κακοφάτιδα βοάν,}
\verso{937}{\notagramatical{κακομέλετον}{κακομέλετος/\allowbreak -ον}{cantado em tom de lamento}{adj., nom., sing., neut.}}{κακομέλετον ἰὰν}
\verso{938}{}{Μαριανδυνοῦ θρηνητῆρος}
\verso{939}{\notagramatical{πέμψω}{πέμπω}{}{fut., ind., at., 1ª s.}}{πέμψω πέμψω,}
\verso{940}{}{πολύδακρυν ἰαχάν.}
\falante{ΞΕΡΞΗΣ}
\verso{941}{\notagramatical{αἰανῆ}{αἰανής/\allowbreak -ές}{}{adj., ac., sing., fem.}\crlf \notagramatical{πάνδυρτον}{πάνδυρτος/\allowbreak -ον}{muito chorado}{adj., ac., sing., fem.}}{ἵετ᾽ αἰανῆ [καὶ] πάνδυρτον}
\verso{942}{}{δύσθροον αὐδάν. δαίμων γὰρ ὅδ᾽ αὖ}
\verso{943}{\notagramatical{μετάτροπος}{μετάτροπος/\allowbreak -ον}{mudado, contrário}{adj., nom., sing., masc.}}{μετάτροπος ἐπ᾽ ἐμοί.}
\falante{ΧΟΡΟΣ}
\verso{944}{}{ἥσω τοι τὰν πάνδυρτον,}
\verso{945}{}{λαοπαθέα σέβων}
\verso{946}{\notagramatical{ἁλίτυπά}{ἁλίτυπος/\allowbreak -ον}{batido pelo mar}{adj., ac., pl., neut.}}{ἁλίτυπά τε βάρη,}
\verso{947}{}{πόλεως γέννας πενθητῆρος.}
\verso{948}{\notagramatical{κλάγξω}{κλάζω}{gritar}{fut., ind., at., 1ª s.}}{<κλάγξω> κλάγξω}
\verso{949}{\notagramatical{ἀρίδακρυν}{ἀρίδακρυς/\allowbreak -υ}{muito lacrimoso}{adj., ac., sing., masc.}}{δὲ γόον ἀρίδακρυν.}
\falante{ΞΕΡΞΗΣ}
\verso{950}{\notagramatical{ἀπηύρα}{ἀπαυράω}{privar de}{imperf., ind., at., 3ª s.}}{Ἰάνων γὰρ ἀπηύρα,}
\verso{951}{}{Ἰάνων ναύφαρκτος}
\verso{952}{\notagramatical{ἑτεραλκὴς}{ἑτεραλκής/\allowbreak -ές}{que favorece um lado}{adj., nom., sing., masc.}}{Ἄρης ἑτεραλκὴς}
\verso{953}{\notagramatical{κερσάμενος}{κείρω}{cortar, devastar}{part., aor., méd., nom., sing., masc.}}{νυχίαν πλάκα κερσάμενος}
\verso{954}{}{δυσδαίμονά τ᾽ ἀκτάν.}
\falante{ΧΟΡΟΣ}
\verso{955}{\notagramatical{βόα}{βοάω}{}{pres., imper., at., 2ª s.}\crlf \notagramatical{ἐκπεύθου}{ἐκπεύθομαι}{informar-se bem}{pres., imper., méd., 2ª s.}}{οἰοιοῖ βόα καὶ πάντ᾽ ἐκπεύθου.—}
\verso{956}{}{ποῦ δὲ φίλων ἄλλος ὄχλος,}
\verso{957}{}{ποῦ δέ σοι παραστάται,}
\verso{958}{}{οἷος ἦν Φαρανδάκης,}
\verso{959}{}{Σούσας, Πελάγων, [καὶ] Δοτάμας, ἠδ᾽ Ἀγ-}
\verso{960}{}{δαβάτας, Ψάμμις, Σουσισκάνης τ᾽}
\verso{961}{\notagramatical{λιπών}{λείπω}{}{part., aor., at., nom., sing., masc.}}{Ἀγβάτανα λιπών;}
\falante{ΞΕΡΞΗΣ}
\verso{962}{\notagramatical{ὀλοοὺς}{ὀλοός/\allowbreak -ή/\allowbreak -όν}{funesto}{adj., ac., pl., masc.}\crlf \notagramatical{ἀπέλειπον}{ἀπολείπω}{deixar para trás}{imperf., ind., at., 1ª s.}}{ὀλοοὺς ἀπέλειπον}
\verso{963}{}{Τυρίας ἐκ ναὸς}
\verso{964}{\notagramatical{ἔρροντας}{ἔρρω}{perecer}{part., pres., at., ac., pl., masc.}\crlf \notagramatical{ἀκταῖς}{ἀκτή, -ῆς, ἡ}{}{dat., pl.}}{ἔρροντας ἐπ᾽ ἀκταῖς}
\verso{965}{}{Σαλαμινιάσι στυφελοῦ}
\verso{966}{}{θείνοντας ἐπ᾽ ἀκτᾶς.}
\falante{ΧΟΡΟΣ}
\verso{967}{}{οἰοιοῖ, <βόα>: ποῦ σοι Φαρνοῦχος}
\verso{968}{}{Ἀριόμαρδός τ᾽ ἀγαθός,}
\verso{969}{}{ποῦ δὲ Σευάλκης ἄναξ,}
\verso{970}{}{ἢ Λίλαιος εὐπάτωρ,}
\verso{971}{}{Μέμφις, Θάρυβις, καὶ Μασίστρας,}
\verso{972}{}{Ἀρτεμβάρης τ᾽ ἠδ᾽ Ὑσταίχμας;}
\verso{973}{\notagramatical{ἐπανερόμαν}{ἐπανέρομαι}{perguntar de novo}{aor., ind., méd., 1ª s.}}{τάδε σ᾽ ἐπανερόμαν.}
\falante{ΞΕΡΞΗΣ}
\verso{974}{}{ἰὼ ἰώ μοι,}
\verso{975}{}{τὰς ὠγυγίους κατιδόντες}
\verso{976}{\notagramatical{πιτύλῳ}{πίτυλος, -ου, ὁ}{golpe de remos}{dat., sing.}}{στυγνὰς Ἀθάνας πάντες ἑνὶ πιτύλῳ,}
\verso{977}{\notagramatical{τλάμονες}{τλήμων/\allowbreak -ον}{}{adj., nom., pl., masc.}\crlf \notagramatical{ἀσπαίρουσι}{ἀσπαίρω}{debater-se}{pres., ind., at., 3ª p.}}{ἐὴ ἐή, τλάμονες ἀσπαίρουσι χέρσῳ.}
\falante{ΧΟΡΟΣ}
\verso{978}{}{ἦ καὶ τὸν Περσᾶν αὐτοῦ}
\verso{979}{\notagramatical{πιστὸν}{πιστός/\allowbreak -ή/\allowbreak -όν}{}{adj., ac., sing., masc.}}{τὸν σὸν πιστὸν πάντ᾽ ὀφθαλμὸν}
\verso{980}{}{μυρία μυρία πεμπαστὰν}
\verso{981}{}{Βατανώχου παῖδ᾽ Ἄλπιστον}
\verso{982}{}{[…]}
\verso{983}{}{τοῦ Σησάμα τοῦ Μεγαβάτα,}
\verso{984}{}{Πάρθον τε μέγαν τ᾽ Οἰβάρην ἔλιπες ἔλιπες;}
\verso{985}{\notagramatical{δᾴων}{δάιος/\allowbreak -α/\allowbreak -ον}{}{adj., gen., pl., masc.}}{ὢ ὢ <ὢ> δᾴων.}
\verso{986}{\notagramatical{ἀγαυοῖς}{ἀγαυός/\allowbreak -ή/\allowbreak -όν}{ilustre}{adj., dat., pl., masc.}}{Πέρσαις ἀγαυοῖς κακὰ πρόκακα λέγεις.}
\falante{ΞΕΡΞΗΣ}
\verso{987}{}{ἴυγγά μοι δῆτ᾽}
\verso{988}{}{ἀγαθῶν ἑτάρων ὑπορίνεις,}
\verso{990}{\notagramatical{ἄλαστα}{ἄλαστος/\allowbreak -ον}{inesquecível, funesto}{adj., ac., pl., neut.}}{<ἄλαστ᾽> ἄλαστα στυγνὰ πρόκακα λέγων.}
\verso{991}{\notagramatical{ἦτορ}{ἦτορ, ἤτορος, τό}{coração}{nom., sing.}}{βοᾷ βοᾷ <μοι> μελέων ἔντοσθεν ἦτορ.}
\falante{ΧΟΡΟΣ}
\verso{992}{}{καὶ μὴν ἄλλους γε ποθοῦμεν,}
\verso{993}{}{Μάρδων ἀνδρῶν μυριοταγὸν}
\verso{994}{}{Ξάνθιν ἄρειόν τ᾽ Ἀγχάρην,}
\verso{995}{}{Δίαιξίν τ᾽ ἠδ᾽ Ἀρσάκην}
\verso{996}{}{ἱππιάνακτας,}
\verso{997}{}{Κηγδαδάταν καὶ Λυθίμναν}
\verso{998}{\notagramatical{ἀκόρεστον}{ἀκόρεστος/\allowbreak -ον}{insaciável}{adj., ac., sing., masc.}}{Τόλμον τ᾽ αἰχμᾶς ἀκόρεστον.}
\verso{999}{}{ἔταφον ἔταφον,}
\verso{1000}{}{οὐκ ἀμφὶ σκηναῖς}
\verso{1001}{\notagramatical{τροχηλάτοισιν}{τροχήλατος/\allowbreak -ον}{montado sobre rodas}{adj., dat., pl., fem.}\crlf \notagramatical{ἑπομένους}{ἕπομαι}{seguir}{part., pres., méd., ac., pl., masc.}}{τροχηλάτοισιν ὄπιθεν ἑπομένους.}
\falante{ΞΕΡΞΗΣ}
\verso{1002}{\notagramatical{βεβᾶσι}{βαίνω}{}{perf., ind., at., 3ª p.}}{βεβᾶσι γὰρ τοίπερ ἀγρέται στρατοῦ.}
\falante{ΧΟΡΟΣ}
\verso{1003}{\notagramatical{βεβᾶσιν}{βαίνω}{}{perf., ind., at., 3ª p.}\crlf \notagramatical{νώνυμοι}{νώνυμος/\allowbreak -ον}{sem nome}{adj., nom., pl., masc.}}{βεβᾶσιν, οἴ, νώνυμοι.}
\falante{ΞΕΡΞΗΣ}
\verso{1004}{}{ἰὴ ἰή, ἰὼ ἰώ.}
\falante{ΧΟΡΟΣ}
\verso{1005}{}{ἰὼ ἰώ, δαίμονες,}
\verso{1006}{\notagramatical{ἄελπτον}{ἄελπτος/\allowbreak -ον}{}{adj., nom., sing., neut.}}{ἔθεσθ᾽ ἄελπτον κακὸν}
\verso{1007}{\notagramatical{δέδορκεν}{δέρκομαι}{olhar}{perf., ind., at., 3ª s.}}{διαπρέπον, οἷον δέδορκεν Ἄτα.}
\falante{ΞΕΡΞΗΣ}
\verso{1008}{\notagramatical{πεπλήγμεθ᾽}{πλήσσω}{golpear}{perf., ind., pass., 1ª p.}}{πεπλήγμεθ᾽ οἵᾳ δι᾽ αἰῶνος τύχᾳ:}
\falante{ΧΟΡΟΣ}
\verso{1009}{\notagramatical{πεπλήγμεθ᾽}{πλήσσω}{}{perf., ind., méd./pass., 1ª p.}}{πεπλήγμεθ᾽: εὔδηλα γάρ:}
\falante{ΞΕΡΞΗΣ}
\verso{1010}{\notagramatical{δύᾳ}{δύη, -ης, ἡ}{dor}{dat., sing.}}{νέᾳ νέᾳ δύᾳ δύᾳ:}
\falante{ΧΟΡΟΣ}
\verso{1011}{}{Ἰαόνων ναυβατᾶν}
\verso{1012}{\notagramatical{κύρσαντες}{κυρέω}{}{part., aor., at., nom., pl., masc.}}{κύρσαντες οὐκ εὐτυχῶς.}
\verso{1013}{\notagramatical{δυσπόλεμον}{δυσπόλεμος/\allowbreak -ον}{desafortunado na guerra}{adj., nom., sing., neut.}}{δυσπόλεμον δὴ γένος τὸ Περσᾶν.}
\falante{ΞΕΡΞΗΣ}
\verso{1014}{}{πῶς δ᾽ οὔ; στρατὸν μὲν τοσοῦ-}
\verso{1015}{\notagramatical{πέπληγμαι}{πλήσσω}{}{perf., ind., méd./pass., 1ª s.}}{τον τάλας πέπληγμαι.}
\falante{ΧΟΡΟΣ}
\verso{1016}{}{τί δ᾽ οὒκ ὄλωλεν, μέγ᾽ ἄλαστε, Περσᾶν;}
\falante{ΞΕΡΞΗΣ}
\verso{1017}{}{ὁρᾷς τὸ λοιπὸν τόδε τᾶς ἐμᾶς στολᾶς;}
\falante{ΧΟΡΟΣ}
\verso{1018}{}{ὁρῶ ὁρῶ.}
\falante{ΞΕΡΞΗΣ}
\verso{1020}{\notagramatical{ὀιστοδέγμονα}{ὀιστοδέγμων/\allowbreak -ον}{que recebe flechas}{adj., ac., sing., masc.}}{τόνδε τ᾽ ὀιστοδέγμονα—}
\falante{ΧΟΡΟΣ}
\verso{1021}{\notagramatical{σεσωσμένον}{σώζω}{}{part., perf., méd./pass., ac., sing., neut.}}{τί τόδε λέγεις σεσωσμένον;}
\falante{ΞΕΡΞΗΣ}
\verso{1022}{}{θησαυρὸν βελέεσσιν.}
\falante{ΧΟΡΟΣ}
\verso{1023}{}{βαιά γ᾽ ὡς ἀπὸ πολλῶν.}
\falante{ΞΕΡΞΗΣ}
\verso{1024}{\notagramatical{ἐσπανίσμεθ᾽}{σπανίζω}{carecer de}{mais-que-perf., ind., méd./pass., 1ª p.}}{ἐσπανίσμεθ᾽ ἀρωγῶν.}
\falante{ΧΟΡΟΣ}
\verso{1025}{}{Ἰάνων λαὸς οὐ φυγαίχμας.}
\falante{ΞΕΡΞΗΣ}
\verso{1026}{}{ἀγανόρειος: κατεῖ-}
\verso{1027}{}{δον δὲ πῆμ᾽ ἄελπτον.}
\falante{ΧΟΡΟΣ}
\verso{1028}{\notagramatical{ναύφρακτον}{ναύφρακτος/\allowbreak -ον}{provido de navios}{adj., ac., sing., masc.}}{τραπέντα ναύφρακτον ἐρεῖς ὅμιλον;}
\falante{ΞΕΡΞΗΣ}
\verso{1030}{}{πέπλον δ᾽ ἐπέρρηξ᾽ ἐπὶ συμφορᾷ κακοῦ.}
\falante{ΧΟΡΟΣ}
\verso{1031}{}{παπαῖ παπαῖ.}
\falante{ΞΕΡΞΗΣ}
\verso{1032}{}{καὶ πλέον ἢ παπαῖ μὲν οὖν.}
\falante{ΧΟΡΟΣ}
\verso{1033}{}{δίδυμα γάρ ἐστι καὶ τριπλᾶ.}
\falante{ΞΕΡΞΗΣ}
\verso{1034}{}{λυπρά, χάρματα δ᾽ ἐχθροῖς.}
\falante{ΧΟΡΟΣ}
\verso{1035}{\notagramatical{ἐκολούσθη}{κολούω}{mutilar, diminuir}{aor., ind., pass., 3ª s.}}{καὶ σθένος γ᾽ ἐκολούσθη—}
\falante{ΞΕΡΞΗΣ}
\verso{1036}{}{γυμνός εἰμι προπομπῶν.}
\falante{ΧΟΡΟΣ}
\verso{1037}{}{φίλων ἄταισι ποντίαισιν.}
\falante{ΞΕΡΞΗΣ}
\verso{1038}{}{δίαινε δίαινε πῆμα: πρὸς δόμους δ᾽ ἴθι.}
\falante{ΧΟΡΟΣ}
\verso{1039}{\notagramatical{διαίνομαι}{διαίνω}{molhar de lágrimas}{pres., ind., méd./pass., 1ª s.}\crlf \notagramatical{γοεδνὸς}{γοεδνός/\allowbreak -ή/\allowbreak -όν}{lamentoso}{adj., nom., sing., masc.}\crlf \notagramatical{ὤν}{εἰμί}{}{part., pres., at., nom., sing., masc.}}{διαίνομαι γοεδνὸς ὤν.}
\falante{ΞΕΡΞΗΣ}
\verso{1040}{}{βόα νυν ἀντίδουπά μοι.}
\falante{ΧΟΡΟΣ}
\verso{1041}{}{δόσιν κακὰν κακῶν κακοῖς.}
\falante{ΞΕΡΞΗΣ}
\verso{1042}{\notagramatical{ἴυζε}{ἰύζω}{}{pres., imper., at., 2ª s.}\crlf \notagramatical{τιθείς}{τίθημι}{pôr, estabelecer}{part., pres., at., nom., sing., masc.}}{ἴυζε μέλος ὁμοῦ τιθείς.}
\falante{ΧΟΡΟΣ}
\verso{1043}{}{ὀτοτοτοτοῖ.}
\verso{1044}{}{βαρεῖά γ᾽ ἅδε συμφορά.}
\verso{1045}{}{οἲ μάλα καὶ τόδ᾽ ἀλγῶ.}
\falante{ΞΕΡΞΗΣ}
\verso{1046}{\notagramatical{ἔρεσσ᾽}{ἐρέσσω}{}{pres., imper., at., 2ª s.}}{ἔρεσσ᾽ ἔρεσσε καὶ στέναζ᾽ ἐμὴν χάριν.}
\falante{ΧΟΡΟΣ}
\verso{1047}{}{αἰαῖ αἰαῖ, δύα δύα.}
\falante{ΞΕΡΞΗΣ}
\verso{1048}{\notagramatical{ἀντίδουπά}{ἀντίδουπος/\allowbreak -ον}{que responde em eco}{adj., nom., pl., neut.}}{βόα νυν ἀντίδουπά μοι.}
\falante{ΧΟΡΟΣ}
\verso{1049}{\notagramatical{πάρεστι}{πάρειμι}{}{pres., ind., at., 3ª s.}}{μέλειν πάρεστι, δέσποτα.}
\falante{ΞΕΡΞΗΣ}
\verso{1050}{\notagramatical{ἐπορθίαζέ}{ἐπορθιάζω}{erguer a voz}{pres., imper., at., 2ª s.}}{ἐπορθίαζέ νυν γόοις.}
\falante{ΧΟΡΟΣ}
\verso{1051}{}{ὀτοτοτοτοῖ.}
\verso{1052}{}{μέλαινα δ᾽ ἀμμεμίξεται,}
\verso{1053}{\notagramatical{στονόεσσα}{στονόεις/\allowbreak -εσσα/\allowbreak -εν}{lamentoso}{adj., nom., sing., fem.}}{οἴ, στονόεσσα πλαγά.}
\falante{ΞΕΡΞΗΣ}
\verso{1054}{\notagramatical{ἄρασσε}{ἀράσσω}{}{pres., imper., at., 2ª s.}}{καὶ στέρν᾽ ἄρασσε καὶ βόα τὸ Μύσιον.}
\falante{ΧΟΡΟΣ}
\verso{1055}{}{ἄνι᾽ ἄνια.}
\falante{ΞΕΡΞΗΣ}
\verso{1056}{\notagramatical{πέρθε}{πέρθω}{}{pres., imper., at., 2ª s.}\crlf \notagramatical{λευκήρη}{λευκήρης/\allowbreak -ες}{branco}{adj., ac., sing., fem.}}{καί μοι γενείου πέρθε λευκήρη τρίχα.}
\falante{ΧΟΡΟΣ}
\verso{1057}{}{ἄπριγδ᾽ ἄπριγδα μάλα γοεδνά.}
\falante{ΞΕΡΞΗΣ}
\verso{1058}{\notagramatical{ἀύτει}{ἀυτέω}{gritar}{pres., imper., at., 2ª s.}}{ἀύτει δ᾽ ὀξύ.}
\falante{ΧΟΡΟΣ}
\verso{1059}{\notagramatical{ἔρξω}{ἔρδω}{}{fut., ind., at., 1ª s.}}{καὶ τάδ᾽ ἔρξω.}
\falante{ΞΕΡΞΗΣ}
\verso{1060}{\notagramatical{ἔρεικε}{ἐρείκω}{rasgar}{pres., imper., at., 2ª s.}}{πέπλον δ᾽ ἔρεικε κολπίαν ἀκμῇ χερῶν.}
\falante{ΧΟΡΟΣ}
\verso{1061}{}{ἄνι᾽ ἄνια.}
\falante{ΞΕΡΞΗΣ}
\verso{1062}{\notagramatical{ψάλλ᾽}{ψάλλω}{arrancar, puxar}{pres., imper., at., 2ª s.}\crlf \notagramatical{ἔθειραν}{ἔθειρα, -ας, ἡ}{cabeleira}{ac., sing.}\crlf \notagramatical{κατοίκτισαι}{κατοικτίζω}{lamentar}{aor., imper., méd., 2ª s.}}{καὶ ψάλλ᾽ ἔθειραν καὶ κατοίκτισαι στρατόν.}
\falante{ΧΟΡΟΣ}
\verso{1063}{}{ἄπριγδ᾽ ἄπριγδα μάλα γοεδνά.}
\falante{ΞΕΡΞΗΣ}
\verso{1064}{}{διαίνου δ᾽ ὄσσε.}
\falante{ΧΟΡΟΣ}
\verso{1065}{\notagramatical{τέγγομαί}{τέγγω}{}{pres., ind., méd., 1ª s.}}{τέγγομαί τοι.}
\falante{ΞΕΡΞΗΣ}
\verso{1066}{}{βόα νυν ἀντίδουπά μοι.}
\falante{ΧΟΡΟΣ}
\verso{1067}{}{οἰοῖ οἰοῖ.}
\falante{ΞΕΡΞΗΣ}
\verso{1068}{\notagramatical{κίε}{κίω}{ir}{pres., imper., at., 2ª s.}}{αἰακτὸς ἐς δόμους κίε.}
\falante{ΧΟΡΟΣ}
\verso{1069}{\notagramatical{δύσβατος}{δύσβατος/\allowbreak -ον}{difícil de atravessar}{adj., nom., sing., fem.}}{ἰὼ ἰώ, Περσὶς αἶα δύσβατος.}
\falante{ΞΕΡΞΗΣ}
\verso{1070}{}{ἰωὰ δὴ κατ᾽ ἄστυ.}
\falante{ΧΟΡΟΣ}
\verso{1071}{}{ἰωὰ δῆτα, ναὶ ναί.}
\falante{ΞΕΡΞΗΣ}
\verso{1072}{}{γοᾶσθ᾽ ἁβροβάται.}
\falante{ΧΟΡΟΣ}
\verso{1073}{}{ἰὼ ἰώ, Περσὶς αἶα δύσβατος.}
\falante{ΞΕΡΞΗΣ}
\verso{1074}{\notagramatical{τρισκάλμοισιν}{τρίσκαλμος/\allowbreak -ον}{}{adj., dat., pl., fem.}}{ἰὴ ἰὴ τρισκάλμοισιν,}
\verso{1075}{\notagramatical{βάρισιν}{βᾶρις, βάριδος, ἡ}{}{dat., pl.}}{ἰὴ ἰή, βάρισιν ὀλόμενοι.}
\falante{ΧΟΡΟΣ}
\verso{1076}{\notagramatical{δυσθρόοις}{δύσθροος/\allowbreak -ον}{de som lamentoso}{adj., dat., pl., masc.}}{πέμψω τοί σε δυσθρόοις γόοις.}
